% BEGIN INCLUDED SGA4_English_current_buildable_source_20260906/source/04/drafts/SGA4_IV_LINES_31_187_English.texfrag
\setcounter{pageoriginal}{298}
\setcounter{section}{-1}
\section{引言}\etiquette[0]{IV.sec0}\pageoriginale

\subsection{}\etiquette[0.1]{IV.subsec0.1}
我们在\hyperref[sga4.exp.II]{第 II 篇}中看到了形如 $\widetilde{𝒞}=$ $𝒞$ 上的集合层范畴的各种正合性质, 其中 $𝒞$ 为小景. 可以说, 这些性质表明此类范畴 (我们将称之为\emph{拓扑斯}) 在许多方面继承了 (小) 集合范畴 $(\Ens)$ 的熟知性质. 另一方面, 经验表明, 数学中的各种情形应当\emph{首先看作构造相应的 (集合) 层范畴的技术手段}, 即\emph{构造相应的「拓扑斯」的技术手段}. 看来, 与景相联系的一切真正重要的概念 (例如在\hyperref[sga4.exp.V]{第 V 篇}中研究的上同调不变量; 其他各种「拓扑」不变量, 如 M. \textsc{Artin} 与 B. \textsc{Mazur} 最近研究的同伦不变量 [1]; 以及 J. \textsc{Giraud} 的非交换上同调著作中研究的概念), 其实都直接以相应的拓扑斯表述. 从这个观点看, 当两个景相应的拓扑斯为等价范畴时, 就应把这两个景看作本质等价的. 给定一个景 (至少在实践中特别重要的情形, 即其拓扑比典范拓扑粗的情形), 应当看作相当于给定一个拓扑斯 $ℰ$ (即由该景上的集合层组成的相应拓扑斯), 再加上 $ℰ$ 的一个对象生成族 (参见\pageoriginale II \hyperref[II.subsec4.9]{4.9} 及下面的 \hyperref[IV.cor1.2.1]{1.2.1}). 这种观点类似于: 从一组生成元以及这些生成元之间的一组关系联系出一个群, 并把兴趣放在这个群的结构上, 而不是用于生成它的生成元与关系系上 (后者看作该情形中的辅助资料). 此外, 「比较引理」III \hyperref[III.theo4.1]{4.1} 提供了许多景对 $𝒞$, $𝒞'$ 的例子, 它们不同构, 甚至作为范畴也不等价, 却给出等价的拓扑斯; 因而应当把 $𝒞$ 与 $𝒞'$ 看作本质等价的.

\subsection{}\etiquette[0.2]{IV.subsec0.2}
本篇用简单的正合性质给出拓扑斯的刻画 (归于 J. \textsc{Giraud}); 受一个拓扑空间到另一个拓扑空间的连续映射概念的启发, 研究拓扑斯态射的自然概念; 并在拓扑斯的框架内发展通常层论中熟悉的某些构造 ($\Hhom$ 层、张量积层、支集). 最后, 按照 M. \textsc{Artin} 的做法, 说明如何从一个「开」部分、其互补「闭」部分, 以及联系两者的某个左正合函子重构拓扑斯, 而该函子实际上几乎可以任意选取.

\subsection{}\etiquette[0.3]{IV.subsec0.3}
这样就得到一个粘合拓扑斯的程序; 把它应用于来自通常拓扑空间的拓扑斯, 一般会得到一个不再属于同一类型的拓扑斯. 这首次表明, 拓扑斯概念在各种自然构造下具有显著的稳定性, 而拓扑空间概念 (拓扑斯概念正由此\pageoriginale 得到启发) 却缺乏这种稳定性. 作为另一个显著的例子, 还要提到拓扑斯中一个群的相对分类拓扑斯 (参见 [1] 或下面的\hyperref[IV.exer5.9]{5.9}), 它受拓扑群的分类空间这一经典概念的启发; 以及与代数几何或解析几何中各种「模问题」相联系的「模拓扑斯」概念 [10] [13].

其他拓扑斯, 如概形的\emph{平展拓扑斯} (本讨论班自第 VII 篇起系统研究它) 或相对概形的\emph{晶体拓扑斯} [6], 则在希望为抽象代数簇 (更一般地, 为概形) 发展可用的上同调理论, 以替代复数域上代数簇的经典 Betti 上同调时自然出现.

\subsection{}\etiquette[0.4]{IV.subsec0.4}
因此可以说, 拓扑斯概念作为拓扑学中\emph{层论观点}的自然产物, 反过来构成了拓扑空间概念的大幅扩充\footnote[1]{关于拓扑斯概念与拓扑空间概念之间的确切关系, 参见 [9] 或下面的 \hyperref[IV.subsec4.1]{4.1} 及 \hyperref[IV.subsec4.2]{4.2}.}, 它涵盖了许多从前不被认为属于拓扑直观范畴的情形. 此类情形的特征在于存在一个「局部化」概念, 而景的概念, 最终则是拓扑斯的概念 (通过景所对应的拓扑斯), 正是对它的形式化. 正如「topos」一词本来就意在暗示的那样, 在本讨论班的作者看来, 将拓扑学视为研究\emph{拓扑斯} (而不仅仅是拓扑空间) 的学科, 是合理而正当的.

\subsection{}\etiquette[0.5]{IV.subsec0.5}
我们认为, 在本篇关于拓扑斯的一般论述中收入相当多的例子是有用的, 尽管其中许多与本讨论班最初的目标 (即研究平展上同调) 只有间接\pageoriginale 联系. 只对平展上同调感兴趣而又急于往下读的读者, 当然可以毫无妨碍地略去这些例子, 以及本篇的大部分内容, 只在有需要时查阅.

\section{拓扑斯的定义与刻画}\etiquette[1]{IV.sec1}

\begin{enonce*}{定义 1.1}\etiquette[1.1]{IV.def1.1}
范畴 $E$ 称为 $𝒰$-拓扑斯, 或在不会引起混淆时简称拓扑斯, 如果存在景 $C\in 𝒰$, 使得 $E$ 等价于 $C$ 上的 $𝒰$-集合层范畴 $C^{\widetilde{\phantom{,}}}$.
\end{enonce*}

\setcounter{subsection}{1}
\subsubsection{}\etiquette[1.1.1]{IV.subsubsec1.1.1}
设 $E$ 为 $𝒰$-拓扑斯. 总把 $E$ 看作赋有典范拓扑 (\hyperrefext[II][II.def2.5]{2.5}); 这使它成为\emph{景}, 并且由下面的 \hyperref[IV.subsubsec1.1.2]{1.1.2} d), 甚至成为 $𝒰$-景 (\hyperrefext[II][II.def3.0.2]{3.0.2}). 除非另有明说, 我们不考虑 $E$ 上上述拓扑以外的其他拓扑.

\subsubsection{}\etiquette[1.1.2]{IV.subsubsec1.1.2}
我们在 \hyperrefext[II][II.prop4.8]{4.8}, \hyperrefext[II][II.cor4.11]{4.11} 中看到, $𝒰$-拓扑斯 $E$ 是满足下述条件的 $𝒰$-范畴 (\hyperrefext[I][I.def1.1]{1.1}):
\begin{enumerate}
\renewcommand{\theenumi}{\alph{enumi}}
\renewcommand{\labelenumi}{\theenumi)}
\item 有限投射极限在 $E$ 中可表.

\item 以 $𝒰$ 的元素为指标的直和在 $E$ 中可表. 它们是不交且泛的 (\hyperrefext[II][II.def4.5]{4.5}).

\item $E$ 中的等价关系都是泛有效的 \sisi{(\hyperrefext[I][I.sec10.9]{10.9}).}{(\hyperrefext[I][I.sec10.10]{10.10}).}

\item $E$ 有一个以 $𝒰$ 的元素为指标的生成族 (\hyperrefext[II][II.subsec4.9]{4.9}).
\end{enumerate}

事实上\pageoriginale, 我们将看到这些内在性质\emph{刻画}了 $𝒰$-拓扑斯:

% END INCLUDED SGA4_English_current_buildable_source_20260906/source/04/drafts/SGA4_IV_LINES_31_187_English.texfrag


% BEGIN INCLUDED SGA4_English_current_buildable_source_20260906/source/04/drafts/SGA4_IV_LINES_188_382_English.texfrag
\begin{enonce*}{定理 1.2}\etiquette[1.2]{IV.theo1.2}
{\rm (J. Giraud).} 设 $E$ 为 $𝒰$-范畴. 以下性质等价:
\begin{itemize}
\item[i)] $E$ 为 $𝒰$-拓扑斯 $(\hyperref[IV.def1.1]{1.1})$.

\item[ii)] $E$ 满足 {\rm \hyperref[IV.subsubsec1.1.2]{1.1.2} 的条件 a), b), c) 和 d).}

\item[iii)] $E$ 上关于典范拓扑的 $𝒰$-层都可表, 且 $E$ 有一个小生成族 (条件 {\rm \hyperref[IV.subsubsec1.1.2]{1.1.2} d)}).

\item[iv)] 存在范畴 $C\in 𝒰$ 以及全忠实函子 $i:E\rightarrow \widehat{C}$ (其中 $\widehat{C}$ 表示 $C$ 上的 $𝒰$-预层范畴), 它有一个左正合的左伴随 a.

\item[${\rm i')}$] 存在景 $C\in 𝒰$, 使得投射极限在 $C$ 中可表, $C$ 的拓扑比其典范拓扑粗 {\rm (II \hyperref[II.def2.5]{2.5}),} %% label!
且 $E$ 等价于 $C$ 上的 $𝒰$-集合层范畴 $C^{\widetilde{\phantom{,}}}$.
\end{itemize}
\end{enonce*}

此外:

\begin{enonce*}{推论 1.2.1}\etiquette[1.2.1]{IV.cor1.2.1}
设 $E$ 为 $𝒰$-拓扑斯, $C$ 为 $E$ 的满子范畴, 赋予 $E$ 的拓扑 (\hyperref[IV.subsubsec1.1.1]{1.1.1}) 诱导的拓扑 {\rm (\hyperrefext[III][III.subsec3.1]{3.1})}. 考虑函子
$$
E\longrightarrow C^{\widetilde{\phantom{,}}}
$$
它将每个 $X\in \ob E$ 映为由 $X$ 所表的函子在 $C$ 上的限制. 此函子为范畴等价, 当且仅当 $\ob C$ 为 $E$ 的生成族.
\end{enonce*}

等价性\pageoriginale i) $\Leftrightarrow$ iv) 直接由 \hyperrefext[II][II.theo5.5]{5.5} 得出, 并且上面已经回顾了 i) $\Rightarrow$ ii). 因为 ${\rm i'}$) $\Rightarrow$ i) 显然, 只须证明 ii) $\Rightarrow$ iii) 和 iii) $\Rightarrow$ ${\rm i'}$), 这将在下面的 \hyperref[IV.subsubsec1.2.4]{1.2.4} 与 \hyperref[IV.subsubsec1.2.3]{1.2.3} 中完成. 然后注意所讨论的函子分解为 $E\rightarrow E^{\widetilde{\phantom{,}}}\rightarrow C^{\widetilde{\phantom{,}}}$, 即得推论; 因为由 (iii), $E\rightarrow E^{\widetilde{\phantom{,}}}$ 为等价, 问题化为确定何时 $E^{\widetilde{\phantom{,}}}\rightarrow C^{\widetilde{\phantom{,}}}$ 为等价. 于是结论由「比较引理」\hyperrefext[III][III.theo4.1]{4.1} 得出; 参见下面 \hyperref[IV.subsubsec1.2.3]{1.2.3} 中的证明.

\setcounter{subsection}{2} \setcounter{subsubsection}{2}
\subsubsection{}\etiquette[1.2.3]{IV.subsubsec1.2.3}
\emph{iii) $\Rightarrow$ ${\rm i'}$) 的证明}. 设 $\mathfrak{X}=(X_i)_{i\in I}$, $I\in 𝒰$, 为 $E$ 的小生成族. 因为 $E$ 为 $𝒰$-范畴, $E$ 中以各 $X_i$ 为对象的有限图表的同构类集是 $𝒰$-小的. 因此, 包含 $\mathfrak{X}$ 的对象的有限投射极限的最小的 $E$ 的对象集 $\mathfrak{X}'$ 是小集的可数并, 因而为小集~\footnote[1]{这里在同构的意义下, 用 $E$ 的对象的类进行论证.}. 置 $\mathfrak{X}^{(n+1)}=\sisi{\mathfrak{X}^{(n)(1)}}{(\mathfrak{X}^{(n)})'}$, 及 $\mathfrak{X}=\bigcup_n \mathfrak{X}^{(n)}$, 可见在扩大生成族之后, 可假设 $\mathfrak{X}$ 对有限投射极限稳定. 设 $𝒱$ 为包含 $𝒰$ 的宇宙, 使得 $E$ 为 $𝒱$-小的. 对 $E$ 的每个对象 $H$, 用 $I(H)$ 表示集合 $\amalg_{i\in I}\Hom(X_i, H)$. 设 $H$ 为 $E$ 的对象, 并设 $H'\hookrightarrow H$ 为 $H$ 关于典范拓扑的子层, 它是各态射 $u:X_i\rightarrow H$, $(u,i)\in I(H)$ 的像的「并」(\hyperrefext[II][II.theo4.1]{4.1}). 因为 $H'$ 为 $𝒰$-层的子层, $H'$ 为 $𝒰$-层, 因而可表. 因为族 $\mathfrak{X}$ 是生成族, 且对每个 $i\in I$, 映射 $\Hom(X_i, H')\rightarrow \Hom(X_i,H)$ 为双射, 态射 $H'\hookrightarrow H$ 为同构 (\hyperrefext[II][II.subsec4.9]{4.9}). 因此 (\sisi{\hyperrefext[II][II.def6.1]{6.1}}{\hyperrefext[II][II.def5.2]{5.2}}), 族 $(u:X_i\rightarrow H)$, $(u,i)\in I(H)$ 对 $E$ 的典范拓扑为覆盖族. 这样\pageoriginale, $E$ 的每个对象都可由 $\mathfrak{X}$ 的对象按 $E$ 的典范拓扑覆盖. 设 $C$ 为由 $\mathfrak{X}$ 的对象定义的 $E$ 的满子范畴, $u:C\hookrightarrow E$ 为包含函子. $E$ 的典范拓扑在 $C$ 上诱导的拓扑 $\Cc$ 比 $C$ 的典范拓扑粗, 且当 $𝒰$ 含有无限势的元素时, 有限投射极限在 $C$ 中可表. 由 \hyperrefext[III][III.theo4.1]{4.1} 可知, 函子 $F\mapsto F\circ u$ 是从 $E$ 到 $C$ 上关于拓扑 $\Cc$ 的 $𝒰$-层范畴的等价, \cqfd

\subsubsection{}\etiquette[1.2.4]{IV.subsubsec1.2.4}
\emph{ii) $\Rightarrow$ iii) 的证明}. 证明分为四步.

设 $𝒱$ 为包含 $𝒰$ 的宇宙, 使得 $E$ 为 $𝒱$ 的元素. 设 $\widetilde{E}$ 为 $E$ 上关于典范拓扑的 $𝒱$-层范畴, $J_E:E\rightarrow \widetilde{E}$ 为典范函子.

\paragraph{}\etiquette[1.2.4.1]{IV.par1.2.4.1}
设 $(g_i:G_i\rightarrow H)$, $i\in I\in 𝒰$, 为 $\widetilde{E}$ 中的满态射族. 若各 $G_i$ 及纤维积 $\fprod{G_i}{G_j}{H}$ 可表, 则层 $H$ 可表.

事实上, 直和 $\coprod_{i\in I}G_i$ 在 $\widetilde{E}$ 中 (\hyperrefext[II][II.theo4.1]{4.1}) 由一个可表层 $G$ 所表 (性质 b) 及 \hyperrefext[II][II.prop4.6]{4.6}). 对直和 $K=\coprod_{(i,j)\in I\times I}\fprod{G_i}{G_j}{H}$ 也同样成立. 此外, 图表 $K\rightrightarrows G\rightarrow H$ 正合, 且 $K$ 为 $G$ 在 $H$ 上的纤维平方. (后一性质在预层范畴中成立, 从而在层范畴中也成立 (\hyperrefext[II][II.theo4.1]{4.1}).) 由 c) 及 \hyperrefext[II][II.prop4.7]{4.7} 得层 $H$ 可表.

\paragraph{}\etiquette[1.2.4.2]{IV.par1.2.4.2}
设 $X_{\alpha}$, $\alpha\in A\in 𝒰$, 为 $E$ 的生成族. 对每个层 $H$, 用 $I(H)$ 表示集合 $\coprod_{\alpha \in A}\Hom_E (X_{\alpha}, H)$. 族 $(u:X_{\alpha}\rightarrow H, (u,\alpha)\in I(H))$\pageoriginale 在 $\widetilde{E}$ 中为满态射族. 当 $H$ 为 $𝒰$-层时, $I(H)$ 为 $𝒰$ 的元素.

事实上, 每个层 $H$ 都是一个以可表层为源的态射满态射族的靶 (\hyperrefext[I][I.prop3.4]{3.4} 及 \hyperrefext[II][II.theo4.1]{4.1}), 因此只须在层 $H$ 可表时证明第一个断言. 于是设 $G$ 为族 $(u:X_{\alpha}\rightarrow H, (u,\alpha)\in I(H))$ 的像. 态射 $G\rightarrow H$ 为单态射. 因而处于第一步的情形, 因为 $\fprod{X_{\alpha}}{X_{\beta}}{G}$ 同构于可表的 $\fprod{X_{\alpha}}{X_{\beta}}{H}$, 且 $I(H)$ 为 $𝒰$ 的元素. 从而 $G$ 可表. 但因为 $X_{\alpha}$ 是生成族, $G\rightarrow H$ 为同构. 最后一个断言直接由 $𝒰$-层的定义得出.

\paragraph{}\etiquette[1.2.4.3]{IV.par1.2.4.3}
可表层的每个子层都可表.

事实上, 设 $G\rightarrow H$ 为可表层的子层. 则 $G$ 为 $𝒰$-层. 因此族 $(u:X_{\alpha}\rightarrow G, (u,\alpha)\in I(G))$ 在 $\widetilde{E}$ 中为满态射族, 且以 $𝒰$ 的元素为指标. 此外, 纤维积 $\fprod{X_{\alpha}}{X_{\beta}}{G}$ 同构于可表的纤维积 $\fprod{X_{\alpha}}{X_{\beta}}{H}$. 因而处于第一步的情形, $G$ 可表.

\paragraph{}\etiquette[1.2.4.4]{IV.par1.2.4.4}
每个 $𝒰$-层都可表.

事实上, 由第一步与第二步, 只须说明纤维积 $\fprod{X_{\alpha}}{X_{\beta}}{H}$ 可表. 但这些纤维积是积 $\fprod{X_{\alpha}}{X_{\beta}}{}$ 的子对象. 因此由第三步即得结论. 这就完成了定理 \hyperref[IV.theo1.2]{1.2} 的证明.

\begin{enonce*}[remark]{注 1.3}\etiquette[1.3]{IV.rem1.3}
当然\pageoriginale, 对给定的 $𝒰$-拓扑斯 $E$, 一般没有特别优先的方式把它在等价的意义下表为 $C^{\widetilde{\phantom{,}}}$, 其中 $C$ 为\emph{小}景; 或者, 当只考虑拓扑比典范拓扑粗的 $C$ 时, 由 \hyperref[IV.cor1.2.1]{1.2.1}, 实质等价地说, $E$ 中没有特别优先的\emph{小}生成族. 另一方面, 当不再要求 $C$ 为小的时, (由 \hyperref[IV.theo1.2]{1.2} iii)) 就可以完全典范地选取一个 $𝒰$-景 $C$, 使得 $E$ 等价于 $C^{\widetilde{\phantom{,}}}$, 即取 $E$ 自身! 这就是实践中只使用小景会不方便的技术原因之一: 事实上, 最重要的景, 即 $𝒰$-拓扑斯, 都不是小的! 此外, 在代数几何的许多问题 (甚至在拓扑学中, 参见 \hyperref[IV.subsec2.5]{2.5}) 出现的生成拓扑斯的景, 也不是小的; 例如概形的平展景 \sisi{(VIII.1).}{(VII.1).}
\end{enonce*}

% END INCLUDED SGA4_English_current_buildable_source_20260906/source/04/drafts/SGA4_IV_LINES_188_382_English.texfrag


% BEGIN INCLUDED SGA4_English_current_buildable_source_20260906/source/04/drafts/SGA4_IV_LINES_383_536_English.texfrag
\begin{enonce*}{命题 1.4}\etiquette[1.4]{IV.prop1.4}
设 $E$ 为 $𝒰$-拓扑斯, $E'$ 为范畴 (不一定为 $𝒰$-范畴), $F$ 为 $E$ 上取值于 $E'$ 的预层. $F$ 为取值于 $E'$ 的层 {\rm(\hyperrefext[II][II.def6.1]{6.1})} 的充分必要条件是, $F$ 将 $E$ 中的 $𝒰$-归纳极限变成 $E'$ 中的投射极限.
\end{enonce*}

设 $𝒱$ 为包含 $𝒰$ 的宇宙, 使得 $E'$ 为 $𝒱$-范畴. 将 $F$ 与 $E'$ 的对象 $X'$ 所定义的函子 $\Hom(X',-):E'\rightarrow 𝒱\text{-}\Ens$ 复合, 就化归到 $E'=𝒱\text{-}\Ens$ 的情形, 其中 $𝒱$ 为宇宙. 假设 $F$ 把 $𝒰$-归纳极限变成投射极限. 则由定义立即可知 $F$ 为层, 因为按 $ℰ$ 的典范拓扑的定义, $ℰ$ 中的覆盖族\pageoriginale $X_i\rightarrow X$ 使得可以把 $X$ 看作图表 $ \left(
\fprod{X_i}{X_j}{X}\raise6pt\hbox{$\mathop{\nearrow}^{\mbox{$X_i$}}_{\displaystyle{\kern-10pt\raise-8pt\hbox{$\searrow_{\mbox{$X_j$}}$}}}$}
\right) $ 的归纳极限. 假设 $F$ 为层, 来证明它将 $𝒰$-归纳极限变成投射极限. 若 $𝒱\subset 𝒰$, 可假设 $𝒱=𝒰$, 应用判别准则 \hyperref[IV.theo1.2]{1.2} iii) 即可. 要处理一般情形, 还须稍费些功夫. 设 $C$ 为 $E$ 中由一个以 $𝒰$ 的元素为指标的生成族所生成的满子范畴, $u:C\rightarrow E$ 为包含函子. 给 $C$ 赋予 $E$ 的典范拓扑所诱导的拓扑 \sisi{(III 4.5).}{(III \hyperref[III.prop3.5]{3.5}).}
用 $C^{\widetilde{\phantom{,}}}_{𝒰}$ 与 $C^{\widetilde{\phantom{,}}}_{𝒱}$ 表示 $C$ 上的层范畴, 用 $E^{\widetilde{\phantom{,}}}_{𝒱}$ 表示 $E$ 上关于典范拓扑的 $V$-层范畴, 用 $J_E:E\rightarrow E_{𝒱}$ 表示典范函子, 并用 $i_{𝒰,𝒱}:C^{\widetilde{\phantom{,}}}_{𝒰}\rightarrow C^{\widetilde{\phantom{,}}}_{𝒱}$ 表示包含函子. 图表
\[\vcenter{
\xymatrix@C=3.25cm@R=1.25cm{
C^{\widetilde{\phantom{a}}}_{\underline{U}} \ar[r]^{i_{\underline{U},\underline{V}}} & C^{\widetilde{\phantom{a}}}_{\underline{V}}\\
E  \ar[r]^{J_E} \ar[u]&
E^{\widetilde{\phantom{a}}}_{\underline{V}}\ar[u]\quad,
}}\leqno{(*)}
\]
交换, 其中竖直箭头由函子 $F\mapsto F\circ u$ 诱导. 由相伴层函子的显式构造 (\hyperrefext[II][II.sec3]{3}) 以及 \hyperrefext[II][II.theo4.1]{4.1} 可知, 函子 $i_{𝒰,𝒱}$ 与 $𝒰$-归纳极限交换. 此外, 由 \hyperrefext[III][III.subsec5.1]{5.1} 及 \hyperref[IV.cor1.5]{1.5} 可知, $(\ast)$ 中的竖直箭头均为范畴等价. 因此, $J_E:E\rightarrow E^{\widetilde{\phantom{,}}}_{𝒱}$ 与 $𝒰$-归纳极限交换. 对 $E$ 的每个对象 $X$, 有
$$
F(X)\simeq
\Hom_{E^{\widetilde{\phantom{,}}}_{𝒱}}(J_E(X),F){\quoi}.
$$
所以\pageoriginale $F$ 将 $E$ 中的 $𝒰$-归纳极限变成投射极限.

\begin{enonce*}{推论 1.5}\etiquette[1.5]{IV.cor1.5}
设 $E$ 为 $𝒰$-拓扑斯, $E'$ 为 $𝒰$-范畴, $f:E\rightarrow E'$ 为函子. $f$ 有右伴随的充分必要条件是, $f$ 与 $𝒰$-归纳极限交换.
\end{enonce*}

这是 \hyperref[IV.prop1.4]{1.4} 在 $E'=𝒰\text{-}\Ens$ 的情形下的直接推论.

\begin{enonce*}{推论 1.6}\etiquette[1.6]{IV.cor1.6}
设 $E$ 与 $E'$ 为两个 $𝒰$-拓扑斯, $f:E\rightarrow E'$ 为函子. 以下条件等价:
\begin{enumerate}
\renewcommand{\theenumi}{\roman{enumi}}
\renewcommand{\labelenumi}{\theenumi)}
\item $f$ 与 $𝒰$-归纳极限交换.

\item $f$ 有右伴随.

\item $f$ 连续 {\rm(\hyperrefext[III][III.def1.1]{1.1}).}
\end{enumerate}
\end{enonce*}

等价性 i) $\Leftrightarrow$ ii) 已在 \hyperref[IV.cor1.5]{1.5} 中证明. 为证明 i) $\Leftrightarrow$ iii), 应用连续函子的定义, 选取宇宙 $𝒱$ 使得 $𝒰\in 𝒱$ (因此 $E$ 与 $E'$ 为 $𝒱$-小景). 须表明, 对 $E'$ 上的每个 $𝒱$-层 $F$, 复合 $F\circ f$ 为 $E$ 上的层, 即 (\hyperref[IV.prop1.4]{1.4}) 它将 $𝒰$-归纳极限变成投射极限. 为此, i) 成立是充分的, 因为由 \hyperref[IV.prop1.4]{1.4}, $F$ 自身把 $𝒰$-归纳极限变成投射极限; 取 $F$ 为可表函子, 也可见 i) 是必要的.

\begin{enonce*}{推论 1.7}\etiquette[1.7]{IV.cor1.7}
用 \hyperref[IV.cor1.6]{1.6} 的记号, $f$ 是拓扑斯态射 $(\hyperref[IV.def3.1]{3.1})$ $u:E'\rightarrow E$ 的逆像函子 $u^{\ast}$ 的充分必要条件是, $f$ 左正合且将满态射族映为满态射族.
\end{enonce*}

必要性\pageoriginale 由定义显然 (注意每个右正合函子都把满态射映为满态射). 充分性由 \sisi{\hyperrefext[III][III.prop2.6]{2.6}}{\hyperrefext[III][III.prop1.6]{1.6}} 得出, 因为它蕴涵在所述条件下 $f$ 连续, 从而由 \hyperref[IV.cor1.6]{1.6} 可知它与 $𝒰$-归纳极限交换.

\begin{enonce*}[remark]{注 1.8}\etiquette[1.8]{IV.rem1.8}
用 \hyperref[IV.cor1.5]{1.5} 的记号, 可证明 $f$ 有左伴随当且仅当它与 $𝒰$-投射极限交换. 换言之, \emph{协变}函子 $E\rightarrow (𝒰\text{-}\Ens)$ 可表当且仅当它与 $𝒰$-投射极限交换\footnote[1]{更一般的陈述见 \hyperrefext[I][I.cor8.12.8]{8.12.8}, \hyperrefext[I][I.exer8.12.9]{8.12.9}; 以下 a) 至 c) 的证明提纲对应于该处给出的证明.}. 这里只指出证明的原则, 分为三步:
\begin{enumerate}
\renewcommand{\theenumi}{\alph{enumi}}
\renewcommand{\labelenumi}{\theenumi)}
\item 标准论证 [5, ${\rm n^{\circ}}$ 195, \S 3] 表明, 若 $F$ 与 $\varprojlim$ 交换, 则它由一个严格投射系 $(T_i)_{i\in I}$ pro-表示, 其中 $I$ 为滤过有序集, \emph{不一定是小的}. 可假定若 $i>j$, 则 $T_i\rightarrow T_j$ 不是同构, 且在此假设下, $F$ 可表当且仅当 $I$ 为小的 (这事实上蕴涵该投射系本质常值).

\item 已知对 $E$ 的每个对象 $X$, 集合 $F(X)=\varinjlim_i \Hom(T_i,X)$ 为小的, 为证明 $I$ 为小的, 只须有一个\emph{小余生成族} $(X_j)_{j\in J}$ (即它是反范畴 $E^{\circ}$ 的生成族). 而可以证明每个 $𝒰$-拓扑斯 $E$ 都有小余生成族.

\item 为证明最后一点, 用标准论证 [Toh] 注意到 $E$ 的每个对象 $X$ 都有一个到「内射」对象的单态射; 再注意到对 $E$ 的每个生成族 $(L_{\alpha})$, 若\pageoriginale 每个 $L_{\alpha}$ 都以这种方式嵌入内射对象 $I_{\alpha}$, 则族 $(I_{\alpha})$ 是余生成族.
\end{enumerate}
\end{enonce*}

% END INCLUDED SGA4_English_current_buildable_source_20260906/source/04/drafts/SGA4_IV_LINES_383_536_English.texfrag


% BEGIN INCLUDED SGA4_English_current_buildable_source_20260906/source/04/drafts/SGA4_IV_LINES_537_947_English.texfrag
\section{拓扑斯的例子}\etiquette[2]{IV.sec2}

\setcounter{subsection}{-1}
\subsection{}\etiquette[2.0]{IV.subsec2.0}
本节汇集了相当多典型的拓扑斯例子, 读者想必已在别处遇到过它们; 目的是帮助读者掌握拓扑斯的「瑜伽」. 关于景上拓扑的其他例子 (取自代数几何), 它们同样各自给出一个拓扑斯, 可参看 \SGA 3 IV 6, 以及 (对平展拓扑斯) 本讨论班第 VII 篇. 因为以后除记号或术语问题外, 几乎不会用到本节, 所以把为增进读者的一般认识而补充在这些例子中的陈述的验证, 留给读者作为练习. 本节的所有例子都将在第 4 节中, 通过考察它们对所给资料的函子性依赖而进一步说明; 也将在以后各节中用来阐明关于拓扑斯的一般概念.

\subsection{拓扑空间所对应的拓扑斯}\etiquette[2.1]{IV.subsec2.1}

设 $X$ 为小拓扑空间, $\Ouv(X)$ 为 $X$ 的开子集范畴, 赋予典范拓扑 \sisi{(III 2.9.1).}{(III \hyperref[III.exem1.9.1]{1.9.1}).}
用 $\Top(X)$ 表示 $\Ouv(X)$ 上的 $𝒰$-层拓扑斯. 将 $X$ 上的每个平展拓扑空间 $X'$ 配以 $\Ouv(X)$ 上的层 $U\mapsto \Gamma(X'/U)=\Hom_X(U,X')$, 就得到此拓扑斯与 $X$ 上的平展拓扑空间范畴的等价 [TF]\nde{这里 [TF] 表示文献 [8].}. 我们有时不免会滥用\pageoriginale 语言, 用同一个字母 $X$ 表示由拓扑空间 $X$ 定义的拓扑斯 $\Top(X)$.

当然, 发展拓扑斯理论时主要以这个例子为指引和直观基础. 不过应当注意, 由拓扑空间得到的拓扑斯具有很特殊的性质, 尤其是由于用来描述它们的景 $C=\Ouv(X)$ 中, \emph{一切态射都是单态射} (从而其底范畴归结为一个预序集). 特别地, $C$ 的对象所表的层是终层的子层, 因而\emph{终层的各子层构成所考虑拓扑斯的生成族}. 大多数在代数几何或代数中自然出现的拓扑斯都不具有这个性质, 参见以下例子. 稍加限制, 这个性质刻画了形如 $\Top(X)$ 的拓扑斯 (参见下面的\hyperref[IV.exer7.1.9]{7.1.9}).

容易验证, 将 $X$ 的每个开子集配以它所表的层的映射 $\Ouv(X)\rightarrow \Top(X)$, 是从 $\Ouv(X)$ 到 $\Top(X)$ 的终对象的子对象集的\emph{双射}; 对自然的序结构而言, 此双射甚至是同构, 即它诱导相应范畴的同构. 这提示, 当只在等价意义下知道 $\Top(X)$ 时, 应能在同胚意义下重构拓扑空间 $X$. 我们将在下面 (\hyperref[IV.subsec4.2]{4.2}) 看到, 只须对 $X$ 加一个很小的限制, 情况确实如此.

\subsection{点拓扑斯或终拓扑斯, 与空拓扑斯或始拓扑斯}\etiquette[2.2]{IV.subsec2.2}\pageoriginale

当 $X$ 为只有一点的拓扑空间时, 函子
$$
F\longrightarrow F(X):\Top(X)\longrightarrow 𝒰\mbox{-}\Ens
$$
是范畴等价. 特别地, 这表明\emph{范畴} $𝒰$-$\Ens$ 是 $𝒰$-\emph{拓扑斯}. 我们在第 II 篇的例子中看到, 从正合性质的观点看, 此 $𝒰$-拓扑斯是典型的: 一般拓扑斯的许多性质 (尤其是正合性质) 的验证, 都化归到这个特殊的拓扑斯. 这里用单点拓扑空间给出的解释, 说明把等价于范畴 $𝒰$-$\Ens$ 的拓扑斯称为\emph{点拓扑斯}这一语言滥用是有根据的 (尽管作为范畴, 它完全不等价于单点范畴!). 此术语符合这些拓扑斯所起作用的正确几何直观. 另一种语言滥用是, 有时把点拓扑斯称为\emph{终拓扑斯}, 参见 \hyperref[IV.subsec4.3]{4.3}; 我们用「\emph{那个}终拓扑斯」指拓扑斯 ($𝒰$-$\Ens$).

当 $X$ 为空拓扑空间时, $\Ouv(X)$ 为单点范畴, $\Ouv(X)$ 上的预层 $F$ 为层当且仅当它在 $\Ouv(X)$ 的唯一对象 $\phi$ 处的值为单元素集. 从而 $\Top(∅)$ 同构于单元素 $𝒰$-集的范畴, 后者等价于单点范畴. 特别地由此得到, 单点范畴 (以及与之等价的任意 $𝒰$-范畴) 为 $𝒰$-拓扑斯. 滥用\pageoriginale 语言, 有时称之为\emph{空拓扑斯}或\emph{始拓扑斯} (参见 \hyperref[IV.subsec4.4]{4.4}); 应注意它不等价于空范畴.

\subsection{带算子空间所对应的拓扑斯}\etiquette[2.3]{IV.subsec2.3}

设 $X$ 为拓扑空间, $G$ 为通过同胚作用于 $X$ 的离散群. 在 \sisi{[Toh. \hyperref[IV.subsec5.1]{5.1}]}{[Toh]\nde{这里 [Toh] 表示文献 [4].} \hyperref[IV.subsec5.1]{5.1}} 中定义了 $X$ 上的 $G$-\emph{层}范畴, 我们也将称之为 $(X,G)$ 上的层; 它们是 $X$ 上的 (集合) 层, \emph{配有}与 $G$ 在 $X$ 上的作用相容的 $G$ 作用. 用 Giraud 判别准则 \hyperref[IV.theo1.2]{1.2} iii), %% \XXX ii) est probablement plus facile à vérifier
立即可见此范畴为拓扑斯 (注意始终默认宇宙 $𝒰$), 简记为 $\Top(X,G)$. 当 $G$ 为平凡群时, 就得到 \hyperref[IV.subsec2.1]{2.1} 的例子; 当 $X$ 为单点空间时, 得到 $G$ 左作用于其上的集合 (或 $G$-\emph{集}) 的拓扑斯, 也称为\emph{离散群 $G$ 的分类拓扑斯}, 记为 $B_G$. 容易验证, 拓扑斯 $B_G$ 的终对象 $e$ 的子对象只有 $e$ 与空层 $\phi$; 特别地, 若 $G$ 不是平凡群, 则分类拓扑斯 $B_G$ 的终对象的子对象不构成 $B_G$ 的生成族. 因而此时 $B_G$ 不等价于 \hyperref[IV.subsec2.1]{2.1} 中所考虑的 $\Top(X)$ 类型的拓扑斯.

$G$-层的概念是在该处为了发展交换 $G$-层的上同调理论而引入的. 将后者解释为拓扑斯 $(X,G)$ 的交换层, 该理论就成为在一般拓扑斯框架中发展的第 V 篇理论的一部分.

更\pageoriginale 一般地, 对配有一个自同构拓扑群 $G$ (不必离散) 的拓扑空间 $X$, 可以设法联系出一个适当的拓扑斯, 从而得到合适的上同调理论. 对微分流形、实或复解析流形, 以及概形, 也可以这样做. 这确实是可能的, 参见下面的 \hyperref[IV.subsec2.5]{2.5}.

\subsection{群对象的分类拓扑斯}\etiquette[2.4]{IV.subsec2.4}

设 $E$ 为拓扑斯, $G$ 为 $E$ 的群对象. 设 $(E,G)$ 为 $E$ 中具有 $G$ 作用的对象所成的范畴. 由 Giraud 判别准则立即可知它是拓扑斯. 称它为群对象 $G$ 的\emph{分类拓扑斯}, 记为 $B_G$. 当 $E$ 为点拓扑斯 (\hyperref[IV.subsec2.2]{2.2}), 即 $G$ 为通常的群时, 就得到 \hyperref[IV.subsec2.3]{2.3} 的分类拓扑斯.

所采用术语的根据在于, 拓扑斯 $B_G$ 在 $G$ 下的「挠子」(或齐性主丛) 的分类中具有泛性质; 更一般地, 对 $G_{E'}=f^{\ast}(G)$ 下的挠子也是如此, 其中 $E'$ 为「$E$ 上」的拓扑斯, 即赋有态射 $f:E'\rightarrow E$ (参见下面的 \hyperref[IV.def3.1]{3.1}). [3 Chap V] 或下面的 \hyperref[IV.exer5.9]{5.9} 中明确说明的这个作用, 表明在拓扑斯的框架中, $B_G$ 所起的作用, 与拓扑空间同伦理论中拓扑群的经典分类空间的作用相同. 后者可以看作 (参见 \hyperref[IV.subsec2.5]{2.5}) 前者的弱化版本, 它只保留分类拓扑斯的「同伦型」, 这里取某个适当的意义, 无须详述.

\subsection{拓扑空间的「大景」与「大拓扑斯」. 拓扑群的分类拓扑斯}\footnote[1]{这些景与拓扑斯是 M. \textsc{Giraud} 引入的, 他还指出了它们优于传统「小」景之处.}\etiquette[2.5]{IV.subsec2.5}\pageoriginale

设 $U-\Esp$, 或简记 $(\Esp)$, 为 $\in 𝒰$ 的拓扑空间的范畴. 已知有限投射极限在 $(\Esp)$ 中可表. 在 $(\Esp)$ 上考虑如下预拓扑 (\hyperrefext[II][II.def1.3]{1.3}): $\Cov(X)$ 为开嵌入的满射族 $u_i:X_i\rightarrow X$ 的集合. 通过上述预拓扑生成的拓扑, 把 $(\Esp)$ 看作景. 对 $(\Esp)$ 的每个对象 $X$, 将 $X$ 上的 $(\Esp)$ 的对象, 即 $X$ 上的拓扑空间所成的范畴
$$
(\Esp)_{/X}
$$
看作景, 其拓扑由遗忘函子 $(\Esp)_{/X}\rightarrow (\Esp)$ 从 $(\Esp)$ 的拓扑诱导而来 (\hyperrefext[III][III.prop5.2]{5.2} 4)). 称此景为与 $X$ 相联系的\emph{大景}. 应注意它不 $\in 𝒰$; 它也不是 \hyperrefext[II][II.def3.0.2]{3.0.2} 意义下的 $𝒰$-景, 所以对它应用通常结果时须谨慎. 为消除此不便, 可以选取宇宙 $𝒱$ 使得 $𝒰\in 𝒱$, 从而 $(\Esp)_{/X}$ 成为 $𝒱$-景, 并使用相应的 $𝒱$-拓扑斯 $(\Esp)_{/X}$, 可记为 $\TOP(X)$, 称为 \emph{$X$ 的大拓扑斯}. 若不愿扩大 $𝒰$, 可选取一个基数 $c$, 为 $X$ 及将在论证中使用的一切拓扑空间的势的界 (通常 \sisi{$\Sup(\card X, \card ℛ)$}{$\Sup(\card X, \card \mathbb{R})$} 就足够!), 并用由 $X$ 上满足 $\card X'\leq c$ 的各 $X'$ 组成的子范畴 $(\Esp)'_{/X}$ 代替 $(\Esp)_{/X}$, 赋予诱导拓扑; 此时用 $\TOP(X)$ 表示这个\pageoriginale 景上的层拓扑斯. 为确定概念, 假设采用第一种定义.

$X$ 的大拓扑斯相对于小拓扑斯的优势在于, 定义它的景包含 $(\Esp)_{/X}$ 作为满子范畴; 因为此景的拓扑显然比典范拓扑粗, 所以从 $(\Esp)_{/X}$ 到 $\TOP(X)$ 的典范函子, 即将 $X$ 上的每个空间 $X'$ 配以它所表的层的函子, 是\emph{全忠实}的. 因此, \emph{$X$ 上的空间 $X'$ 由它所定义的层 $(\in \Top(X))$ 在 $X$-同构的意义下确定}; 从而 ($X$ 的大景上) 的层的概念可以看作 $X$ 上的拓扑空间概念的\emph{推广}, 借此层论的一切构造对 $X$ 上的拓扑空间都有了意义.

因此, 当 $G$ 为 $X$ 上的拓扑空间范畴 $(\Esp)_{/X}$ 的群对象时, 可联系出它的分类拓扑斯 $B_G$ (\hyperref[IV.subsec2.4]{2.4}), 因而得到分类上同调群、分类同伦群等等 (定义为 $𝒱$-拓扑斯 $B_G$ 的相应不变量). 特别地, 当 $X$ 为单点空间时, $G$ 等同于通常的拓扑群. 在保证各笛卡儿积 $G^n$ 的奇异上同调与层意义下的上同调 (例如取常值系数) 一致的通常局部条件下, 例如 $G$ 局部可缩时, 可以验证 $G$ 的分类拓扑斯的上同调, 与拓扑学家意义下 $G$ 的分类空间的上同调典范同构.

通过「大景」引入\pageoriginale 分类拓扑斯, 比分类空间更有利之处在于提供了更丰富的理论, 尤其是对比常值或局部常值系数更一般的系数, 它们也提供有用的上同调不变量. 此外, 这里所考虑的定义显然适用于其他常见的框架: 微分流形、实或复解析流形或空间, 以及概形. 尤其是把「经典群」看作来自整数环上定义的概形, 这个观点就能联系起从传统观点与「算术」观点出发的示性类研究; 关于这一方向的提示, 参见 [\hyperref[IV.sec7]{7}]. 同样, J. \textsc{Giraud} [\hyperref[IV.sec3]{3}] 在十分一般而灵活的拓扑斯框架中, 关于群对象扩张的分类的一般结果, 可借助「大拓扑斯」特殊化为关于拓扑群, 或实或复 Lie 群的扩张分类的结果; 除了核交换的扩张情形 [\hyperref[IV.sec11]{11}] 以外, 拓扑学家似乎几乎不了解这些结果.

\subsection{形如 $\widehat{C}$ 的拓扑斯}\etiquette[2.6]{IV.subsec2.6}

设 $C$ 为小范畴. 则 $C$ 上的 $U$-预层范畴 $\widehat{C}$ 显然是 $𝒰$-拓扑斯, 因为它形如 $C^{\widetilde{\phantom{,}}}$, 其中 $C$ 赋予混沌拓扑. 下面将详细说明 $C$ 与 $\widehat{C}$ 之间的关系. 这里只指出, 等价于形如 $\widehat{C}$ 的拓扑斯的拓扑斯 $E$ 性质相当特殊, 因为\emph{它有一个由连通投射对象组成的小生成族 $(X_i)$}, 即由这样的对象 $X$ 组成: 函子 $Y\mapsto \Hom(X,Y)$\pageoriginale 将满态射映为满态射, 将余积映为余积; 事实上, 在 $\widehat{C}$ 中取由 $X\in \ob C$ 所表的函子组成的生成族即可. 此外注意, 若拓扑斯 $E$ 中有一个覆盖族 $X_i\rightarrow X$, 其中各 $X_i$ 投射且连通, 则 $X$ 的其他每个覆盖族都被此覆盖族加细 (\hyperrefext[I][I.sec4.3.2]{4.3.2}, \hyperrefext[I][I.sec4.3.3]{4.3.3}). 因而在形如 $\widehat{C}$ 的拓扑斯 $E$ 中, 每个对象 $X$ 都有一个加细所有其他覆盖族的覆盖族. 形如 $\Top(X)$ 的拓扑斯 (\hyperref[IV.subsec2.1]{2.1}), 其中 $X$ 为点都是闭的拓扑空间, 只有在 $X$ 离散时才有上述性质.

当范畴 $C$ 只有一个对象时, $C$ 等同于幺半群 $G$. 此时 $C$ 上的预层等同于 $G$ \emph{右作用}于其上的集合 (因为它是函子 $G^{\circ}\rightarrow (\Ens)$), 而 $\widehat{C}$ 是带有右算子幺半群的集合的拓扑斯, 根据 \hyperref[IV.subsec2.3]{2.3} 也可记为 $B_{G^{\circ}}$: 它是\emph{以幺半群 $G^{\circ}$ ($G$ 的反幺半群) 为算子的集合的拓扑斯}. 当 $G$ 为群时, 用从 $G$ 到 $G^{\circ}$ 的同构 $g\mapsto g^{-1}$, 就得到 \hyperref[IV.subsec2.3]{2.3} 的分类拓扑斯 $B_G$.

\subsection{pro-群的分类拓扑斯}\etiquette[2.7]{IV.subsec2.7}

\subsubsection{}\etiquette[2.7.1]{IV.subsubsec2.7.1}
设 $𝒢=(G_i)_{i\in I}$ 为群的投射系, 其中 $G_i$, $I\in 𝒰$. 假定此投射系\emph{严格}, 即过渡态射 $G_j\rightarrow G_i$ 都为满射. 若 $E$ 为集合, 则 $𝒢$ \emph{在 $E$ 上的作用} (例如左作用) 指下述结构: a) $E$ 的子集族 $(E_i)_{i\in I}$, 其并为 $E$; b) 对每个 $i\in I$, 群 $G_i$ 在集合 $E_i$ 上的作用; 此外要求这些资料满足\pageoriginale 以下条件: 对 $j\geq i$, $E_i$ 是 $E_j$ 中在 $G_j\rightarrow G_i$ 的核的作用下不动的元素所成的子集. 若给定了 $𝒢$ 在 $E$ 上的 (左) 作用, 也说 $𝒢$ (从左) \emph{作用于} $E$. 赋有 $𝒢$ 作用的 $\in 𝒰$ 的集合, 以显然的方式构成范畴. 由 Giraud 判别准则立即可见此范畴为 $𝒰$-拓扑斯. 记为 $B_{𝒢}$, 称为 \emph{$𝒢$ 的分类拓扑斯}. 当 $I$ 有始对象 $i_{\circ}$ 时, 置 $G=G_{i_\circ}$, 就得到 \hyperref[IV.subsec2.3]{2.3} 的分类拓扑斯.

\subsubsection{}\etiquette[2.7.2]{IV.subsubsec2.7.2}
另一个重要例子是各群 $G_i$ 有限的情形, 此时
$$
G=\varprojlim G_i
$$
为紧的全不连通拓扑群, 或称\emph{pro-有限群}. 此时 $𝒢$ 在 $E$ 上的作用相当于 $G$ 在 $E$ 上的连续作用, 或等价地, $E$ 的每个点的稳定子都是 $G$ 的开子群的作用. 分类拓扑斯 $B_{𝒢}$ 也记为 $B_{G}$, 其中当然应把 $G$ 看作赋有 pro-有限拓扑的群.

\subsubsection{}\etiquette[2.7.3]{IV.subsubsec2.7.3}
用 \hyperref[IV.subsec2.6]{2.6} 的说明, 容易验证, 由群的严格投射系 $𝒢=(G_i)_{i\in I}$ 定义的拓扑斯 $B_{𝒢}$, 只有在此投射系本质常值时才等价于形如 $\widehat{C}$ 的拓扑斯; 对 pro-有限群的情形, 这就是说该群实际上有限.

\subsubsection{}\etiquette[2.7.4]{IV.subsubsec2.7.4}
离散群 $G$ 的分类拓扑斯 $B_G$ 的以下几何解释, 有助于形成关于这些拓扑斯的正确几何直观. (同一思路还可参见下面的 \hyperref[IV.subsec4.5]{4.5}, \hyperref[IV.subsec5.8]{5.8}, \hyperref[IV.exer5.9]{5.9}, 及 \hyperref[IV.subsec7.2]{7.2}.)
设 $X$ 为连通、局部连通且\pageoriginale 局部单连通的拓扑空间, $x$ 为 $X$ 的点, $G$ 为它在 $x$ 处的基本群. (注意, 已知在同构意义下, 每个离散群 $G$ 都可这样得到.) 则 $X$ 的覆盖的 Galois 理论给出 $G$-集范畴 $B_G$ 与 $X$ 的\emph{平展覆盖}范畴之间的等价; 后者即 $X$ 上这样的空间 $X'$ 所成的范畴: 它局部地 $X$-同构于形如 $X\times I$ 的 $X$-空间, 其中 $I$ 为离散空间. (比较 \SGA 1 V 4,5.) 此外, 后一范畴可解释为 $X$ 上的局部常值层范畴, 即拓扑斯 $\Top(X)$ 的局部常值对象 (IX \hyperref[IX.subsec2.0]{2.0}) 的范畴.

当 $G$ 为 pro-有限群时, $B_G$ 也有类似的几何解释: 它是连通概形 $X$ 的有限平展覆盖的余积所成的 $X$-概形范畴, 其中 $X$ 配有几何点 $x$ 以及同构 $G\simeq \pi_1(X,x)$. 也已知每个 pro-有限群都可由适当的连通概形 (若愿意, 可取域的谱) 的基本群得到. 最后, pro-群 (不一定 pro-有限或本质常值) 也出现在连通、局部连通而非局部单连通空间的覆盖分类中, 以及连通非正规概形的不一定有限或 ind-有限的平展覆盖分类中. 后一情形参见 \SGA 3 X 6.

\begin{enonce*}[remark]{练习 2.7.5}\etiquette[2.7.5]{IV.exer2.7.5}
对拓扑斯 $E$, 定义连通、局部连通~\footnote[1]{参见 \sisi{\hyperref[IV.exer8.7]{8.7} 1).}{\hyperref[IV.exer8.7]{8.7} l).}}、单连通及局部单连通的概念. 定义 $E$ 的常值对象及局部常值对象的概念 (参见 IX \hyperref[IX.subsec2.0]{2.0}). 设 $f:E'\rightarrow E$\pageoriginale 为拓扑斯态射 (\hyperref[IV.def3.1]{3.1}), 其中 $E'$ 单连通 (例如 $E'$ 为点拓扑斯 (\hyperref[IV.subsec2.2]{2.2})), $E$ 连通且局部连通. 定义一个严格 pro-群 $\pi_1(E,f)=(G_i)_{i\in I}=𝒢$ (称为 \emph{$E$ 在 $f$ 处的基本 pro-群}), 以及 $B_{𝒢}$ 与 $E$ 的局部常值对象范畴之间的范畴等价.\footnote[2]{可参考 \SGA 3 X 6.} 证明当 $E$ 局部单连通时, $\pi_1(E,f)$ 本质常值, 从而等同于通常的离散群 $\pi_1(E,f)$, 称为 \emph{$E$ 在 $f$ 处的基本群}. 证明每个严格 pro-群 $𝒢$ 都同构 (作为 pro-群) 于适当的连通且局部连通拓扑斯在适当的 $f$ 处的基本 pro-群, 其中 $E'$ 为点拓扑斯 (取 $E=B_{𝒢}$, 并由遗忘函子 $f^{\ast}:E\rightarrow E'=(\Ens)$ 定义 $f:E'\rightarrow E$). 当 $𝒢$ 本质常值, 即同构 (作为 pro-群) 于通常的离散群时, 证明上面的 $E$ 可取为局部单连通的 (仍取 $E=B_G$).
\end{enonce*}

\subsection{伪拓扑斯的例子}\etiquette[2.8]{IV.subsec2.8}

设 $𝒢=(G_i)_{i\in I}$ 为严格 pro-群, 其中 $I$ 为滤过有序集, 且 $i>j$ 蕴涵 $G_i\rightarrow G_j$ 不是同构. 假设 $\card(I)\not\in 𝒰$. 考虑 $𝒢$ 从左作用于其上的集合 $E\in 𝒰$ 的范畴 (\hyperref[IV.subsubsec2.7.1]{2.7.1}). 它是 $𝒰$-范畴, 并且如同 \hyperref[IV.subsubsec2.7.1]{2.7.1}, 可见此范畴满足 \hyperref[IV.subsubsec1.1.2]{1.1.2} 的条件 a), b) 和 c). 然而它不是 $𝒰$-拓扑斯, 因为它没有 $𝒰$-小的生成族. 同样, 对任何宇宙 $𝒱$, 它都不是 $𝒱$-拓扑斯.

% END INCLUDED SGA4_English_current_buildable_source_20260906/source/04/drafts/SGA4_IV_LINES_537_947_English.texfrag


% BEGIN INCLUDED SGA4_English_current_buildable_source_20260906/source/04/drafts/SGA4_IV_LINES_948_1319_English.texfrag
\section{拓扑斯态射}\etiquette[3]{IV.sec3}\pageoriginale

\begin{enonce*}{定义 3.1}\etiquette[3.1]{IV.def3.1}
设 $E$ 与 $E'$ 为两个 $𝒰$-拓扑斯. 从 $E$ 到 $E'$ 的态射\nde{有些作者称为\emph{几何态射}.}, 有时也 (滥用语言) 称为从 $E$ 到 $E'$ 的连续映射, 是一个三元组 $u=(u_{\ast}, u^{\ast},\varphi)$, 由函子
$$
u_{\ast}:E\longrightarrow E'\quav u^{\ast}:E'\longrightarrow E
$$
以及关于 $X'\in \Ob E'$, $Y\in \Ob E$ 的双函子的「伴随」同构 $\varphi$
$$
\varphi:\Hom_E (u^{\ast}(X'),Y)\xrightarrow{\sim} \Hom_{E'}(X',
u_{\ast}(Y)){\quoi},
$$
组成, 此外要求函子 $u^{\ast}$ 左正合, 即与有限投射极限交换. 函子 $u_{\ast}$ 称为拓扑斯态射 $u$ 的正像函子, 函子 $u^{\ast}$ 称为拓扑斯态射 $u$ 的逆像函子, 同构 $\varphi$ 称为 $u$ 的伴随同构.
\end{enonce*}

\setcounter{subsection}{1}
\subsubsection{}\etiquette[3.1.1]{IV.subsubsec3.1.1}
除非另有明说, 对拓扑斯态射 $u:E\rightarrow E'$, 今后都用 $u_{\ast}$ 与 $u^{\ast}$ 表示相应的正像和逆像函子\footnote[1]{有时也宜用 $u^{-1}$ 代替 $u^{\ast}$, 参见 \hyperref[IV.subsubsec13.2.3]{13.2.3}.}. 注意, 由伴随同构 $\varphi$, $u_{\ast}$ 是 $u^{\ast}$ 的右伴随, 而 $u^{\ast}$ 是 $u_{\ast}$ 的左伴随; 因而由伴随函子的熟知一般事实 [\hyperref[IV.sec14]{14}], 两个函子 $u_{\ast}$, $u^{\ast}$ 中的每一个都在唯一同构的意义下确定另一个. 实际上, 依情形而定, 指定 $u^{\ast}:E'\rightarrow E$ 或指定 $u_{\ast}:E\rightarrow E'$ 中的一种方式, 可能更便于定义拓扑斯态射 $u:E\rightarrow E'$; 在第一种情形, 只须验证给定函子 $u^{\ast}$ 有\pageoriginale 右伴随且左正合. 在第二种情形, 须验证给定函子 $u_{\ast}$ 有一个左正合的左伴随. 在两种情形中, 从部分资料出发, \emph{选取}一个伴随函子及一个伴随同构, 就得到拓扑斯态射 $u:E\rightarrow E'$; 而由资料 $u^{\ast}$, 分别由 $u_{\ast}$, 后者在一个足够清楚、并将在下面\hyperlink{sga4.exp.IV.subsubsection.3.2.1}{(3.2.1)} 中完全说清的意义下, 将是「在唯一同构的意义下唯一的」.

\subsubsection{}\etiquette[3.1.2]{IV.subsubsec3.1.2}
若 $u:E\rightarrow E'$ 为拓扑斯态射, 则由伴随函子的性质 (\hyperrefext[I][I.prop2.11]{2.11}) 可知, 正像函子 $u_{\ast}:E\rightarrow E'$ 与投射极限交换, 函子 $u^{\ast}:E'\rightarrow E$ 与归纳极限交换\footnote[2]{此外 (\hyperref[IV.cor1.5]{1.5} 及 \hyperref[IV.rem1.8]{1.8}), 对给定函子 $u_{\ast}:E\rightarrow E'$, 分别对 $u^{\ast}:E'\rightarrow E$, 此函子有左伴随 (分别有右伴随) 当且仅当它与 $𝒰$-投射极限 (分别与 $𝒰$-归纳极限) 交换.}\MakeLinkTarget*{sga4.exp.IV.footnote.3.1.1.2}\label{sga4[1]}. 因为还假定后者左正合, 即与有限投射极限交换, 特别地可知 $u^{\ast}$ \emph{正合}. 因而可以说, 在对 $(u_{\ast},u^{\ast})$ 中, \emph{逆像函子 $u^{\ast}$ 具有最显著的正合性质}. 这些性质保证: 对每种代数结构类型 $\Sigma$, 如果其资料可用底集之间以及从底集反复作有限投射极限与任意归纳极限所得到的集合之间的箭头来描述, 则每个「赋有 $\Sigma$ 型结构的 $ℰ'$ 的对象」(借助拓扑斯 $ℰ'$ 的内在正合性质 (\hyperrefext[II][II.theo4.1]{4.1}), 这个概念是有意义的), 其在 $u^{\ast}$ 下的像都赋有同样的结构. 与其承担精确说明这句话的意义并形式地证明它这一费力不讨好的工作,\pageoriginale 不如建议读者对群、环、环上的模、余环上的余模、环上的双代数以及群下的挠子等结构类型, 亲自将它明确写出并确信其有效. (这些例子中, 前三种结构类型完全用有限投射极限定义, 而其他概念隐含涉及了也使用归纳极限的构造.) 此外, 函子 $u^{\ast}$ 与对这些结构所作的一切通常函子性运算「\emph{交换}」; 更确切地说, 它与一切可用 $\varinjlim$ 和有限 $\varprojlim$ 表达的运算交换: 由一个对象生成的自由对象的构造 (例如自由群或自由模)、张量积 (参见下面的 \hyperref[IV.sec12]{\S 12}) 等等.

至于正像函子 $u_{\ast}:E\rightarrow E'$, 由于它与投射极限交换, 就「保持」$E$ 的一个对象 (或对象族) 上的每种完全用投射极限定义的代数结构 (例如前述例子中的群、环、环上的模的结构). 另一方面, 函子 $u_{\ast}$ 一般不右正合, 即一般不与有限归纳极限交换, 甚至一般不把满态射映为满态射 (此外, 正是 $u_{\ast}$ 正合性的这种缺失, 成为了它的上同调性质的来源; 下一篇将从交换同调代数的观点研究这些性质). 因而它一般不能扩充为余模、\pageoriginale 双代数或群下的挠子类型的对象上的函子, 也一般不与「由……生成的自由模」、模的张量积等运算交换.

\subsubsection{}\etiquette[3.1.3]{IV.subsubsec3.1.3}
实际上, 当给定一个从 $𝒰$-拓扑斯到另一个 $𝒰$-拓扑斯的函子 $f:E\rightarrow F$ 时, 应当明确它的一切正合性质, 包括可能存在的左或右伴随函子 (参见第 \pageref{sga4[1]} 页脚注), 以了解 $f$ 的「几何本性」; 这种了解一般是形成关于该情形的正确几何直观所必不可少的指引. 例如, 若发现 $f$ 与任意归纳极限和有限投射极限交换, 则应将 $f$ 写为
$$
f=u^{\ast}{\quoi},
$$
其中
$$
u:F\longrightarrow E
$$
为拓扑斯态射, 即把 $f$ 解释为拓扑斯的一个「连续映射」的「逆像」函子. 当 $f$ 与任意投射极限交换, 因而有左伴随, 且若后者 (它先验地与任意归纳极限交换) \emph{还}与有限投射极限交换时, 就应把 $f$ 写为
$$
f=v_{\ast}{\quoi},
$$
其中
$$
v:E\longrightarrow F
$$
为拓扑斯态射. 某些情形下 $f$ 可能同时满足\pageoriginale 所考虑的两种性质 (例子见 \hyperref[IV.subsec4.10]{4.10}). 此时应同时引入拓扑斯态射
$$
u:F\longrightarrow E \quav v:E\longrightarrow F{\quoi},
$$
它们给出三个伴随函子的序列 (\hyperrefext[I][I.sec5.3]{5.3}):
$$
v^{\ast} \quav v_{\ast}=u^{\ast}=f\quav u_{\ast}{\quoi}.
$$
此时必须仔细区分拓扑斯态射 $u$ 与 $v$, 以免失去关于该情形的几何直观.

与此相关, 注意当两个拓扑斯 $E$, $F$ 之间有三个伴随函子的序列
$$
e, f, g\; (e,g: F\rightrightarrows E, f:E\rightarrow F){\quoi},
$$
则 $f$ 与任意归纳极限及投射极限交换, 因而总可写为 $u^{\ast}$, 其中 $u:F\rightarrow E$ 为拓扑斯态射. 因而 $g$ 可写为 $g=u_{\ast}$. 另一方面, 当然只有当 $e$ 也与有限投射极限交换时, 才可将 $e$ 写为 $v^{\ast}$ (此时 $f$ 写为 $v_{\ast}$). 若 $e$ 本身是第四个函子 $d$ 的右伴随, 情况显然如此. 如果没有这样的条件, 当逆像函子 $f=u^{\ast}$ 有左伴随时, 常用
$$
e=u_{!}
$$
表示它的左伴随; 此记号由拓扑空间的开浸入 $u:X\rightarrow Y$ 的例子所启发. 同样, 若 $e$ 形如 $v^{\ast}$, 即 $f$ 形如 $v_{\ast}$, 则正像函子 $v_{\ast}$ 有右伴随时, 有时用 $g=v^{!}$ 表示它\footnote[1]{交换层的情形见下面的 \hyperref[IV.subsec14.4]{14.4}.}.

\subsection{}\etiquette[3.2]{IV.subsec3.2}
设\pageoriginale $E$, $E'$ 为两个 $𝒰$-拓扑斯, 并设
$$
u=(u_{\ast},u^{\ast},\varphi)\quav
v=(v_{\ast},v^{\ast},\Psi):E\rightrightarrows E'
$$
为从 $E$ 到 $E'$ 的两个拓扑斯态射. 从 $u$ 到 $v$ 的\emph{态射}, 就是从 $u_{\ast}$ 到 $v_{\ast}$ 的任意态射 (在从 $E$ 到 $E'$ 的函子范畴 $\Hhom(E,E')$ 的意义下). 拓扑斯态射的态射以显然的方式复合, 从而定义一个范畴, 事实上是一个 $\sheaf{\mathcal{U}}$-\emph{范畴} (\hyperrefext[I][I.prop7.8]{7.8}), 记为
$$
\Hhomtop(E,E'){\quoi},
$$
称为\emph{从 $E$ 到 $E'$ 的态射范畴} (或\emph{连续映射范畴}). 于是有一个显然的函子
$$
u\longmapsto u_{\ast}:\Hhomtop(E,E')\longrightarrow
\Hhom(E,E'){\quoi},
$$
根据左边范畴中箭头的定义, 该函子全忠实 (但它在对象上不是单射).

\hypertarget{sga4.exp.IV.subsubsection.3.2.1}{}\subsubsection{}
须注意, 若如上给定 $u$ 与 $v$, 则伴随函子理论给出典范双射
$$
\Hom(u_{\ast},v_{\ast})\xrightarrow{\sim}\Hom(v^{\ast},
u^{\ast}){\quoi};
$$
特别地, 得到 $\Hhomtop(E,E')$ 上的\emph{反变函子}:
$$
u\longmapsto u^{\ast}:\Hhomtop(E,E')^{\circ}\longrightarrow
\Hhom(E,E'){\quoi}.
$$
按几何直观, 正像函子 $u_{\ast}$ 与产生它的连续映射「方向相同」, 因此应当\emph{根据正像函子}定义拓扑斯态射之间态射的箭头方向, 而不根据逆像函子 (尽管正如已经看到的, 后者具有拓扑斯态射概念的特征性的正合性质).

\subsubsection{}\etiquette[3.2.2]{IV.subsubsec3.2.2}
定义了从拓扑斯 $E$ 到拓扑斯 $E'$ 的态射范畴 $\Hhomtop(E,E')$ 之后\pageoriginale, 也就定义了从 $E$ 到 $E'$ 的两个态射 $u$, $v$ 之间\emph{同构}的概念. 实际上, 没有必要在本质上区分两个同构的拓扑斯态射, 至少当有一个两者之间的\emph{典范}同构时如此 (正如当给定一个范畴的两个对象之间的典范同构时, 常无须区分它们). 与此相关, 指出在讨论拓扑斯态射的图表及此类图表的交换性问题 (由\hyperlink{sga4.exp.IV.subsection.3.3}{(3.3)}, 这个概念有意义) 时, 大多数时候只指在 (「典范」) 同构的意义下交换; 此时滥用语言, 把这些图表当作实际交换的图表.

可能有人试图引入在同构意义下从 $E$ 到 $E'$ 的态射所成的集合 \\ $\Homtop(E,E')/\Isom$, 并把这样的同构类称为态射, 而不遵循定义 \hyperref[IV.subsec3.2]{3.2}, 来为这种语言滥用辩护. 但这就会遇到十分严重的困难, 每当没有典范同构却要把一个范畴的两个同构对象等同时, 都会出现这些困难. 经验表明, 这种观点行不通, 必须保留拓扑斯态射之间态射的「精细」概念 \hyperref[IV.def3.1]{3.1}, 即使有时不得不为典范同构之间的相容性费心\footnote[1]{关于这类费心的例子 (看来成功了), 请参看 M. \textsc{Hakim} 女士关于相对概形的著作 [9].}.

\hypertarget{sga4.exp.IV.subsection.3.3}{}\setcounter{subsection}{3} \setcounter{subsubsection}{0}
\subsubsection{}\etiquette[3.3.1]{IV.subsubsec3.3.1}
设\pageoriginale $E$, $E'$, $E''$ 为三个 $𝒰$-拓扑斯, 考虑拓扑斯态射
$$
u:E\longrightarrow E'\quav v:E'\longrightarrow E''{\quoi}.
$$
此时伴随函子理论根据对 $(u_{\ast},u^{\ast})$ 与 $(v_{\ast},v^{\ast})$ 的伴随同构, 给出复合函子 $v_{\ast}u_{\ast}$ 与 $u^{\ast}v^{\ast}$ 之间的伴随同构. 另一方面, 函子 $u^{\ast}v^{\ast}$ 为两个左正合函子的复合, 因而左正合. 所以得到一个从 $E$ 到 $E''$ 的态射, 称为\emph{态射 $u$ 与 $v$ 的复合}, 记为 $vu$:
$$
vu:E\longrightarrow E''{\quoi}.
$$
立即可验证, 态射的复合满足结合律, 且对每个 $𝒰$-拓扑斯 $E$, 都有一个从 $E$ 到自身的态射, 它是复合运算的双侧单位元: 这就是态射 $(\id_E,\id_E,\varphi)$, 其中 $\varphi$ 为 $\id_E$ 与自身的显然伴随同构. 设 $𝒱$ 为满足 $𝒰\in 𝒱$ 的宇宙. 定义范畴
$$
(𝒱\text{-}𝒰\text{-}\Top){\quoi},
$$
其对象为属于 $𝒱$ 的 $𝒰$-拓扑斯, 箭头为这些 $𝒰$-拓扑斯之间的态射, 箭头的复合即刚才描述的运算.

\subsubsection{}\etiquette[3.3.2]{IV.subsubsec3.3.2}
事实上, 复合映射
$$
\Homtop(E, E')\times \Homtop(E',E'')\longrightarrow \Homtop(E,E'')
$$
是一个「态射复合函子」在对象上诱导的映射:
$$
\Hhomtop(E,E')\times \Hhomtop(E',E'')\longrightarrow
\Hhomtop(E,E''){\quoi},
$$
它在箭头上的\pageoriginale 作用为函子 (这里是正像函子) 之间态射的通常「卷积积」运算. 对四个拓扑斯 $E$, $E'$, $E''$, $E'''$, 这些复合函子满足一个 (严格的) 结合性质, 它细化了拓扑斯态射复合的结合性. 也可以用开始为人熟悉的语言 [2], [9] 说, $𝒰$-拓扑斯是一个 2-\emph{范畴}的对象 (或 0-箭头), 其 1-箭头为拓扑斯态射, 2-箭头为拓扑斯态射的态射.

$𝒰$-拓扑斯 (作为宇宙 $𝒱$ 的元素) 构成 2-范畴, 而不只是像通常拓扑空间那样构成普通范畴, 从技术观点看, 这正是拓扑斯理论与拓扑空间理论最重要的区别. 这个事实既是前已提及的某些技术复杂性的来源, 也是相对于传统拓扑学而言某些本质上新颖的事实的来源.

\subsection{}\etiquette[3.4]{IV.subsec3.4}
$𝒰$-拓扑斯 (作为宇宙 $𝒱$ 的元素) 构成 2-范畴 (\hyperref[IV.subsubsec3.3.2]{3.3.2}), 因而特别可以定义\emph{两个 $𝒰$-拓扑斯 $E$, $E'$ 的等价}的概念: 若存在拓扑斯态射 $u:E\rightarrow E'$ 与 $v:E'\rightarrow E$, 使得复合 $uv$ 与 $vu$ 分别同构于 $E$ 与 $E'$ 的恒等态射, 就说 $E$ 与 $E'$ \emph{等价}; 此时称态射 $u$ 与 $v$ \pageoriginale \emph{互为拟逆等价}.

立即可见, 拓扑斯态射 $u:E\rightarrow E'$ 为等价的充分必要条件是 $u_{\ast}$ 为等价, 或者等价地, $u^{\ast}$ 为等价. (用以下事实: 两个拓扑斯之间的一个等价函子 $f:E\rightarrow E'$, 既形如 $u_{\ast}$, 也形如 $v^{\ast}$, 其中 $u$ 与 $v$ 为拓扑斯态射); 而 $E$ 与 $E'$ 在前一段意义下等价的充分必要条件是, 它们作为范畴等价 (即作为 2-范畴 $(𝒱\text{-}\cat)$ 的对象等价). 正如应有的那样, 这表明所引入的等价概念与 \hyperref[IV.subsubsec3.3.1]{3.3.1} 中宇宙 $𝒱$ 的选取无关.

\subsubsection{}\etiquette[3.4.1]{IV.subsubsec3.4.1}
实际上, 大多数时候没有理由在本质上区分等价的 $𝒰$-拓扑斯, 正如常常没有理由在本质上区分两个等价范畴---不过前提是有一个从其中一个到另一个的明确等价, 或至少有一个在唯一同构意义下定义的等价. 在这里, 拓扑斯等价的概念替代了两个拓扑空间之间同胚的传统概念. 关于两个概念的确切关系, 见下面 \hyperref[IV.subsec4.2]{4.2} 的例子.

% END INCLUDED SGA4_English_current_buildable_source_20260906/source/04/drafts/SGA4_IV_LINES_948_1319_English.texfrag


% BEGIN INCLUDED SGA4_English_current_buildable_source_20260906/source/04/drafts/SGA4_IV_LINES_1320_1659_English.texfrag
\section{拓扑斯态射的例子}\etiquette[4]{IV.sec4}

这里利用拓扑斯态射的概念重新考察第 2 节的例子. \hyperref[IV.subsec2.0]{2.0} 的说明同样适用于本节. 一般默认宇宙 $𝒰$.

\subsection{随拓扑空间 $X$ 变化的拓扑斯 $\Top(X)$}
\etiquette[4.1]{IV.subsec4.1}\pageoriginale

\subsubsection{}\etiquette[4.1.1]{IV.subsubsec4.1.1}
设
$$
f:X\longrightarrow Y
$$
为拓扑空间的连续映射. 我们将典范地联系出一个拓扑斯态射
$$
\Top(f)\mbox{ 或 } f:\Top(X)\rightarrow \Top(Y){\quoi},
$$
其中使用 \hyperref[IV.subsec2.1]{2.1} 的记号. 当将 $\Top(X)$ 定义为 $\Ouv(X)^{\widetilde{\phantom{,}}}$ 时, 最方便的是用层的正像函子
$$
f_{\ast}:\Top(X)\longrightarrow \Top(Y){\quoi},
$$
描述 $\Top(f)$, 该函子由公式
$$
f_{\ast}(F)=F\circ f^{-1}{\quoi},
$$
定义, 其中
$$
f^{-1}:\Ouv(Y)\longrightarrow \Ouv(X)
$$
为显然的函子 $U\rightarrow f^{-1}(U)$. 我们已经注意到此函子连续且左正合, 因而确实定义了上述函子 $f_{\ast}$, 它有左正合的左伴随 $f^{\ast}$ (\hyperrefext[III][III.exem1.9.1]{1.9.1}). 当然, 严格地说, 拓扑斯态射 $\Top(f)$ 依赖于 $f_{\ast}$ 的左伴随 $f^{\ast}$ 的选取, 因而只在典范同构的意义下定义. 以下将不再明确指出此类现象.

当采用「平展空间」的观点定义 $\Top(X)$ 时, 最方便的是用逆像函子
$$
f^{\ast}:\Top(Y)\longrightarrow \Top(X)
$$
定义拓扑斯态射 $\Top(f)$,\pageoriginale 只须对 $Y$ 上的每个平展空间 $Y'$, 置
$$
f^{\ast}(Y')=\fprod{X}{Y'}{Y}
$$
即可; 显然该纤维积确是 $X$ 上的平展空间, 如此得到的函子 $f^{\ast}$ 左正合且与任意 $\varinjlim$ 交换, 因而定义拓扑斯态射 $\Top(f)$. 所得到的两种定义之间的相容性, 参见 [TF].

对两个可复合的连续映射
$$
X\xrightarrow{f}Y\xrightarrow{g}Z{\quoi},
$$
有拓扑斯态射的典范同构 $\Top(gf)\simeq\Top(g)\Top(f)$. 对三个可复合的连续映射 $f$, $g$, $h$, 这些\emph{传递同构}满足一个相容关系, 这里不再写出, 它恰是 \SGA 1 VI 7.4 B) 中所考虑的关系 (取 $\epsilon = (\Esp)^{\circ}$). 可以表述为: 当 $X$ 在范畴 $(\Esp)$ 中变化时,
$$
X\longmapsto \Top(X)
$$
为「伪函子」
\begin{equation*}
(𝒰\text{-}\esp)\longrightarrow
(𝒱\text{-}𝒰\text{-}\top){\quoi},\tag{4.1.1.1}
\end{equation*}\etiquette[4.1.1.1]{IV.eq4.1.1.1}
或者用 2-范畴的术语说, 有一个 2-范畴的\emph{非严格}函子 [9]. 实际上, 我们通常允许自己滥用语言, 将 $\Top(gf)$ 与 $\Top(g)\Top(f)$ 等同, 即如同 \eqhyperref[IV.eq4.1.1.1]{4.1.1.1} 是普通范畴之间的真正函子那样推理. 以下处理的其他例子中也将允许类似的语言滥用.\pageoriginale

\subsubsection{}\etiquette[4.1.2]{IV.subsubsec4.1.2}
上述讨论立即推广到带算子群的拓扑空间所对应的拓扑斯 (\hyperref[IV.subsec2.3]{2.3}). 若 $f=(f^{\es},f^{\gr})$,
$$
f:(X,G)\longrightarrow (Y,H)
$$
为带算子空间的态射 (其中
$$
f^{\es}:X\longrightarrow Y\quav f^{\gr}:G\longrightarrow H
$$
分别为连续映射与群同态, 在显然的意义下相容), 就联系出一个拓扑斯态射
$$
\Top(f)\mbox{ 或 } f:\Top(X,G)\longrightarrow \Top(Y,H){\quoi},
$$
其定义留给读者. 当群 $G$ 与 $H$ 平凡时, 就得到 \hyperref[IV.subsubsec4.1.1]{4.1.1} 的定义; 另一方面, 当空间 $X$ 与 $Y$ 各自只有一点时, 逆像函子 $f^{\ast}$ 是「限制算子群」函子
$$
f^{\ast}:B_H\longrightarrow B_G{\quoi},
$$
它将每个 $H$-集配以它通过 $f:G\rightarrow H$ 所定义的 $G$-集. 我们将在 \hyperref[IV.subsec4.5]{4.5} 及 \hyperref[IV.subsubsec4.6.1]{4.6.1} 中以其他形式再次遇到此例.

\subsubsection{}\etiquette[4.1.3]{IV.subsubsec4.1.3}
给定拓扑空间的连续映射 $f:X\rightarrow Y$, 还可联系出相应「大拓扑斯」(\hyperref[IV.subsec2.5]{2.5}) 之间的态射
$$
\TOP(f)\mbox{ 或 } f:\TOP(X)\longrightarrow \TOP(Y){\quoi},
$$
最方便的是\pageoriginale 用逆像函子
$$
f^{\ast}:\TOP(Y)\longrightarrow \TOP(X){\quoi},
$$
定义它, 这就是\emph{限制函子}. 此态射 $\TOP(f)$ 是诱导拓扑斯的所谓「包含」态射的特例, 后者将在第 5 节中研究. 因此, 在与 \hyperref[IV.subsubsec4.1.1]{4.1.1} 同样的限制下, 当 $X$ 在 $(\Esp)$ 中变化时, 可以把拓扑斯 $\TOP(X)$ 看作 $X$ 的函子.

\subsection{$X\mapsto \Top(X)$ 的忠实性性质}\etiquette[4.2]{IV.subsec4.2}

我们要确切说明, 能在多大程度上从拓扑斯 $\Top(X)$ 重构拓扑空间 $X$. 为此, 对两个空间 $X$ 与 $Y$, 宜描述从 $\Top(X)$ 到 $\Top(Y)$ 的态射范畴 (\hyperref[IV.subsec3.2]{3.2}), 以确定「函子」$X\mapsto \Top(X)$ 的忠实性性质. 我们只陈述所得结果, 详情请参看 [9]. 希望亲自验证的读者可参看练习 \hyperref[IV.exer7.8]{7.8}.

\subsubsection{}\etiquette[4.2.1]{IV.subsubsec4.2.1}
回想, 如果拓扑空间 $X$ 的每个不可约闭子集恰有一个一般点, 就称 $X$ 为\emph{素空间} (sober). 实际上使用的空间几乎都是素空间; 特别地, Hausdorff 空间是素空间, 更一般地, 所有点都闭的空间, 以及概形的底空间, 都是素空间. 对拓扑空间 $X$, 可以联系出 (见前引文献或 \EGA $0_I$, 新版) 一个素拓扑空间 $X_{\sob}$ 及连续映射\pageoriginale
\begin{equation*}
\varphi: X\longrightarrow X_{\sob}\tag{4.2.1.1}
\end{equation*}\etiquette[4.2.1.1]{IV.eq4.2.1.1}
它对从 $X$ 到素空间的连续映射是\emph{泛的}; 换言之, 构造了从素空间范畴到「任意」拓扑空间范畴的包含函子 $(\Espsob)\rightarrow (\Esp)$ 的左伴随 $X\mapsto X_{\sob}$ (引号提醒这里涉及一个宇宙!). 显式构造如下: 取 $X$ 的不可约闭子集为 $X_{\sob}$ 的点, 取形如 $U'$ 的集合为它的开子集, 其中 $U$ 为 $X$ 的开子集, $U'\subset X_{\sob}$ 表示与 $U$ 相交的 $X$ 的不可约闭子集所成的集合. 将每个 $x\in X$ 配以 $\{x\}$ 的闭包, 即得到映射 \eqhyperref[IV.eq4.2.1.1]{4.2.1.1}. 空间 $X$ 为素空间当且仅当上述映射为双射, 从而为同胚.

可以验证, $\varphi$ 诱导的函子
$$
\varphi^{-1}:\Ouv(X_{\sob})\longrightarrow \Ouv(X)
$$
为同构, 这蕴涵 $\varphi$ 所定义的拓扑斯态射
$$
\Top(\varphi):\Top(X)\longrightarrow \Top(X_{\sob})
$$
也为同构. 这先验地解释了, 从 $\Top(X)$ 重构 $X$ 的问题中为什么必然要出现 $X_{\sob}$: 因为前者在同构意义下只依赖于 $X_{\sob}$, 故只有当 $X$ 为素空间时, 问题才可能有肯定答案. 下面 (\hyperref[IV.subsec7.1]{7.1}) 将确切说明如何把 $X_{\sob}$ 的点解释为拓扑斯 $\Top(X)$ 的「点」(或者等价地, 为纤维函子), 从而确实由 $\Top(X)$ 重构\pageoriginale $X_{\sob}$.

\subsubsection{}\etiquette[4.2.2]{IV.subsubsec4.2.2}
在每个拓扑空间上, 宜引入如下序关系 $\leq$:
$$
x\leq y\Longleftrightarrow \bar{\{x\}}\subset \bar{\{y\}}\mbox{ 即
} x\in \bar{\{y\}}
$$
(也表述为 $x$ 是 $y$ 的\emph{特殊化}, 或 $y$ 是 $x$ 的\emph{一般化}). 对形如 $X_{\sob}$ 的空间, 这就是 $X$ 的不可约闭子集之间的包含关系.

据此, 在从空间 $X$ 到另一个空间 $Y$ 的映射集上, 宜引入由 $Y$ 的序关系推出的, 称为「特殊化」的序关系, 即
$$
f\leq g\Longleftrightarrow f(x)\leq g(x)\mbox{ 对每个 } x\in
X{\quoi}.
$$
在这些约定下, 有以下结果:

\subsubsection{}\etiquette[4.2.3]{IV.subsubsec4.2.3}
\emph{设 $X$, $Y$ 为两个拓扑空间, 其中 $Y$ 为素空间, 并设 $f$ 与 $g$ 为从 $X$ 到 $Y$ 的两个连续映射. 则从 $\Top(f)$ 到 $\Top(g)$ 至多有一个态射, 且这样的态射存在当且仅当 $g$ 为 $f$ 的特殊化. 最后, 每个拓扑斯态射 $\Top(X)\rightarrow \Top(Y)$ 都同构于形如 $\Top(f)$ 的态射, 其中 $f:X\rightarrow Y$ 为连续映射 (由第一个断言可知它唯一确定).}

如果\pageoriginale 对有序集 $I$, 以当 $i\geq j$ 时从 $i$ 到 $j$ 恰有一个箭头的规定来定义相应范畴 $\cat(I)$, 则上述结果可概括为存在典范范畴等价
$$
\cat(\Hom_{(\esp)}(X,Y))\xrightarrow{\approx}\Hhomtop(\Top(X),\Top(Y))
$$
(其中右边在 \hyperref[IV.subsec3.2]{3.2} 中定义).

以下为这些结果的形式推论:

\begin{enonce*}{推论 4.2.4}\etiquette[4.2.4]{IV.cor4.2.4}
\begin{enumerate}
\renewcommand{\theenumi}{\alph{enumi}}
\renewcommand{\labelenumi}{\theenumi)}
\item 设 $f:X\rightarrow Y$ 为连续映射. 则 $\Top(f):\Top(X)\rightarrow \Top(Y)$ 为拓扑斯等价当且仅当 $f_{\sob}:X_{\sob}\rightarrow Y_{\sob}$ 为同胚 (从而当 $X$ 与 $Y$ 为素空间时, 当且仅当 $f$ 为同胚).

\item 设 $X$ 与 $Y$ 为拓扑空间. 拓扑斯 $\Top(X)$ 与 $\Top(Y)$ 等价当且仅当 $X_{\sob}$ 与 $Y_{\sob}$ 同胚 (从而若 $X$ 与 $Y$ 为素空间, 当且仅当 $X$ 与 $Y$ 同胚).
\end{enumerate}
\end{enonce*}

\subsection{到终拓扑斯的态射: 拓扑斯的常值对象, 截面函子}\etiquette[4.3]{IV.subsec4.3}

用 $P$ (「point」的首字母) 表示标准终拓扑斯, 即 $P=(\Ens)$ (\hyperref[IV.subsec2.2]{2.2}). 设 $E$ 为任意拓扑斯. 我们将看到, \emph{在唯一同构的意义下, 存在唯一的拓扑斯态射}
\hypertarget{sga4.exp.IV.equation.4.3.1}{}$$
f:E\rightarrow P
$$
或者更确切地说, \emph{范畴 $\Hhomtop(E,P)$ 等价于单点范畴}: 因而此范畴的任意两个对象之间都有\pageoriginale 唯一的箭头, 且为同构. (这在一定程度上说明了「终拓扑斯」之名的根据.) 为此, 回想\hyperlink{sga4.exp.IV.subsubsection.3.2.1}{(3.2.1)}, $\Hhomtop(E,P)$ 等价于 $\Hhom(P,E)^{\circ}$ 中由与 $\varinjlim$ 交换且左正合的函子
$$
f^{\ast}:P=(\Ens)\longrightarrow E
$$
所成的满子范畴. 设 $e$ 为单元素集. 每个集合 $X$ 都能典范地写成「$X$ 份 $e$ 的余积」, 从而与 $\varinjlim$ 交换的函子 $g:(\Ens)\rightarrow E$ 的范畴, 通过将每个 $g$ 配以 $E$ 的对象 $g(e)$ 而等价于范畴 $E$. 从 $T=g(e)$ 出发, 用 $g(I)=T\times I$ 可在唯一同构的意义下重构函子 $g$, 其中置 $T\times I=\amalg_{i\in I}T_i$, 对每个 $i\in I$ 有 $T_i=T$. 由 $T$ 如此定义的关于 $I$ 的函子确实与 $\varinjlim$ 交换, 这是因为它显然是从 $E$ 到 $(\Ens)$ 的函子 $X\mapsto \Hom(T,X)$ 的\sisi{右}{左}伴随. 据此, $g$ 左正合的一个显然必要条件是 $g(e)=T$ 为 $E$ 的终对象 (因为 $e$ 为 $(\Ens)$ 的终对象), 而由 $E$ 中「余积是泛的」(\hyperref[IV.subsubsec1.1.2]{1.1.2} b)) 这一事实, 容易推出该条件也充分. 因而逆像函子 $f^{\ast}$ 的范畴等价于 $E$ 的终对象范畴, 后者显然等价于终范畴.

由以上所述可知, 选取态射 \hyperlink{sga4.exp.IV.equation.4.3.1}{(4.3.1)} 实质上等价于选取 $E$ 的一个终对象, 记为 $e_E$. 通过此对象, 有函子的典范同构
\begin{equation*}
 f^{\ast}(I)\simeq e_E\times
 I=\mbox{ } I \mbox{ 份 } e_E \mbox{ 的余积 } \quad \mbox{ 对 } I\in \ob
 (\Ens),\tag{4.3.2}
\end{equation*}\etiquette[4.3.2]{IV.eq4.3.2}
以及\pageoriginale
\begin{equation*}
f_{\ast}(X)=\Hom(e_E,X) \quad \quad \mbox{ 对 } X\in \Ob E.\tag{4.3.3}
\end{equation*}\etiquette[4.3.3]{IV.eq4.3.3}


\setcounter{subsubsection}{3}
\subsubsection{}\etiquette[4.3.4]{IV.subsubsec4.3.4}
上述两个函子将在下文中起重要作用. 对集合 $I$, \eqhyperref[IV.eq4.3.2]{4.3.2} 中 $E$ 的对象 $f^{\ast}(I)=e_E\times I$ 称为 $E$ 中\emph{取值 $I$ 的常值对象} (或者当 $E$ 通过景 $C$ 表为范畴 $C^{\widetilde{\phantom{,}}}$ 时, 称为 \emph{$C$ 上取值 $I$ 的常值层}). 它还常记为 $I_E$, 或当 $E$ 由景 $C$ 定义时记为 $I_C$. $I\mapsto I_E$ 是一个拓扑斯态射的逆像函子, 这个事实说明了它的正合性质; 特别地, 这些性质蕴涵该函子保持所有通常的代数结构类型, 将群变为 $E$ 的群对象, 等等 (\hyperref[IV.subsubsec3.1.2]{3.1.2}). 例如, 给定 $E$ 的群 $G$, 若它同构于形如 $G_{\circ E}$ 的群, 其中 $G_{\circ}$ 为通常的群, 则称它为\emph{常值群} (或在适当情形称为\emph{常值群层}). 对 \hyperref[IV.subsubsec3.1.2]{3.1.2} 中 (或多或少) 规定的意义下的其他每种「代数」结构, 都使用同样的术语.

% END INCLUDED SGA4_English_current_buildable_source_20260906/source/04/drafts/SGA4_IV_LINES_1320_1659_English.texfrag


% BEGIN INCLUDED SGA4_English_current_buildable_source_20260906/source/04/drafts/SGA4_IV_LINES_1660_1806_English.texfrag
\subsubsection{}\etiquette[4.3.5]{IV.subsubsec4.3.5}
注意函子 $I\mapsto I_E$ 不一定全忠实 (甚至不一定忠实: 取 $E$ 为「空拓扑斯」(\hyperref[IV.subsec2.2]{2.2})); 因而 $E$ 的常值对象一般不能在唯一同构意义下确定它所来自的集合 $I$. 如果函子 $I\mapsto I_E$ 全忠实, 就说 $E$ 是 $0$-\emph{无环的}, 或\emph{连通且非空的}. 由伴随函子的一般性质, 这相当于说伴随态射
$$
I\longrightarrow f_{\ast}(f^{\ast}(I))=f_{\ast}(I_E)=\Hom(e_E, I_E)
$$
为\pageoriginale 同构 (从而函子 \eqhyperref[IV.eq4.3.3]{4.3.3} 恢复了 $E$ 的常值对象的「值」). 容易验证, 此条件成立当且仅当 $e_E$ 不是 $E$ 的始对象 $\phi_E$, 即 $E$ 不是「空拓扑斯」(这表述了函子 $I\mapsto I_E$ 的\emph{忠实性}\footnote[1]{或者等价地, 此函子为保守函子 (\hyperrefext[I][I.cor6.3]{6.3}).}), \emph{并且} $e_E$ 为 $E$ 的\emph{连通}对象, 即它不是 $E$ 的两个非「空」对象 (即不是 $E$ 的始对象的对象) 的和.

\subsubsection{}\etiquette[4.3.6]{IV.subsubsec4.3.6}
函子 \eqhyperref[IV.eq4.3.3]{4.3.3} 还常称为\emph{截面函子}, 记为 $\Gamma_E$, $\Gamma(E,-)$, 或简记 $\Gamma$:
\hypertarget{sga4.exp.IV.equation.4.3.6.1}{}\begin{equation*}
\Hom(e_E,X)=\Gamma_E(X)=\Gamma(E,X)=\Gamma(X){\quoi}.\tag{4.3.6.1}
\end{equation*}
它与任意投射极限交换, 但一般不右正合; 下一篇将研究它 (在交换群对象上) 的导出函子.

\subsection{从「空拓扑斯」出发的态射}\etiquette[4.4]{IV.subsec4.4}

设 $\phi_{\text{top}}$ 为空拓扑斯, 对应于等价于终范畴的层范畴 $\Phi$ (\hyperref[IV.subsec2.2]{2.2}). 设 $E$ 为拓扑斯. 从 $E$ 到 $\Phi$ 的函子范畴显然等价于单点范畴, 而每个这样的函子都与归纳极限和投射极限交换 (这并没有什么功劳), 因而可看作拓扑斯态射 $\phi_{\text{top}}\rightarrow E$ 的逆像函子. 从而\emph{范畴 $\Hhomtop(\phi_{\text{top}},E)$ 等价于单点范畴}, 特别地, \emph{在唯一\pageoriginale 同构的意义下, 恰有一个拓扑斯态射}
\begin{equation*}
\phi_{\text{top}}\longrightarrow E{\quoi}.\tag{4.4.1}
\end{equation*}\etiquette[4.4.1]{IV.eq4.4.1}
这在一定程度上说明了 \hyperref[IV.subsec2.2]{2.2} 所引入的「始拓扑斯」这一术语的根据.

也可确定拓扑斯态射
\begin{equation*}
E\longrightarrow \phi_{\text{top}}{\quoi}\ ;\tag{4.4.2}
\end{equation*}\etiquette[4.4.2]{IV.eq4.4.2}
立即可验证, 这样的态射存在当且仅当 $E$ 的始对象也是终对象, 即当且仅当 $E$ 本身为「空拓扑斯」; 此时范畴 $\Hhomtop(E,\phi_{\text{top}})$ 仍等价于单点范畴. 于是唯一的 (在同构意义下) 态射 \eqhyperref[IV.eq4.4.2]{4.4.2} 为拓扑斯等价.

\subsection{随群 $G$ 变化的分类拓扑斯 $B_G$}\etiquette[4.5]{IV.subsec4.5}

\subsubsection{}\etiquette[4.5.1]{IV.subsubsec4.5.1}
设 $E$ 为拓扑斯, 并设
$$
f:G\longrightarrow H
$$
为 $E$ 中群对象的态射. 它诱导一个「限制算子群」函子
$$
f^{\ast}:B_H\longrightarrow B_G{\quoi},
$$
其中使用 \hyperref[IV.subsec2.4]{2.4} 的记号. 此函子显然与归纳极限及投射极限交换; 更可以将它解释为拓扑斯态射
$$
B_f\mbox{ 或 } f:B_G\longrightarrow B_H{\quoi}.
$$
所对应的逆像函子. 相应\pageoriginale 的正像函子
$$
f_{\ast}:B_G\longrightarrow B_H
$$
容易通过公式
$$
f_{\ast}(X)=\Hhom_G(H_s,X){\quoi},
$$
明确给出, 其中 $X$ 为具有 $G$ 左作用的 \sisi{$X$}{$E$} 的对象, $H_s$ 为 $H$ 配上由 \sisi{$F$}{$f$} 诱导的 $G$ 左作用, 而 $\Hhom_G$ 表示下面 (\hyperref[IV.subsec10.2]{10.2}) 定义的对象 $\Hhom$ 的显然子对象; $H$ 在 $\Hhom_G(H_s,X)$ 上的左作用, 由 $H$ 通过\emph{右}平移在 $H_s$ 上的右作用诱导.

因为逆像函子 $f^{\ast}$ 与任意 \sisi{$\varinjlim$}{$\varprojlim$} (而不只是有限 \sisi{$\varinjlim$}{$\varprojlim$}) 交换, 所以它本身为函子
$$
f_{!}:B_G\longrightarrow B_H{\quoi},
$$
的右伴随, 从而如 \hyperref[IV.subsubsec3.1.3]{3.1.3} 有三个伴随函子的序列:
$$
f_{!},f^{\ast}, f_{\ast}{\quoi}.
$$
函子 $f_{!}$ 容易通过公式
$$
f_{!}(X)=H\times {}^GX{\quoi},
$$
明确给出, 其中右边表示由 $G$ 在 $X$ 上的作用 (左作用) 与在 $H$ 上的作用 (通过右平移和 $f$ 的右作用) 所诱导的「缩积」, 定义为 $H\times X$ 关于 $G$ 的作用
$$
g\circ (h,x)=(hg^{-1}, gx){\quoi}.
$$
的商.

函子 $f_{!}$ 为左伴随, 显然与归纳极限交换, 但一般不左正合 (即它\pageoriginale 不能再看作拓扑斯态射 $B_H\rightarrow B_G$ 的逆像函子). 事实上, 容易验证, 它只有在 $f:G\rightarrow H$ 为同构时才可能左正合. 同样, 与投射极限交换的函子 $f_{\ast}$ 一般不右正合, 因而更不会一般地有右伴随. 至少当 $E$ 为点拓扑斯, 即 $G$ 与 $H$ 为通常的群时, $f_{\ast}$ 只有在 $f$ 为同构时才右正合.

再给定一个群态射 $g:H\rightarrow K$, 就如 \hyperref[IV.subsubsec4.1.1]{4.1.1} 得到在典范同构意义下的传递性, 因此在通常的保留说明下, 可以把分类拓扑斯 $B_G$ 看作\emph{函子性地}依赖于群 $G$. 把这种函子性推广到 $G$ \emph{与}拓扑斯 $E$ 同时变化的情形, 留给读者.


% END INCLUDED SGA4_English_current_buildable_source_20260906/source/04/drafts/SGA4_IV_LINES_1660_1806_English.texfrag


% BEGIN INCLUDED SGA4_English_current_buildable_source_20260906/source/04/drafts/SGA4_IV_LINES_1807_1981_English.texfrag
\subsubsection{随 pro-群 $𝒢$ 变化的拓扑斯 $B_{𝒢}$}\etiquette[4.5.2]{IV.subsubsec4.5.2}

请读者循 \hyperref[IV.subsubsec4.5.1]{4.5.1} 的叙述, 明确拓扑斯 $B_{𝒢}$ (\hyperref[IV.subsec2.7]{2.7}) 关于 $𝒢$ 的协变性. 不过应注意, 对并非本质常值的 pro-群之间的态射 $f:𝒢\rightarrow ℋ$, 相应的拓扑斯态射 $B_{𝒢}\rightarrow B_{ℋ}$ 一般不能定义函子 $f_{!}$ (其右伴随应为限制算子 pro-群的函子 $f^{\ast}$). 为简化陈述, 假设 $𝒢$ 与 $ℋ$ 由 \emph{pro-有限}群 $G$ 与 $H$ 定义. 此时容易看出, $f_{!}$ 存在当且仅当所考虑态射 $f:G\rightarrow H$ 的像在 $H$ 中具有有限指数, 且在\pageoriginale 此情形, $f_{!}$ 由与 \hyperref[IV.subsec4.5]{4.5} 中相同的公式给出.

\subsection{随范畴 $C$ 变化的拓扑斯 $\widehat{C}$}\etiquette[4.6]{IV.subsec4.6}

\subsubsection{}\etiquette[4.6.1]{IV.subsubsec4.6.1}
设
$$
f:C\longrightarrow C'
$$
为从范畴 $C\in 𝒰$ 到另一个范畴 $C'$ 的函子, 因而有
$$
f^{\ast}:\widehat{C}'\longrightarrow \widehat{C}\quav
f^{\ast}(F)=F\circ f{\quoi}.
$$
此函子显然与投射极限及归纳极限交换; 因而更可将它看作拓扑斯态射
$$
\widehat{f}\mbox{ 或 } f:\widehat{C}\rightarrow \widehat{C}'{\quoi}.
$$
的逆像函子. 相应的正像函子
$$
f_{\ast}:\widehat{C}\longrightarrow \widehat{C}'
$$
正是 I \hyperref[I.prop5.1]{5.1} 中同样记为 $f_{\ast}$ 的函子. 此外 (从 $f^{\ast}$ 也与任意 $\varprojlim$ 交换这一事实可以预料), $f^{\ast}$ 还有左伴随
$$
f_{!}:\widehat{C}\longrightarrow \widehat{C}'{\quoi},
$$
(它在 \hyperrefext[I][I.prop5.1]{5.1} 中也记为 $f_{!}$). 从而得到三个伴随函子的序列
$$
f_{!}, f^{\ast}, f_{\ast}{\quoi},
$$
其中第一个函子并且是 $f:C\rightarrow C'$ 到预层范畴的扩充 (对 $C$ 与 $C'$ 到范畴 $\widehat{C}$ 与 $\widehat{C}'$ 的通常嵌入而言). 应注意 $f_{!}$ 一般不左正合, $f_{\ast}$ 也一般不右正合; 这消除了与变化的范畴 $C$ 相联系的拓扑斯 $\widehat{C}$ 的变性方向\pageoriginale 的歧义: 此拓扑斯为 $C$ 的\emph{协变}函子, 但要作由传递同构引起的通常保留说明 (参见 \hyperref[IV.subsubsec4.1.1]{4.1.1}). 当 $C$ 与 $C'$ 的对象集都为单元素集, 从而 $C$ 与 $C'$ 等同于幺半群 $G$ 与 $G'$ 时, 相应的拓扑斯为分类拓扑斯 $B_{G^{\circ}}$ 与 $B_{{G'}^{\circ}}$, 就重新得到它们关于变化的 $G$ 的变性; 我们已从其他观点在 \hyperref[IV.subsubsec4.1.2]{4.1.2} 及 \hyperref[IV.subsec4.5]{4.5} 中见过它 (其中限于群而非幺半群并非本质的).

\subsubsection{}\etiquette[4.6.2]{IV.subsubsec4.6.2}
对两个范畴 $C,C'\in 𝒰$, 引入典范函子
\begin{equation*}
\Hhom(C,C')\longrightarrow
\Hhomtop(\widehat{C},\widehat{C}'){\quoi}.\tag{4.6.2.1}
\end{equation*}\etiquette[4.6.2.1]{IV.eq4.6.2.1}
就能进一步明确拓扑斯 $\widehat{C}$ 对 $C$ 的 2-函子性依赖. 还须在箭头上定义此函子. 为此注意, $f\mapsto f^{\ast}$ 将右边 (在范畴等价意义下) 等同于 $\Hhom(\widehat{C}',\widehat{C})^{\circ}$ 的一个满子范畴\hyperlink{sga4.exp.IV.subsubsection.3.2.1}{(3.2.1)}. 若 $f,g:C\rightarrow C'$ 为两个函子, 则函子的每个态射 $u:f\rightarrow g$ 对 $F\in\ob\widehat{C}'$ 都定义函子的态射 $F\circ g\rightarrow F\circ f$ (因为 $F$ 为\emph{反}变函子), 因而定义态射 $g^{\ast}\rightarrow f^{\ast}$, 从而定义态射 $\widehat{f}\rightarrow\widehat{g}$, 如所断言的那样.

当 $C$ 为单点范畴时, $\widehat{C}$ 为点拓扑斯 (\hyperref[IV.subsec2.2]{2.2}), 记为 $P$, 而 \eqhyperref[IV.eq4.6.2.1]{4.6.2.1} 解释为从 $C'$ 到 $\widehat{C}'$ 的「\emph{点范畴}」的自然函子
\begin{equation*}
C'\longrightarrow \Hhomtop(P,\widehat{C}') \mathop{=}^{\dfn}
\Ppoints(\widehat{C}')\tag{4.6.2.2}
\end{equation*}\etiquette[4.6.2.2]{IV.eq4.6.2.2}
这个点范畴将在第 6 节中研究. 此\pageoriginale 函子不一定为范畴等价\hyperlink{sga4.exp.IV.subsubsection.7.3.3}{(7.3.3)}; 因而 \eqhyperref[IV.eq4.6.2.1]{4.6.2.1} 更不一定为范畴等价.

另一方面, 函子 \eqhyperref[IV.eq4.6.2.1]{4.6.2.1} 总是\emph{全忠实}的. 为看到这一点, 注意若 $f,g:E\rightarrow E'$ 为两个拓扑斯态射, 且 $f_{!}$ 与 $g_{!}$ 已定义 (\hyperref[IV.subsubsec3.1.3]{3.1.3}), 则伴随函子理论给出典范同构
$$
\Hom(f_{!},g_{!})\simeq \Hom(g^{\ast}, f^{\ast})\simeq
\Hom(f_{\ast}, g_{\ast}) \mathop{=}^{\dfn} \Hom(f,g){\quoi}.
$$
将此应用于与 $f,g:C\rightarrow C'$ 相联系的形如 $\widehat{f}$ 与 $\widehat{g}$ 的函子, 并注意自然扩充映射 $\Hom(f,g)\rightarrow \Hom(f_{!},g_{!})$ 为双射 (I \hyperref[I.prop7.8]{7.8}), 就得到所断言的结果.

\subsubsection{}\etiquette[4.6.3]{IV.subsubsec4.6.3}
可以问, 函子 \eqhyperref[IV.eq4.6.2.1]{4.6.2.1} 何时为范畴等价, 即何时本质满射. 这是确定一切拓扑斯态射 $\widehat{C}\rightarrow\widehat{C}'$ 这一问题的特例. 更一般地, 若 $C$ 为 $𝒰$ 中的范畴, $E$ 为拓扑斯, 可以设法确定拓扑斯态射
$$
f:\widehat{C}\longrightarrow E{\quoi}.
$$
这些态射的范畴等价于与任意归纳极限和有限投射极限交换的函子 $f^{\ast}:E\rightarrow\widehat{C}$ 的范畴的反范畴. 函子 $E\rightarrow\widehat{C}$ 可解释为函子 $E\times C^{\circ}\rightarrow(𝒰\text{-}\Ens)=(\Ens)$, 或等价地, 解释为函子 $F:C^{\circ}\rightarrow\Hhom(E,(\Ens))$. 所讨论的正合性质此时表为: 对 $C$ 的每个对象 $X$, 要求函子 $F(X):E\rightarrow(\Ens)$ 与任意归纳极限和有限投射极限交换 (或如第 6 节将说的, $F(X)$ 是 $E$ 上的「\emph{纤维\pageoriginale 函子}」). 这相当于说 $F(X)$ 是拓扑斯态射 $P\rightarrow E$ 的逆像函子, 因而最终得到典范范畴等价
\begin{equation*}
\Hhomtop(\widehat{C}, E) \xrightarrow{\approx} \Hhom(C,
\Ppoints(E)){\quoi},\tag{4.6.3.1}
\end{equation*}\etiquette[4.6.3.1]{IV.eq4.6.3.1}
其中为简便起见, 用
$$
\Ppoints(E)=\Hhomtop(P,E)
$$
表示拓扑斯 $E$ 的点范畴.

当 $E$ 形如 $\widehat{C}'$ 时, 由定义立即可知, \eqhyperref[IV.eq4.6.2.1]{4.6.2.1} 与上述等价 \eqhyperref[IV.eq4.6.3.1]{4.6.3.1} 的复合, 是由 $F\mapsto i\circ F$ 定义的函子
\begin{equation*}
\Hhom(C,C')\longrightarrow
\Hhom(C,\Ppoints(\widehat{C}'))\tag{4.6.3.2}
\end{equation*}\etiquette[4.6.3.2]{IV.eq4.6.3.2}
其中 $i:C'\rightarrow\Ppoints(\widehat{C}')$ 为典范嵌入 \eqhyperref[IV.eq4.6.2.2]{4.6.2.2}. 立即可得, \emph{\eqhyperref[IV.eq4.6.2.1]{4.6.2.1} 为等价的充分必要条件是 $C$ 为空, 或者 \eqhyperref[IV.eq4.6.2.2]{4.6.2.2} 本质满射} (从而为范畴等价). 关于 $C'$ 的后一个条件在某些有意义的情形中满足, 尤其当 $C'$ 为群 $G$ 所定义的单对象范畴时 (\hyperref[IV.subsubsec7.2.5]{7.2.5}).



% END INCLUDED SGA4_English_current_buildable_source_20260906/source/04/drafts/SGA4_IV_LINES_1807_1981_English.texfrag


% BEGIN INCLUDED SGA4_English_current_buildable_source_20260906/source/04/drafts/SGA4_IV_LINES_1982_2188_English.texfrag
\begin{enonce*}[remark]{注 4.6.4}\etiquette[4.6.4]{IV.rem4.6.4}\nde{关于拓扑斯的同伦型概念的较新观点, 见 [15] $\S$ \hyperref[IV.subsubsec7.1.6]{7.1.6}.}
将任意范畴 $C$ 配以拓扑斯 $\widehat{C}$, 提示范畴 $C$ 像拓扑斯一样具有拓扑性质的不变量 (上同调群、同伦群等等). $\widehat{C}$ 以交换群对象 $F$ 为系数的上同调群 (按第 V 篇中研究的一般意义), 恰为拓扑学家已熟悉的函子 $\varprojlim$ 的导出函子 $\varprojlim^{(n)}$ 的值.\pageoriginale J. L. Verdier 与 D. G. Quillen (相互独立地) 验证, 当限于常值, 或更一般地局部常值系数时, 这些上同调群等同于与 $C$ 典范相联系的半单纯集 $\Nerf(C)$ 的上同调群 [5, ${\rm n^{\circ}}$ 212, 命题 4.1], 而且在 [2] 的「同伦范畴」中同构的意义下, 每个半单纯集都可由某个范畴 $C$ 得到. 使用一个适当的拓扑斯同伦型概念 (这里不加详述), 可以说, 拓扑学家的半单纯同伦型恰是特殊形状 $\widehat{C}$ 的拓扑斯的同伦型, 更一般地, 是每个对象都容许一个加细其他一切覆盖的覆盖的拓扑斯 $E$ 的同伦型 (\hyperref[IV.subsec2.6]{2.6})\footnote[1]{还可更一般地, 是在一个留给读者明确的显然意义下「局部 $\infty$-连通」的拓扑斯的同伦型.}. 若 $E$ 不满足这个条件, 则它的同伦型至多可用一个适当的半单纯集的\emph{投射系}来表出 [1].
\end{enonce*}

\subsection{随景 $C$ 变化的拓扑斯 $C^{\sim}$ (余连续函子)}\etiquette[4.7]{IV.subsec4.7}

设
$$
f:C\longrightarrow C'
$$
为 $𝒰$ 中的景之间的余连续函子 (\hyperrefext[III][III.def2.1]{2.1}), 即从 $\widehat{C}$ 到 $\widehat{C}'$ 的函子 $\widehat{f}_{\ast}$ (\hyperref[IV.subsubsec4.6.1]{4.6.1}) 将 $C^{\sim}$ 映入 $C'^{\sim}$ 的函子, 换言之, 它诱导函子
\begin{equation*}
\widetilde{f}_{\ast}:\widetilde{C}\longrightarrow
{C'}^{\widetilde{\phantom{,}}}\tag{4.7.1}
\end{equation*}\etiquette[4.7.1]{IV.eq4.7.1}
使下面的函子图表交换:
\[\vcenter{
\xymatrix@C=2.5cm@R=1.25cm{ C^{\widetilde{\phantom{a}}}
\ar[r]^{\widetilde{f}_{\ast}} \ar[d]_{\sisi{i}{i_{\ast}}}&
C'^{\widetilde{\phantom{a}}} \ar[d]^{\sisi{i_{\ast}}{i_{\ast}'}}\\
\widehat{C}  \ar[r]^{\widehat{f}_{\ast}}       & \widehat{C}'\quad,
}}\leqno{(4.7.2)}
\]\etiquette[4.7.2]{IV.eq4.7.2}
其中\pageoriginale \sisi{$i$}{$i_{\ast}$} 与 \sisi{$i'$}{$i_{\ast}'$} 为包含函子. 我们已经看到 (\hyperrefext[III][III.prop2.3]{2.3}), 此时函子 $\widetilde{f}_{\ast}$ 有左伴随 $\widetilde{f}^{\ast}$, 且后者左正合. 换言之, $\widetilde{f}_{\ast}$ 为拓扑斯态射
$$
\widetilde{f}\mbox{ 或 } f:C^{\sim}\longrightarrow
{C'}^{\widetilde{\phantom{,}}}{\quoi},
$$
所对应的正像函子; 相应的逆像函子当然为 $\widetilde{f}^{\ast}$. 对所涉及的函子取左伴随, 交换图表 \eqhyperref[IV.eq4.7.2]{4.7.2} 还给出一个在同构意义下交换的图表:
\[\vcenter{
\xymatrix@C=2.5cm@R=1.5cm{
C^{\widetilde{\phantom{a}}} &  C'^{\widetilde{\phantom{a}}} \ar[l]_{\widetilde{f}^{\ast}}\\
\widehat{C}  \ar[u]^{\underline{a}}  & \widehat{C}'
\ar[l]_{\widehat{f}^{\ast}}\ar[u]_{\underline{a}'}\quad,
}}\leqno{(4.7.3)}
\]\etiquette[4.7.3]{IV.eq4.7.3}
其中 $\underline{a}$ 与 $\underline{a}'$ 为「相伴层」函子. 此图表立即重新给出公式 (\hyperrefext[III][III.prop2.3]{2.3})
$$
\widetilde{f}^{\ast}=\underline{a}\widehat{f}^{\ast} i'{\quoi}.
$$

拓扑斯态射 $\widehat{f}:\widehat{C}\rightarrow \widehat{C}'$ 的传递性质, 蕴涵与余连续函子相联系的拓扑斯态射 $\widetilde{f}:C^{\widetilde{\phantom{,}}}\rightarrow{C'}^{\widetilde{\phantom{,}}}$ 具有同样性质. 因此, 若取余连续函子为景的「态射」, 就可以说\emph{拓扑斯 $C^{\widetilde{\phantom{,}}}$ 以\emph{协变}方式随 $C$ 函子性地变化}.

在\pageoriginale 至此遇到的一切情形中, 所用的余连续函子 $f$ 也连续, 即 (\hyperrefext[III][III.def1.1]{1.1}) 它「扩充」为与归纳极限交换的函子
$$
\widetilde{f}_{!}:C^{\widetilde{\phantom{,}}}\longrightarrow
{C'}^{\widetilde{\phantom{,}}}
$$
它是 $\widetilde{f}^{\ast}$ 的左伴随, 因而有三个伴随函子的序列
$$
\widetilde{f}_{!}, \widetilde{f}^{\ast},
\widetilde{f}_{\ast}{\quoi}.
$$
应注意函子 $\widetilde{f}_{!}$ 一般不左正合, $\widetilde{f}_{\ast}$ 也一般不右正合. 这消除了当 $C$ 通过仅假定余连续、或甚至连续且余连续的函子变化时, 拓扑斯 $C^{\widetilde{\phantom{,}}}$ 的变性方向的歧义.

\begin{enonce*}[remark]{注 4.7.4}\etiquette[4.7.4]{IV.rem4.7.4}
给定拓扑斯态射
$$
F:E\longrightarrow E'{\quoi},
$$
存在 $F^{\ast}$ 的左伴随函子 $F_{!}$ 的充分必要条件是, 能够把 $E$ 与 $E'$ (在等价意义下)「实现」为两个景 $C,C'\in 𝒰$ 的 $C^{\widetilde{\phantom{,}}}$ 与 ${C'}^{\widetilde{\phantom{,}}}$, 并找到\emph{连续且余连续}函子 $f:C\rightarrow C'$, 使得 $F$ 等同于 $\widetilde{f}$. 这显然充分; 对必要性, 只须分别取 $E$ 与 $E'$ 的小满生成子范畴 $C$ 与 $C'$, 使得
$$
F_{!}(\ob C)\subset \ob C'{\quoi},
$$
给它们赋予由 $E$ 与 $E'$ 的拓扑诱导的拓扑, 再取 $F_{!}$ 诱导的函子为 $f$ 即可 (参见 \hyperref[IV.cor1.2.1]{1.2.1}). 回想 (\hyperref[IV.rem1.8]{1.8}), $F_{!}$ 存在也表示 $F^{\ast}$ 与投射极限交换, 或者\pageoriginale 在这里因该函子左正合而等价地说, 它与积交换.
\end{enonce*}

\setcounter{subsubsection}{4}
\subsubsection{}\etiquette[4.7.5]{IV.subsubsec4.7.5}
若问哪些拓扑斯态射 $F:E\rightarrow E'$ 可由景的余连续 (不一定连续) 函子实现, 类似地得到以下回答: $E$ 中使从 $E'$ 到 $(\Ens)$ 的函子 $X'\mapsto \Hom(X,F^{\ast}(X'))$ 与投射极限 (或等价地与积) 交换的那些对象 $X$, 所成的集必须为 \emph{$E$ 的生成族}.\nde{这个陈述看来不正确: 所表述的条件等价于断言 $F^{\ast}$ 与任意积交换, 或等价地, $F^{\ast}$ 有左伴随 $F_{!}$. 对每个景 $C$, 当右边的范畴赋予平凡拓扑时, 恒等函子 $F=id:C\longrightarrow C$ 是余连续的. 此时函子 $F^{\ast}:\hat{C}\longrightarrow\widetilde{C}$ 为相伴层函子, 而它一般不与任意积交换.}

\subsection{景 $C$ 的拓扑斯态射 $\widetilde{C}\rightarrow\widehat{C}$}\etiquette[4.8]{IV.subsec4.8}

设 $C$ 为小景, 因而与之联系着两个拓扑斯 $C^{\widetilde{\phantom{,}}}$ 与 $\widehat{C}$ (后者不依赖于 $C$ 上所赋的拓扑). 在 II \hyperref[II.theo3.4]{3.4} 中, 我们定义了「相伴层」函子
$$
\underline{a}:\widehat{C}\longrightarrow
C^{\widetilde{\phantom{,}}}{\quoi},
$$
并证明它左正合且与归纳极限交换. 因而它是拓扑斯态射
$$
p:\widetilde{C}\longrightarrow \widehat{C},\;
p^{\ast}=\underline{a}{\quoi},
$$
所对应的逆像函子, 相应的正像函子为典范包含
$$
i=p_{\ast}:C^{\widetilde{\phantom{,}}}\longrightarrow
\widehat{C}{\quoi}.
$$
众所周知, 后一个函子一般不右正合 (它在交换对象上的导出函子给出 V 2 的上同调预层 $\Xx^n(F)$), 而 $\underline{a}$ 一般不与任意 $\varprojlim$ 交换. 这消除了 $\widehat{C}$ 与 $C^{\widetilde{\phantom{,}}}$ 之间「自然」拓扑斯态射方向的歧义.

给定景的余连续函子
$$
f:C\longrightarrow C'
$$
就得到\pageoriginale 拓扑斯态射图表
\[
\xymatrix@C=2.25cm@R=1.25cm{ C^{\widetilde{\phantom{a}}} \ar[d]_{\sisi{p}{p_{\ast}}}
\ar[r]^{\widetilde{f}} &
C'^{\widetilde{\phantom{a}}} \ar[d]^{\sisi{p'}{p'_{\ast}}}\\
\widehat{C} \ar[r]^{\widehat{f}}&  \widehat{C}' }
\]
它在典范同构意义下交换: 这正是函子图表 \eqhyperref[IV.eq4.7.2]{4.7.2} 的交换性所表达的. 因此可以说, 当 $C$ 通过景之间的\emph{余连续}函子变化时, 拓扑斯态射 $p:C^{\widetilde{\phantom{,}}}\rightarrow \widehat{C}$ 关于 $C$ 是函子性的.


% END INCLUDED SGA4_English_current_buildable_source_20260906/source/04/drafts/SGA4_IV_LINES_1982_2188_English.texfrag


% BEGIN INCLUDED SGA4_English_current_buildable_source_20260906/source/04/drafts/SGA4_IV_LINES_2189_2354_English.texfrag
\subsection{景的连续函子的作用. 景态射}\etiquette[4.9]{IV.subsec4.9}

\subsubsection{}\etiquette[4.9.1]{IV.subsubsec4.9.1}
设
$$
f:C\longrightarrow C'
$$
为从 $C$ 到 $C'$ 的连续函子, 即函子 $\widehat{f}^{\ast}:\widehat{C}'\rightarrow\widehat{C}$ 将 ${C'}^{\widetilde{\phantom{,}}}$ 映入 $C^{\widetilde{\phantom{,}}}$, 因而诱导函子
$$
f_s:{C'}^{\widetilde{\phantom{,}}}\longrightarrow
C^{\widetilde{\phantom{,}}}{\quoi}.
$$
我们已经看到 (\hyperrefext[III][III.prop1.2]{1.2}), 后者有左伴随
$$
f^s:C^{\widetilde{\phantom{,}}}\longrightarrow
{C'}^{\widetilde{\phantom{,}}}{\quoi},
$$
此外它在显然意义下「扩充」$f$. 应注意, 一般 $f^s$ 不\sisi{右}{左}正合 (即使 $f$ 也余连续), $f_s$ 也不与归纳极限交换, 因而若无进一步假设, 给定 $f$ 并不能描述 $C^{\widetilde{\phantom{,}}}$ 与 ${C'}^{\widetilde{\phantom{,}}}$ 之间任一方向的拓扑斯态射. $f$ 余连续的情形, 即 $f_s$ 与归纳极限交换, 从而\pageoriginale 可以看作拓扑斯态射 $\widetilde{f}:C^{\widetilde{\phantom{,}}}\rightarrow{C'}^{\widetilde{\phantom{,}}}$ 的逆像函子的情形, 已在 \hyperref[IV.subsec4.7]{4.7} 中考察. 我们将考察函子 $f^s$ 左正合的情形, 此时可将它看作\emph{相反方向}的拓扑斯态射的逆像函子:
\begin{equation*}
\Top(f)=g:{C'}^{\sim}\longrightarrow C^{\sim}{\quoi}. \tag{4.9.1.1}
\end{equation*}\etiquette[4.9.1.1]{IV.eq4.9.1.1}
注意, \emph{按照 \hyperref[IV.subsubsec3.1.3]{3.1.3} 的一般建议, 我们有意不用字母 $\widetilde{f}$ 或 $f$ 表示此态射, 以免与 \hyperref[IV.subsec4.7]{4.7} 的情形混淆.} 当函子 $f^s$ 左正合时, 有时称连续函子 $f:C\rightarrow C'$ 为\emph{从 $C'$ 到 $C$ 的景态射} (注意: \emph{不是从 $C$ 到 $C'$}); 换言之, 存在拓扑斯态射 \eqhyperref[IV.eq4.9.1.1]{4.9.1.1}, 使相应逆像函子
$$
g^{\ast}:C^{\widetilde{\phantom{,}}}\longrightarrow
{C'}^{\widetilde{\phantom{,}}}
$$
\emph{扩充}函子 $f$, 即使下列函子图表在同构意义下交换:
\[
\xymatrix@C=2cm@R=1.5cm{
C \ar[d]_{\epsilon} \ar[r]^{f}& C' \ar[d]^{{\epsilon}'}\\
C^{\widetilde{\phantom{a}}} \ar[r]^{g^{\ast}} &
C'^{\widetilde{\phantom{a}}} &, }
\]
%% widetilde + \phantom: horrible hack \XXX
其中 $\epsilon$ 与 $\epsilon'$ 为 II \hyperref[II.subsubsec4.4.0]{4.4.0} 的典范函子.

\subsubsection{}\etiquette[4.9.2]{IV.subsubsec4.9.2}
实际上, 从有限投射极限在 $C$ 中可表, 且 $f$ 与它们交换这一事实, 可认出连续函子 $f:C\rightarrow C'$ 是从 $C'$ 到 $C$ 的景态射 (\hyperrefext[III][III.prop1.3]{1.3} 5)). 对\pageoriginale $C$ 作上述假设 (实践中几乎总满足), 并再假设 $C'$ 的拓扑比典范拓扑粗 (也几乎总满足), 则 $f$ 是从 $C'$ 到 $C$ 的景态射的上述充分条件 ($f$ 左正合) 也是必要的.

\subsubsection{}\etiquette[4.9.3]{IV.subsubsec4.9.3}
按上述讨论的思路, 若 $C$ 与 $C'$ 为两个 $𝒰$-景, 则从 $C$ 到 $C'$ 的景态射范畴 $\Mmorsite(C',C)$ 应定义为从 $C$ 到 $C'$ 的函子范畴的\emph{反范畴 $\Hhom(C,C')^{\circ}$} 中, 由那些为景态射 (从 $C'$ 到 $C$) 的函子所成的满子范畴. 这样就得到典范函子 (在唯一同构的意义下定义)
$$
\Mmorsite(C',C)\longrightarrow
\Hhomtop({C'}^{\widetilde{\phantom{,}}},
C^{\widetilde{\phantom{,}}}){\quoi}.
$$

\begin{enonce*}{命题 4.9.4}\etiquette[4.9.4]{IV.prop4.9.4}
设 $E$ 为 $𝒰$-拓扑斯, $C$ 为 $𝒰$-景. 将每个拓扑斯态射 $f:E\rightarrow C^{\widetilde{\phantom{,}}}$ 配以相应逆像函子 $f^{\ast}:C^{\widetilde{\phantom{,}}}\rightarrow E$ 在 $C$ 上的「限制」的函子 $f\mapsto f^{\ast}| C=f^{\ast}\circ\epsilon_C$, 诱导范畴等价
$$
\Hhomtop(E,C^{\widetilde{\phantom{,}}})\xrightarrow{\approx}\Mmorsite(E,C)(\hookrightarrow
\Hhom(C,E)^{\circ}){\quoi}.
$$
当有限 $\varprojlim$ 在 $C$ 中可表时, 相应的全忠实函子
$$
\Hhomtop(E,C^{\widetilde{\phantom{,}}})\longrightarrow
\Hhom(C,E)^{\circ}
$$
的本质像为左正合且连续的函子 $g:C\rightarrow E$, 或等价地, 为左正合且把\pageoriginale 覆盖族映为覆盖族的函子.
\end{enonce*}

最后一个断言由第一个断言及 \hyperref[IV.subsubsec4.9.2]{4.9.2} 得出. 另一方面, 由 \hyperrefext[III][III.prop1.2]{1.2} iv) 及 IV \hyperref[IV.theo1.2]{1.2} iii) 可知, 函子 $G\rightarrow G\circ\epsilon$ 在从 $C^{\widetilde{\phantom{,}}}$ 到 $E$ 的\emph{连续}函子范畴与从 $C$ 到 $E$ 的\emph{连续}函子范畴之间诱导等价, 按该处的记号, 拟逆由 $g\rightarrow g^s$ 给出. 另一方面, 按定义 $g$ 为景态射当且仅当 $g^s$ 为拓扑斯态射 $E\rightarrow C^{\widetilde{\phantom{,}}}$ 所对应的逆像函子, 从而由\hyperlink{sga4.exp.IV.subsubsection.3.2.1}{3.2.1} 得到结论; 「或等价地」由 \hyperrefext[III][III.prop1.6]{1.6} 得出.

应从 \hyperref[IV.prop4.9.4]{4.9.4} 记住的要点是, (当有限 $\varprojlim$ 在 $C$ 中可表时) 给定拓扑斯态射 $f:E\rightarrow C^{\widetilde{\phantom{,}}}$, 与给定左正合且把覆盖族映为覆盖族的函子 $g:C\rightarrow E$, 「是一回事」.

\setcounter{subsubsection}{4}
\subsubsection{}\etiquette[4.9.5]{IV.subsubsec4.9.5}
用 \hyperref[IV.subsubsec4.9.1]{4.9.1} 及 \hyperref[IV.subsubsec4.9.3]{4.9.3} 的讨论, 可如通常那样看到, 通过刚才明确的景态射与景态射之间态射的概念, 拓扑斯 $C^{\widetilde{\phantom{,}}}$ 函子性地 (或更确切地, 2-函子性地) 依赖于变化的 $𝒰$-景 $C$. 注意, 按刚才引入的术语, $C^{\widetilde{\phantom{,}}}$ \emph{协变}地依赖于景 $C$.

\subsubsection{}\etiquette[4.9.6]{IV.subsubsec4.9.6}
显然 (与余连续函子概念的情形不同, 参见 \hyperref[IV.subsec4.7]{4.7} 4), \emph{每个}拓扑斯态射 $F:E\rightarrow E'$ 都可用景态射 $f:C\rightarrow C'$ 实现 (即连续函子 $f:C'\rightarrow C$, 满足 \ldots): 只须分别选取 $E$ 与 $E'$ 的小满生成子范畴 $C$\pageoriginale 与 $C'$, 赋予 $E$ 与 $E'$ 的拓扑诱导的拓扑, 使得
$$
F^{\ast}(\ob C')\subset \ob C{\quoi},
$$
并取由 $F^{\ast}$ 诱导的函子为 $f$. 这说明了为什么实践中遇到的大多数拓扑斯态射确实由景态射描述 (而不是如 \hyperref[IV.subsec4.7]{4.7} 用余连续函子描述).


% END INCLUDED SGA4_English_current_buildable_source_20260906/source/04/drafts/SGA4_IV_LINES_2189_2354_English.texfrag


% BEGIN INCLUDED SGA4_English_current_buildable_source_20260906/source/04/drafts/SGA4_IV_LINES_2355_2641_English.texfrag
\subsection{拓扑空间 $X$ 的小拓扑斯与大拓扑斯之间的关系}\etiquette[4.10]{IV.subsec4.10}

沿用 \hyperref[IV.subsec2.5]{2.5} 的记号. 特别地, $𝒱$ 为满足 $𝒰\in 𝒱$ 的宇宙. 与 \hyperref[IV.subsec2.1]{2.1} 的约定不同, 这里用 $\Top(X)$ 表示 $X$ 上的 $𝒱$-层 (而非 $𝒰$-层) 的拓扑斯. 因而我们将使用 $𝒱$-拓扑斯, 而非 $𝒰$-拓扑斯. (注意, 也可以对 $\TOP(X)$ 的定义采用适当约定, 使后者为 $𝒰$-拓扑斯, 从而保留 $𝒰$-拓扑斯, 参见 \hyperref[IV.subsec2.5]{2.5}.) 据此, 我们将定义两个拓扑斯态射
\begin{equation*}
\begin{cases}
 f: \Top(X)\longrightarrow \TOP(X)\\
 g: \TOP(X)\longrightarrow \Top(X)\quav gf\simeq \id_{\Top(X)}{\quoi},
\end{cases}\tag{4.10.1}
\end{equation*}\etiquette[4.10.1]{IV.eq4.10.1}
它们给出三个伴随函子的序列
\begin{equation*}
g^{\ast}=f_{!}\quav g_{\ast}=f^{\ast}\quav
g^{!}=f_{\ast}{\quoi},\tag{4.10.1.1}
\end{equation*}\etiquette[4.10.1.1]{IV.eq4.10.1.1}
其中使用 \hyperref[IV.subsubsec3.1.3]{3.1.3} 的记号. 将依次定义\pageoriginale 此序列的前两个函子, 再将第三个定义为第二个的右伴随.

\setcounter{subsubsection}{1}
\subsubsection{}\etiquette[4.10.2]{IV.subsubsec4.10.2}
函子
$$
g_{\ast}=f^{\ast}:\TOP(X)\longrightarrow \Top(X)
$$
就定义为从 $(\Esp)_{/X}$ 到 $X$ 的开子集景 $\Ouv(X)$ 的\emph{限制函子}; 由定义立即可知, 它确实将层映为层. 对 $X$ 的大景上的层使用带下划线的字母, 对这样的层 $\underline{F}$, 用 $\underline{F}_X$ 表示它在景 $\Ouv(X)$ 上的限制, 从而得到函子
\begin{equation*}
\Restr:\underline{F}\longmapsto \underline{F}_X:\TOP(X)\longrightarrow
\Top(X){\quoi}.\tag{4.10.2.1}
\end{equation*}\etiquette[4.10.2.1]{IV.eq4.10.2.1}
显然此函子与 $𝒱$-投射极限交换, 因为这些极限逐点计算 (也可用下面 \hyperref[IV.subsubsec4.10.4]{4.10.4} 中构造的左伴随的存在性). 我断言它还与 $𝒱$-归纳极限交换. 为看到这一点, 将用 $X$ 上的拓扑空间 $X'$ 上的通常层 (注意这里指 $𝒱$-层), 对 $X$ 上的「大」层, 即 $\TOP(X)$ 的对象, 给出十分方便的解释.

\subsubsection{}\etiquette[4.10.3]{IV.subsubsec4.10.3}
对 $X$ 上的大层 $\underline{F}$, 以及 $X$ 上的每个空间 $X'$ (默认 $X'\in 𝒰$), 如 \hyperref[IV.subsubsec4.10.2]{4.10.2}, 以显然的方式定义「小」层 $\underline{F}_{X'}$, 即 $\underline{F}$ 在 $X'$ 上的限制. 若
$$
u:X''\longrightarrow X'
$$
为\pageoriginale $(\Esp)_{/X}$ 中的态射, 则以显然的方式定义态射 $u_{\ast}(\underline{F}_{X''})\rightarrow \underline{F}_{X'}$, 或者等价地, 定义「\emph{过渡态射}」
\begin{equation*}
\varphi_u :u^{\ast}(\underline{F}_{X'})\longrightarrow
\underline{F}_{X''}{\quoi}. \tag{4.10.3.1}
\end{equation*}\etiquette[4.10.3.1]{IV.eq4.10.3.1}
当 $u$ 变化时, 对 $(\Esp)_{/X}$ 中的态射复合
$$
X'''\xrightarrow{u}X''\xrightarrow{v}X'
$$
这些态射满足一个显然的传递条件, 其明确表述留给读者. 这样得到一个从 $X$ 上的大层范畴 $\TOP(X)$ 到如下系统范畴的自然函子:
$$
(\underline{F}_{X'})\quad(X'\in \ob (\Esp)_{/X}), (\varphi_u)\quad(u\in
Fl(\Esp)_{/X}){\quoi},
$$
这些系统满足刚才描述的传递条件, 并且对每个为\emph{开浸入} (或更一般地为局部同胚) 的态射 $u:X''\rightarrow X'$, 过渡态射 $\varphi_u$ 都为同构. 因为用于描述各 $\varphi_u$ 的函子 $u^{\ast}:\Top(X')\rightarrow \Top(X'')$ 与 $𝒱$-归纳极限交换, 所以在上述用「小」层 $F_{X'}$ 描述 $\TOP(X)$ 的对象的方式下, $𝒱$-归纳极限逐点计算, 即形如 $\underline{F}\rightarrow \underline{F}_{X'}$ 的函子与 $𝒱$-归纳极限交换.

特别地, 对 \hyperref[IV.subsubsec4.10.2]{4.10.2} 中考虑的函子 $\underline{F}\mapsto\underline{F}_X$ 也成立. 因而它确有右伴随 (\hyperref[IV.cor1.5]{1.5}). 还须构造左伴随 (其存在性也先验地由 \hyperref[IV.rem1.8]{1.8} 得出), 并验证后者左正合. 这样\pageoriginale 就通过函子 \eqhyperref[IV.eq4.10.1.1]{4.10.1.1}, 完成了 \eqhyperref[IV.eq4.10.1]{4.10.1} 中所预告的拓扑斯态射的定义.

\subsubsection{}\etiquette[4.10.4]{IV.subsubsec4.10.4}
函子
$$
g^{\ast}=f_{!}:\Top(X)\longrightarrow \TOP(X)
$$
如下得到: 将 $X$ 上的每个小层 $F$ 配以它通过结构态射 $X'\rightarrow X$ 的逆像系统 $(F_{X'})$, $X'\in \ob (\Esp)_{/X}$, 其中对 $u:X''\rightarrow X'$, 过渡态射 $\varphi_u$ 为层的逆像的传递同构 (\hyperref[IV.subsubsec4.1.1]{4.1.1}). 立即可见, 这给出函子
$$
\Prol:\Top(X)\longrightarrow \TOP(X){\quoi},
$$
它\emph{全忠实}, 且其本质像由 $X$ 上使 \sisi{4.11.3}{\eqhyperref[IV.eq4.10.3.1]{4.10.3.1}} 的一切过渡态射 $\varphi_u$ 均为同构的那些大层 $ℱ$ 组成. 在此典范延拓函子与 \hyperref[IV.subsubsec4.10.2]{4.10.2} 的限制函子之间定义伴随同构, 以证明后者为前者的右伴随, 留给读者. 这直接来自 \hyperref[IV.subsubsec4.10.3]{4.10.3} 对范畴 $\TOP(X)$ 的描述.

\begin{enonce*}[remark]{注 4.10.5}\etiquette[4.10.5]{IV.rem4.10.5}
\begin{enumerate}
\renewcommand{\theenumi}{\alph{enumi}}
\renewcommand{\labelenumi}{\theenumi)}
\item \hyperref[IV.subsubsec4.10.4]{4.10.4} 中的构造也表明 $g^{\ast}=f_{!}$ 全忠实 (或者等价地, $g_{\ast}=f^{\ast}$ 为局部化函子, 又或者 $g^{!}=f_{\ast}$ 全忠实). 属于函子 $g^{\ast}=f_{!}=\Prol$ 的本质像的 $X$ 上的大层, 理应称为 $X$ 上的大\emph{平展层}, 因为它们组成的范畴等价于 $X$ 上的通常层范畴, 或等价于\pageoriginale $X$ 上的平展空间范畴. 此外容易看到, 若 $X'$ 为 $X$ 上的拓扑空间, 则它所表的 $X$ 上的大层在上述意义下为平展的, 当且仅当 $X'$ 为 $X$ 上的平展空间.

\item $f_{\ast}$ 的全忠实性也可表述为伴随态射 $f^{\ast}f_{\ast}\rightarrow \id_{\Top(X)}$ 为同构, 即 $g_{\ast}f_{\ast}\simeq \id$, 即 $gf\simeq \id$. 因此 $g$ 使 $\TOP(X)$ 成为 $\Top(X)$ \emph{上的拓扑斯}, 容许 $\Top(X)$ 上的一个「截面」$f$.

\item 函子 $g^{\ast}=f_{!}$ 全忠实、正合且与 $𝒱$-归纳极限交换, 这一事实部分说明了 J. \textsc{Giraud} 提出的一个十分方便的观点的根据: 在 $X$ 上的层论中, 几乎所有问题里都可无害地用相应的「大」层代替通常的层或「小」层. 这尤其适用于上同调问题, 因为 $g_{\ast}$ 正合, 函子 $R^i g_{\ast}$ (\sisi{V\quad}{\hyperrefext[V][V.sec5]{5}}) 当 $i>0$ 时为零, 因而 (\sisi{V\quad}{\hyperrefext[V][V.prop5.4]{5.4}}) 对 $X$ 上的每个\emph{大}层 $ℱ$, 有典范同构
$$
H^i(\TOP(X),ℱ)\simeq H^i(\Top(X),
ℱ_X)=H^i(X,F){\quoi}.
$$
将此应用于形如 $\Prol(F)$ 的层, 其中 $F$ 为 $X$ 上的小层, 就得到 (因为典范地 $\Prol(F)_X\simeq F$) 典范同构
$$
H^i(X,F)=H^i(\TOP(X), \Prol(F)){\quoi}.
$$
因此, \emph{通过 $X$ 的小拓扑斯或大拓扑斯计算的 $X$ 的上同调不变量, 本质上相同}. 此外在非交换上同调中同样成立.
\end{enumerate}
\end{enonce*}

\begin{enonce*}[remark]{练习 4.10.6}\etiquette[4.10.6]{IV.exer4.10.6}
设\pageoriginale $S$ 为拓扑比典范拓扑粗的 $𝒰$-景, $\underline{M}$ 为 $Fl$ $S$ 的子集, 满足以下条件:
\begin{enumerate}
\renewcommand{\theenumi}{\alph{enumi}}
\renewcommand{\labelenumi}{\theenumi)}
\item $ℳ$ 的态射可作拉回 (\sisi{\hyperrefext[I][I.sec10.3]{10.3}}{\hyperrefext[I][I.sec10.7]{10.7}}), 且 $\underline{M}$ 对基变换稳定.

\item $\underline{M}$ 包含恒等箭头, 且对复合稳定.

\item 若箭头 $u:X\rightarrow Y$ 满足: 存在覆盖族 $Y_i\rightarrow Y$, 使各箭头 $\fprod{X}{Y_i}{Y}\rightarrow Y_i$ 属于 $\underline{M}$, 则它本身属于 $\underline{M}$.

\item 对每个 $X\in \ob S$, $X$ 的每个覆盖族都被某个满足 $f_i\in \underline{M}$ 的覆盖族 $(f_i:X_i\rightarrow X)$ 加细.
\end{enumerate}

对每个 $X\in \ob S$, 考虑景 $S(X)$ (「$X$ 的小景」), 它的底范畴为 $S_{/X}$ 中由结构态射属于 $\underline{M}$ 的对象所成的满子范畴, 赋予由 $S$ 的拓扑诱导的拓扑 (\hyperrefext[III][III.subsec3.1]{3.1}).
\begin{itemize}
\item[$1^{\circ})$] 对 $S$ 的每个箭头 $u:X\rightarrow Y$, 证明从 $S(Y)$ 到 $S(X)$ 的沿 $u$ 的基变换为连续函子, 因而给出与归纳极限交换的函子 (\sisi{III}{\hyperrefext[III][III.prop1.2]{1.2} iv)})
$$
S(u)^x:S(Y)^{\sim}\longrightarrow S(X)^{\sim}{\quoi}.
$$

\item[$2^{\circ})$] 定义拓扑斯 $S^{\sim}$ 与以下系统范畴之间的等价:
$$
(F_X)\quad (X\in \ob S)\quav (\varphi_u)\quad (u\in Fl S)
$$
它们\pageoriginale 由对象 $F_X\in \ob S(X)^{\sim}$ 以及对 $S$ 的每个箭头 $u:X\rightarrow Y$ 的态射 $\varphi_u:S(u)^x(F_Y)\rightarrow F_X$ 组成; 这些系统须对 $S$ 的箭头的复合 $v\circ u$ 满足传递条件, 并满足 $u\in \underline{M}$ 蕴涵 $\varphi_u$ 为同构的条件.

\item[$3^{\circ})$] 定义「限制」与「延拓」函子
$$
\Res_X:(S_{/X})^{\sim}\longrightarrow S(X)^{\sim},
\Prol_X:S(X)^{\sim}\longrightarrow (S_{/X})^{\sim}
$$

证明 $\Res_X$ 与小归纳极限及投射极限交换, $\Prol_X$ \emph{全忠实}, 且其本质像由这样的层 $F$ 组成: 对 $S_{/X}$ 的每个箭头 $u$, $\varphi_u$ 都为同构.

\item[$4^{\circ})$] 定义伴随态射, 使 $\Res_X$ 为 $\Prol_X$ 的右伴随. 推出存在拓扑斯态射
$$
f:S(X)^{\sim}\rightarrow (S_{/X})^{\sim}{\quoi}
$$
使得 $S(X)^{\sim}$ 成为 $(S_{/X})^{\sim}$ 的子拓扑斯, 并满足
\begin{align*}
f_{!} &= \Prol_X\\
f^x &= \Res_X{\quoi}.
\end{align*}

证明 $\Prol_X$ 将交换层映为交换层.

\item[$5^{\circ})$] 证明若 $S(X)$ 容许纤维积, 则 $\Prol_X$ 正合. 推出此时存在拓扑斯态射 $g:(S_{/X})^{\sim}\rightarrow S(X)^{\sim}$, 它为 $f$ 的左收缩, 即满足
\begin{align*}
g^x &= \Prol_X\\
g_x &= \Res_X{\quoi}.
\end{align*}
(对\pageoriginale $S(X)$ 不容许纤维积的例子, 取 $S$ 为赋有平展拓扑的概形范畴, $M$ 为光滑态射, $X$ 为维数 $>0$ 的 Noether 概形.)

\item[$6^{\circ})$] 证明对 $(S_{/X})^{\sim}$ 的每个交换层 $F$, 有同构
$$
H^q(X,F)\simeq H^q(X,\Res_X F)\qquad \forall q{\quoi}.
$$
例如用超覆盖证明, 对 $S(X)^{\sim}$ 的每个交换层 $G$, 有同构
$$
H^q(X,G)\simeq H^q(X,\Prol_X G)\qquad \forall q{\quoi}.
$$

\item[$7^{\circ})$] 给编者买一枚巧克力奖章.
\end{itemize}
\end{enonce*}

% END INCLUDED SGA4_English_current_buildable_source_20260906/source/04/drafts/SGA4_IV_LINES_2355_2641_English.texfrag


% BEGIN INCLUDED SGA4_English_current_buildable_source_20260906/source/04/drafts/SGA4_IV_LINES_2642_2987_English.texfrag
\section{诱导拓扑斯}\etiquette[5]{IV.sec5}

\subsection{}\etiquette[5.1]{IV.subsec5.1}
设 $E$ 为拓扑斯, $X$ 为 $E$ 的对象. 则 $X$ 上的 $E$ 的对象范畴 $E_{/X}$ 为拓扑斯, 例如可直接由 Giraud 判别准则 \hyperref[IV.theo1.2]{1.2} ii) 得到. 或者, 由 \hyperref[IV.cor1.2.1]{1.2.1}, 可将 $E$ 实现为层范畴 $C^{\sim}$, 其中 $C$ 为 $E$ 的生成子范畴\pageoriginale, 并可选为使得 $X\in \ob C$; 这时已知 $E_{/X}=C^{\sim}_{/X}$ 等价于 $(C_{/X})^{\sim}$ (\hyperrefext[III][III.prop5.4]{5.4}), 因而为拓扑斯.

称拓扑斯 $E_{/X}$ 为 $E$ 的对象 $X$ 上的\emph{诱导拓扑斯}.

\subsection{}\etiquette[5.2]{IV.subsec5.2}
将定义典范拓扑斯态射
\begin{equation*}
j_X:E_{/X}\longrightarrow E{\quoi},\tag{5.2.1}
\end{equation*}\etiquette[5.2.1]{IV.eq5.2.1}
称为诱导拓扑斯 $E_{/X}$ 到环境拓扑斯 $E$ 的\emph{包含态射}, 或者更恰当地 (参见 \hyperref[IV.subsec5.7]{5.7}) 称为 \emph{$E$ 在 $X$ 处的局部化态射}. 此态射对应于三个伴随函子的序列 (参见 \hyperref[IV.subsubsec3.1.3]{3.1.3})
\begin{equation*}
j_{X!}\quav j^{\ast}_X\quav j_{X^{\ast}}{\quoi},\tag{5.2.2}
\end{equation*}\etiquette[5.2.2]{IV.eq5.2.2}
可明确描述如下:
\begin{enumerate}
\renewcommand{\theenumi}{\alph{enumi}}
\renewcommand{\labelenumi}{\theenumi)}
\item 函子
$$
j_{X!}:E_{/X}\longrightarrow E
$$
是从 $E_{/X}$ 到 $E$ 的\emph{遗忘结构箭头}的函子.

\item 函子
$$
j^{\ast}_X:E\longrightarrow E_{/X}
$$
定义为
$$
j^{\ast}_X(Z)=(X\times Z, \pr_1){\quoi},
$$
其中 $\pr_1:X\times Z\rightarrow X$ 为第一投影. 也可将它解释为相对于态射
$$
X\longrightarrow e_E{\quoi},
$$
的\emph{基变换函子}, 其中 $e_E$ 为 $E$ 的终对象. 由基变换函子的定义, 立即可知 $j^{\ast}_X$ 确为 $j_{X!}$ 的右伴随. 特别地, 它与投射极限交换. 由拓扑斯的特殊正合性质 (\hyperrefext[II][II.sec4]{4}), 它还与归纳极限交换.\pageoriginale

\item 由最后这个事实 (\hyperref[IV.cor1.6]{1.6}), $j^{\ast}_X$ 有右伴随
$$
j_{X^{\ast}}:E_{/X}\longrightarrow E{\quoi}.
$$
\begingroup\tolerance=9999\emergencystretch=4em
对 $X$ 上的对象 $X'$, 对象 $j_{X^{\ast}}(X')$ 有时还记为 $\prod_{X/e_E}(X'/X)$, $\Hhom_{X/e_E}(X,X')$ 或 $\Res_{X/e_E}(X'/X)$ (「Weil 限制」) 之一; 这些记号中想必有一些已为本谦逊著作的博学读者所熟悉.\par
\endgroup
\end{enumerate}

\subsection{}\etiquette[5.3]{IV.subsec5.3}
应注意, 函子 $j_{X!}$ 与纤维积交换, 并将单态射映为单态射, 但并不因此就一般地左正合. 它只有在 $X$ 为 $E$ 的终对象时才左正合 (因为 $X=j_{X!}(X,\id_X)$, 而 $(X,\id_X)$ 为 $E_{/X}$ 的终对象), 即当且仅当 $j_X$ 事实上是拓扑斯等价 (从而 \eqhyperref[IV.eq5.2.2]{5.2.2} 中的函子全为等价).\footnote[1]{关于 $j_{X!}$ 在拓扑斯变换下的交换性质, 参见 XVII \hyperref[XVII.lem5.1.2]{5.1.2}.}

同样, 函子 $j_{X^{\ast}}$ 一般不右正合, 也不一定把满态射映为满态射.

由这些观察可知, 对拓扑斯 $E$ 的对象 $X$, 联系 $E_{/X}$ 与 $E$ 的拓扑斯态射的方向毫无歧义地确定.

\subsection{}\etiquette[5.4]{IV.subsec5.4}
函子 $j^{\ast}_X$ 常称为\emph{局部化函子}或\emph{限制函子}. 后一种术语的根据在于, 把 $E$ 及 $E_{/X}$ 的对象分别等同于 $E$ 及 $E_{/X}$ 上的层 (\hyperref[IV.theo1.2]{1.2} iii)), 并注意在这种等同下, $j_{X!}$ (遗忘函子) 与 $j^{\ast}_X$ 之间的伴随公式, 可解释为 $j^{\ast}_X(F)$ \emph{是\pageoriginale $E$ 上的层 $F$ 在 $E_{/X}$ 上的限制} (这里「限制」取广义: 与「包含」函子 $j_{X!}:E_{/X}\rightarrow E$ 复合). 按照其他情境中熟悉的记号, 我们也将经常互换地写:
\begin{equation*}
j^{\ast}_X(F)=F| X=F_X=\mbox{ } F\mbox{ 在 }
X\mbox{ 上的限制}{\quoi}.\tag{5.4.1}
\end{equation*}\etiquette[5.4.1]{IV.eq5.4.1}

同样, 当 $E$ 形如 $C^{\sim}$, 其中 $C$ 为 (小) 景, 且 $X$ 形如 $\epsilon(S)$, 其中 $S\in\ob C$, 从而有 \hyperref[IV.subsec5.1]{5.1} 中回顾的范畴等价
$$
E_{/X}=C^{\sim}_{/\epsilon(S)}\xrightarrow{\approx}(C_{/S})^{\sim}
$$
(其中 $C_{/S}$ 赋予 $C$ 的拓扑诱导的拓扑), 函子 $j^{\ast}_X$ 就等同于将 $C$ 上变化的层 $F$「限制」到范畴 $C_{/S}$ 的函子.

这些讨论还使人预见, 在导致我们用「$E$ 上的局部化」来论证的所有问题中 (参见第 8 节), 即实际上在一切涉及景或拓扑斯的问题中, 诱导拓扑斯与局部化态射 \eqhyperref[IV.eq5.2.1]{5.2.1} 将起重要作用.

\subsection{}\etiquette[5.5]{IV.subsec5.5}
设
$$
f:X\longrightarrow Y
$$
为 $E$ 的箭头, 因而可将 $X$ (或更确切地说 $(X,f)$) 看作 $E_{/Y}$ 的对象. 显然有典范同构
\begin{equation*}
(E_{/Y})_{/X}\xrightarrow{\sim}E_{/X}\tag{5.5.1}
\end{equation*}\etiquette[5.5.1]{IV.eq5.5.1}
(\emph{诱导拓扑斯的传递性}). 以 $E_{/Y}$ 代替 $E$ 应用 \hyperref[IV.subsec5.1]{5.1} 的构造, 就得到典范拓扑斯态射
\begin{equation*}
\loc(f)\mbox{ 或 } f:E_{/X}\longrightarrow E_{/Y}{\quoi},\tag{5.5.2}
\end{equation*}\etiquette[5.5.2]{IV.eq5.5.2}
也称为\pageoriginale \emph{(与 $f$ 相联系的) 局部化态射}. 当 $Y$ 为 $E$ 的终对象时, 实质上重新得到 \eqhyperref[IV.eq5.2.1]{5.2.1}. $f$ 的局部化态射与三个伴随函子的序列相联系:
\begin{equation*}
f_{!}\quav f^{\ast}\quav f_{\ast}{\quoi},\tag{5.5.3}
\end{equation*}\etiquette[5.5.3]{IV.eq5.5.3}
它们分别可解释为遗忘函子、限制函子, 以及一个可随意 (以 $X'$ 为自变量) 记为 $\prod_{X/Y}(X'/X)$, $\Hhom_{X/Y}(X,X')$ 或 $\Res_{X/Y}(X'/X)$ 的函子.

\subsection{}\etiquette[5.6]{IV.subsec5.6}
由遗忘函子的传递性, 对 $E$ 的两个可复合态射
$$
X\xrightarrow{f} Y \xrightarrow{g} Z
$$
有拓扑斯态射的典范同构
$$
\loc(gf)\simeq \loc(g)\loc(f):E_{/X}\longrightarrow
E_{/Y}\longrightarrow E_{/Z}{\quoi}.
$$
此外, 对三个可复合态射, 上述传递同构满足通常的相容关系. 因而在惯常的语言滥用下, 可将
$$
X\longmapsto E_{/X}
$$
看作从 $E$ 到范畴 ($𝒱$-$𝒰$-$\text{top}$) 的协变函子 (\hyperref[IV.subsubsec3.3.1]{3.3.1}). 更确切地, 得到一个从 $E$ 到 2-范畴 ($𝒱$-$𝒰$-$\text{top}$) 的 2-函子 (一般不严格).

在继续诱导拓扑斯的一般理论之前 (\hyperref[IV.subsec5.10]{5.10}), 先给出几个有启发性的例子.

\subsection{}\etiquette[5.7]{IV.subsec5.7}
设\pageoriginale $X$ 为拓扑空间, 因而有拓扑斯 $\Top(X)$ (\hyperref[IV.subsec2.1]{2.1}). 设 $X'$ 为 $\Top(X)$ 的对象, 将它解释为 $X$ 上的平展空间 $p:X'\rightarrow X$ 较为方便. 此时范畴 $\Top(X)_{/X'}$ 等同于配有 $X$-态射 $X''\rightarrow X'$ 的 $X$ 上平展空间 $X''$ 的范畴. 已知这样的态射使 $X''$ 成为 $X'$ 上的平展空间, 因而得到拓扑斯的典范等价
$$
\Top(X)_{/X'}\xrightarrow{\approx} \Top(X'){\quoi}.
$$
于是局部化态射等同于一个拓扑斯态射 $\Top(X')\rightarrow \Top(X)$, 立即可见后者正是由结构连续映射 $p:X'\rightarrow X$ 所联系出的态射
$$
\Top(p):\Top(X')\longrightarrow \Top(X)
$$
因此直观地说, 局部化态射就是将平展映射 $p:X'\rightarrow X$ 翻译成拓扑斯的语言. 这既说明了「局部化态射」这一术语的恰当性, 也提示须慎用「包含态射」这一术语, 因为它似乎主要适用于 $p$ 为开浸入的情形. 因而在 \hyperref[IV.subsec5.1]{5.1} 中, 将「包含态射」保留给态射 $X\rightarrow e_E$ 为单态射, 即 $X$ 等同于 $E$ 的终对象的子对象的情形, 是审慎的.

\subsection{}\etiquette[5.8]{IV.subsec5.8}
设 $E$ 为拓扑斯, $G$ 为 $E$ 中的群, $H$ 为 $G$ 的子群, 并设
$$
X=G/H
$$
为商齐性空间, 看作有 $G$ 左作用的 $E$ 的对象, 即分类拓扑斯 $B_G$ 的对象 (\hyperref[IV.subsec2.4]{2.4}). 我们\pageoriginale 将确定诱导拓扑斯
$$
(B_G)_{/X}\quad (X=G/H){\quoi},
$$
为此定义拓扑斯等价
\begin{equation*}
c:B_H \xrightarrow{\approx} (B_G)_{/X}{\quoi},\tag{5.8.1}
\end{equation*}\etiquette[5.8.1]{IV.eq5.8.1}
使得下列图表在典范同构意义下交换:
\[\vcenter{
\xymatrix@C=2cm{ B_H \ar[r]^{c} \ar[r] \ar
@{^{(}->}@<-9pt>[rr]_{B_i}& {(B_{G})_{/X}} \ar[r]^{j_X}& B_G
}}\leqno{(5.8.2)}
\]\etiquette[5.8.2]{IV.eq5.8.2}
其中 $j_X$ 为 $B_G$ 中的局部化态射,
$$
B_i:B_H\longrightarrow B_G
$$
为由包含 $i:H\rightarrow G$ 推出的态射 (\hyperref[IV.subsec4.5]{4.5}).

为定义 \eqhyperref[IV.eq5.8.1]{5.8.1}, 我们只描述逆像函子 $c^{\ast}$ 与正像函子 $c_{\ast}$, 把验证它们确实互为拟逆等价、且给出拓扑斯态射的交换图表 \eqhyperref[IV.eq5.8.2]{5.8.2} 留给读者. 设 $e$ 为 $E$ 中底层对象 $X$ 的子对象, 它是从 $G$ 的单位截面得出的 $X$ 的截面 (在 $E$ 的一个选定的终对象 $e_E$ 上) 的像. 若 $X'$ 为 \sisi{$G$}{$X$} 上的 $B_G$ 的对象, 则用 $c^{\ast}(X')$ 表示 $e$ 在 $X'$ 中的逆像,
$$
c^{\ast}(X')=\fprod{X'}{e}{X}{\quoi},
$$
它在限制 $G$ 于 $X'$ 的给定作用所得到的 $H$ 于 $X'$ 的左作用下稳定. 这确实定义了函子
$$
c^{\ast}:(B_G)_{/X}\longrightarrow B_H{\quoi}.
$$
另一方面\pageoriginale, 若 $X'$ 为 $B_H$ 的对象, 则从 $X'$ 到 $B_H$ 的终对象 $e_H$ (即 $E$ 的终对象 $e_E$ 配上 $H$ 的平凡作用) 的态射, 经函子 $i_{!}$ (\hyperref[IV.subsec4.5]{4.5}) 给出 $B_G$ 中的态射
$$
i_{!}(X')\longrightarrow i_{!}(e_H)\simeq X{\quoi},
$$
使得可以把 $i_{!}(X')$ 看作 $(B_G)_{/X}$ 的对象, 因而给出所求的函子
$$
c_{\ast}:X'\longmapsto (i_{!}(X')\longrightarrow
X):B_H\longrightarrow (B_G)_{/X}{\quoi}.
$$

\setcounter{subsubsection}{2}
\subsubsection{}\etiquette[5.8.3]{IV.subsubsec5.8.3}
因此, 由 \eqhyperref[IV.eq5.8.2]{5.8.2} 可见, 若 $G$ 为拓扑斯中的群, $H$ 为其子群, 则\emph{「将算子从 $G$ 限制到 $H$」的函子可解释为局部化函子}. 特别地, 取 $H$ 为单位子群, 可知\emph{遗忘 $G$ 作用}的函子可解释为局部化函子. 更一般地由此可得, 若 $(X,G)$ 为有 $G$ 作用的 $E$ 的对象 (即分类拓扑斯 $B_G$ 的对象), 则「遗忘 $G$ 作用」的函子
$$
E'=(B_G)_{/(X,G)}\longrightarrow E_{/X}
$$
可解释为相对于 $E'$ 中覆盖终对象的某个适当对象 $E_{G,(X,G)}$ 的局部化函子. 只须取 $E_{G,(X,G)}=E_G\times (X,G)$, 给出笛卡儿图表 (其中 $E_G=G_S=G$ 配有 $G$ 的左平移作用)
\[\vcenter{
\xymatrix@C=2cm@R=1.5cm{
E_G \ar[d]& **[r]E_{G,(X,G)} = Z \ar[l] \ar[d]\\
e_G & (X,G) \ar[l]          \quad , }}\leqno{(5.8.3.1)}
\]\etiquette[5.8.3.1]{IV.eq5.8.3.1}
并\pageoriginale 注意由诱导拓扑斯的传递性, 诱导拓扑斯 $E'=((B_G)_{X/G)/Z}$ 等价于拓扑斯 $((B_G)_{/E_G})_{/Z}$; 因为 $(B_G)_{/E_G}$ 通过 \eqhyperref[IV.eq5.8.1]{5.8.1} 的函子 $c^{\ast}$ (取 $H=e$) 等价于 $E$, 只须验证此等价将 $Z$ 映为 $E$ 的对象 $X$ (这显然成立), 就推出所求范畴等价
$$
E'_{/Z}\xrightarrow{\approx}E_{/X}{\quoi}.
$$
把验证此等价与局部化函子
$$
E'=(B_G)_{/(X,G)}\longrightarrow E'_{/Z}
$$
的复合为「遗忘 $G$ 作用」的函子留给读者. 这样可见, 通过取诱导拓扑斯, 上面的笛卡儿图表给出一个在典范同构意义下交换的拓扑斯态射图表 (\hyperref[IV.subsec5.6]{5.6}):
\[\vcenter{
  \xymatrix@C=1.25cm@R=1cm{
**[r](B_G)_{/E_G}\approx E \ar[d]_{j} & **[r](B_G)_{/Z} \approx
E'_{/Z} \ar[l]_(.3){i}
  \approx E_{/X} \ar[d]^{j'}\\
B_G & **[r](B_G)_{/(X,G)} =  E' \ar[l]_(.5){f}\quad,
}}\leqno{(5.8.3.2)}
\]\etiquette[5.8.3.2]{IV.eq5.8.3.2}
其中与竖直箭头 $j$, $j'$ 相联系的逆像函子为「遗忘 $G$ 作用」的函子, 逆像函子 $i^{\ast}$ 则等同于局部化函子 $E\rightarrow E_{/X}$. 此图表\pageoriginale 还在下面 \hyperref[IV.prop5.11]{5.11} 的意义下为「2-笛卡儿」图表.

\subsubsection{}\etiquette[5.8.4]{IV.subsubsec5.8.4}
把上述讨论推广到 pro-群 $𝒢=(G_i)$ 的情形, 其中限于 $H$ 定义为 $𝒢\rightarrow G_{i_{\circ}}$ 的「核」的特殊情形, 留给读者. (对 pro-有限群, 这表示只考虑 $G$ 的\emph{开}子群 $H$, 即闭且具有\emph{有限指数}的子群.)

% END INCLUDED SGA4_English_current_buildable_source_20260906/source/04/drafts/SGA4_IV_LINES_2642_2987_English.texfrag


% BEGIN INCLUDED SGA4_English_current_buildable_source_20260906/source/04/drafts/SGA4_IV_LINES_2988_3412_English.texfrag
\begin{enonce*}[remark]{练习 5.9}\etiquette[5.9]{IV.exer5.9}
设 $E$ 为拓扑斯, $G$ 为 $E$ 的群, $B_G$ 为 $G$ 的分类拓扑斯 (\hyperref[IV.subsec2.4]{2.4}), 并设
$$
\pi:B_G\longrightarrow E
$$
为由从 $G$ 到 $E$ 的单位群 $e$ 的同态推出的拓扑斯态射 (这里考虑到 $B_e\approx E$).
\begin{enumerate}
\renewcommand{\theenumi}{\alph{enumi}}
\renewcommand{\labelenumi}{\theenumi)}
\item 设
$$
E_G=G_s
$$
为 $B_G$ 的对象, 其底层 $E$-对象为 $G$, $G$ 的作用由左平移定义. 将 $\pi^{\ast}(G)$ 看作 $B_G$ 中的群 (它是 $E$ 的群 $G$, 配以 $G$ 在其上的平凡作用). 证明由 $G$ 在 $G_s$ 上的右平移定义的 $E$ 的态射
$$
G_s\times G\longrightarrow G_s
$$
与 $G$ 在 $G_s$ 及 $G=\pi^{\ast}(G)$ 上的作用相容, 并定义 $B_G$ 的态射
$$
E_G\times \pi^{\ast}(G)\longrightarrow E_G{\quoi}.
$$
证明此态射给 $E_G$ 赋予 $B_G$-群 $\pi^{\ast}(G)$ 的右作用, 且此作用使 $E_G$ 成为 $\pi^{\ast}(G)$ 下的\pageoriginale \emph{挠子} (即 $B_G$ 的终对象可由对象 $X_i$ 覆盖, 使对每个 $i$, $E_G$ 在 $X_i$ 上的限制同构于带算子的平凡丛 $(\pi^{\ast}(G)_{X_i})_d$).

\item 设
$$
q:E'\longrightarrow E
$$
为 $E$ 上的拓扑斯. 对 $E$ 上的每个拓扑斯 $B$, 定义从 $E'$ 到 $B$ 的拓扑斯态射的范畴 $\Hhomtop_E(E',B)$, 要求它们在一个指定的同构的意义下与结构态射 $E'\rightarrow E$ 及 $B\rightarrow E$ 相容. 取 $B=B_G$, 用公式 $f\mapsto f^{\ast}(E_G)$ 定义范畴等价
$$
\Hhomtop_E(E',B_G)\longrightarrow \Ttors(E',q^{\ast}(G)){\quoi},
$$
其中右边表示 $E'$ 上的 (右) $q^{\ast}(G)$-挠子范畴.

\item 设 $H$ 为 $G$ 的子群, 从而 $\pi^{\ast}(H)$ 为 $\pi^{\ast}(G)$ 的子群, 因而通过限制算子群而作用于 $E_G$, 给出商对象
$$
X=E_G/\pi^{\ast}(H){\quoi},
$$
它就是对象 $G/H$ 配以通常的 $G$ 左作用. 设 $e_G$ 为 $B_G$ 的终对象 (即 $E$ 的终对象 $e_E$ 配以 $G$ 的必然平凡的作用), 并在 $B_G$ 中考虑态射图表
\[
\xymatrix@C=-1cm@R=.75cm{
E_G \ar[dd] \ar[rd]& \\
    & \qquad\qquad\quad X = E_{G}/\pi^{*}(H)\ar[ld]\\
e_G & }
\]
它在\pageoriginale 转为 $B_G$ 在这些对象上的诱导拓扑斯后, 给出一个在典范同构意义下交换的拓扑斯态射图表 (\hyperref[IV.subsec5.6]{5.6}). 证明此图表等价于图表
\[
\xymatrix@C=.05cm@R=.85cm{
E & = & B_e \ar[dd] \ar[rrrd]& & &\\
  &   &                    &  & &B_H \ar[llld]\\
  &   & B_G                &  & & &,
}
\]
其三个箭头为 \hyperref[IV.subsec4.5]{4.5} 中与 $E$ 的群态射 $e\rightarrow H\rightarrow G$ 相联系的箭头. (因而, \emph{分类拓扑斯 $B_H$ 可直观地解释为 $B_G$ 上的, 以 $\pi^{\ast}(G)$ 为群, 与 $B_G$ 上的泛挠子} (= 齐性主丛) \emph{$E_G$ 相联系的齐性空间}.)

\item 当 $X'$ 为 $X$ 的\emph{平展覆盖} (局部平凡, 参见 \hyperref[IV.subsubsec2.7.4]{2.7.4}) 时, 根据本练习的讨论解释 \hyperref[IV.subsec5.7]{5.7} 的例子, 并重新考察它.
\end{enumerate}
\end{enonce*}

\setcounter{subsection}{9}
\subsection{诱导拓扑斯的逆像}\etiquette[5.10]{IV.subsec5.10}

设
$$
f:E'\rightarrow E
$$
为 $𝒰$-拓扑斯态射, $X$ 为 $E$ 的对象, 并设 $X'=f^{\ast}(X)$ 为它在 $E'$ 中的逆像. 我们将定义拓扑斯态射图表
\[\vcenter{
\xymatrix@C=2.5cm@R=1.5cm{
E_{/X}  \ar[d]_{j_X} & E'_{/X'} \ar[l]_{f_{/X}} \ar[d]^{j_{X'}} \\
E      &  E'\ar[l]_{f}\quad , }}\leqno{(5.10.1)}
\]\etiquette[5.10.1]{IV.eq5.10.1}
其中\pageoriginale $j_X$ 与 $j_{X'}$ 为局部化态射 (\hyperref[IV.subsec5.2]{5.2}), $f_{/X}$ 则在下面定义; 此图表将在典范同构意义下交换. 这里用逆像函子
\begin{equation*}
(f_{/X})^{\ast}:E_{/X}\longrightarrow E'_{/X'}{\quoi},\tag{5.10.2}
\end{equation*}\etiquette[5.10.2]{IV.eq5.10.2}
定义 $f_{/X}$, 将该函子取为在显然意义下由 $f^{\ast}$「诱导」的函子. 因为包含函子 $j_{X!}:E_{/X}\rightarrow E$ 与 $j_{X'!}:E'_{/X'}\rightarrow E'$ 和 $𝒰$-归纳极限及纤维积交换, 且为保守函子, 所以立即可知函子 $(f_{/X})^{\ast}$ 与 $𝒰$-归纳极限和纤维积交换, 正如诱导它的函子 $f^{\ast}$ 一样. 此外, $(f_{/X})^{\ast}$ 显然把终对象 $X$ 映为终对象 $X'$, 所以它左正合, 从而对应于拓扑斯态射
\begin{equation*}
f_{/X}:E'_{/X'}\longrightarrow E_{/X}{\quoi},\tag{5.10.3}
\end{equation*}\etiquette[5.10.3]{IV.eq5.10.3}
它在唯一同构意义下定义 (\hyperref[IV.subsubsec3.1.1]{3.1.1}). 另一方面, 由构造及 $f^{\ast}$ 与积 $X\times Y$ 交换这一事实, 从 \eqhyperref[IV.eq5.10.1]{5.10.1} 转为逆像函子所得到的图表 (在典范同构意义下) 交换. 因而 \eqhyperref[IV.eq5.10.1]{5.10.1} 本身也在唯一同构意义下交换, 该同构在两个复合 $f\circ j_{X'}$ 与 $j_X\circ f_{/X}$ 的逆像函子上诱导典范同构\hyperlink{sga4.exp.IV.subsubsection.3.2.1}{(3.2.1)}.

\setcounter{subsubsection}{3}
\subsubsection{}\etiquette[5.10.4]{IV.subsubsec5.10.4}
当 $f:E'\rightarrow E$ 对应于拓扑空间的连续映射 $Y'\rightarrow Y$ (\hyperref[IV.subsec4.1]{4.1}), 从而 $X$ 等同于 $Y$ 上的平展空间, $X'$ 等同于纤维积 $\fprod{X}{Y'}{Y}$ 时, 将诱导拓扑斯 $\Top(Y)_{/X}$ 与 $\Top(Y')_{/X'}$ 分别等同于 $\Top(X)$ 与 $\Top(X')$ (\hyperref[IV.subsec5.7]{5.7}), 可知刚构造的态射 $f_{/X}$\pageoriginale 就是 (在唯一同构意义下) 从典范连续映射 $X'\rightarrow X$ 推出的态射 $\Top(X')\rightarrow\Top(X)$. 由此例可见, 可以说拓扑斯 $E'_{/X'}$ 起着 $E_{/X}$ 与 $E'$ 在 $E$ 上的纤维积的作用. 以下结果使这一直观精确化:

\begin{enonce*}{命题 5.11}\etiquette[5.11]{IV.prop5.11}
用 \hyperref[IV.subsec5.10]{5.10} 的记号, 图表 \eqhyperref[IV.eq5.10.1]{5.10.1} 配以相容同构 $\alpha:f\circ j_{X'}\rightarrow j_X\circ f_{/X}$, 为「$2$-笛卡儿」图表. 更确切地说, 对每个 $𝒰$-拓扑斯 $F$, 若将每个拓扑斯态射 $g:F\rightarrow E'_{/X'}$ 配以三元组 $(f_{/X}\circ g, j_{X'}\circ g,\alpha\ast g)=(g_1,g_2,\beta)$, 其中 $g_1:F\rightarrow E_{/X}$ 与 $g_2:F\rightarrow E'$ 为拓扑斯态射, $\beta:j_X\circ g_1\rightarrow f\circ g_2$ 为从 $F$ 到 $E$ 的拓扑斯态射的同构, 则得到范畴 $\Hhomtop(F,E'_{/X'})$ 与上述一切三元组 $(g_1,g_2,\beta)$ 所成的范畴
\sisi{$\fprod{\Hhomtop(F,E_{/X})^2}{\Hhomtop(F,E')}{\Hhomtop(F,E)}$}{$\Hhomtop(F,E_{/X})\times^{2}_{\Hhomtop(F,E)}\Hhomtop(F,E')$}
之间的等价.
\end{enonce*}

(注意. 在笛卡儿积符号上方写 $2$, 是为了提醒这不是普通的范畴纤维积, 而是在陈述中所明确的意义下的「$2$-纤维积」.)

为证明 \hyperref[IV.prop5.11]{5.11}, 将通过相应的逆像函子明确描述拓扑斯态射. 为此需要以下 \hyperref[IV.subsubsec5.11.1]{5.11.1} 至 \hyperref[IV.prop5.12]{5.12} 中的辅助结果.

\setcounter{subsection}{11} \setcounter{subsubsection}{0}
\subsubsection{}\etiquette[5.11.1]{IV.subsubsec5.11.1}
一般地, 考虑给定两个范畴 $E$ 与 $E'$, 以及 $E$ 的对象 $X$ 的情形, 假设 $X$ 可作拉回, 即函子
$$
j:A\longmapsto A\times X:E\longrightarrow E_{/X}
$$
有\pageoriginale 定义, 最后给定左正合函子
\begin{equation*}
\varphi:E_{/X}\longrightarrow E'{\quoi}.\tag{$\ast$}
\end{equation*}
因为 $E_{/X}$ 有终对象, 即 $(X,\id_X)=X$, 所以 $E'$ 也有终对象
$$
e'\simeq \varphi(X){\quoi}.
$$
选定 $E'$ 的终对象 $e'$ 后, 将上述每个 $\varphi$ 配以一对
\begin{equation*}
(\Psi,u),\mbox{ 其中 } \Psi=\varphi\circ j:E\rightarrow E',\mbox{ 且
} u\in \Hom(e',\Psi(X)){\quoi}.\tag{${\ast}{\ast}$}
\end{equation*}
为定义 $u$, 注意
$$
j(X)=X\times X
$$
在 $E_{/X}$ 的终对象 $X$ 上有一个典范截面 $\delta_X$, 即对角截面, 并定义
$$
u=\varphi(\delta_X):\varphi(X)=e'\longrightarrow \varphi(X\times
X)=\Psi(X){\quoi}.
$$
我断言知道对 $({\ast}{\ast})$, 就能在唯一同构意义下重构 $(\ast)$ 的函子 $\varphi$. 为此, 对 $E_{/X}$ 的每个对象
$$
p:X'\longrightarrow X
$$
考虑 $E_{/X}$ 中的下列笛卡儿图表:
\[
\xymatrix@C=1cm@R=1.5cm{
X' \ar[d] \ar[r] & **[r]j(X') = X'\times X \ar[d]^{j(p)}\\
X  \ar[r]^{\delta_X}& **[r]j(X) = X \times X   &, }
\]
其中第一个水平箭头为图态射 $\Gamma_p=(\id_{X'},p)$. 对此图表应用左正合函子 $\varphi$, 并用 $\Psi$ 定义为 $\varphi\circ j$, 得到笛卡儿图表
\[
\xymatrix@C=2.25cm@R=1.25cm{
\varphi(X',p) \ar[d] \ar[r]& \Psi (X') \ar[d]^{\Psi(p)}\\
e'  \ar[r]          & \Psi (X) \quad , }
\]
或者\pageoriginale 换言之, 典范同构
\begin{equation*}
\varphi(X',p)\simeq
\fprod{\Psi(X')}{e'}{\Psi(X)}{\quoi}\tag{5.11.1.1}
\end{equation*}\etiquette[5.11.1.1]{IV.eq5.11.1.1}
立即可见它关于 $E_{/X}$ 的对象 $(X',p)$ 是函子性的, 这明确描述了如何从 $({\ast}{\ast})$ 的对 $(\Psi,u)$ 在同构意义下重构函子 $\varphi:E_{/X}\rightarrow E'$. 还注意到函子 $\Psi$ 左正合. 更一般地, $\Psi$ 与 $\varphi$ 所交换的每种投射极限交换, 因为 $j$ 与投射极限交换.

反之, 现在假设有限投射极限在 $E$ 与 $E'$ 中可表, 并从一对
$$
(\Psi,u),\Psi:E\longrightarrow E', u\in \Hom(e',\Psi(X)){\quoi},
$$
出发, 其中函子 $\Psi$ 左正合. 此时公式 \eqhyperref[IV.eq5.11.1.1]{5.11.1.1} 定义函子 $\varphi:E_{/X}\rightarrow E'$. 立即可验证, 此函子左正合, 更一般地, 与在 $E$ 中可表且与 $\Psi$ 交换的每种投射极限交换. 这直接由联系\pageoriginale 在 $E$ 与 $E_{/X}$ 中所计算的投射极限的笛卡儿图表
\[
\xymatrix@C=2.5cm@R=1cm{ \displaystyle{\varprojlim^{E}_{I} X'_i}
\ar[d] &
\displaystyle{\varprojlim^{E_{/X}}_{I} X'_i} \ar[d] \ar[l]\\
\displaystyle{\varprojlim^{E}_{I} X} & X \ar[l]_{\delta} }
\]
得出 (其中 $\delta$ 表示到取值 $X$ 的常值函子的投射极限的对角态射): 对其应用函子 $\Psi$, 得到 $E'$ 中的笛卡儿图表, 再与从 $\Psi(X)\leftarrow e'$ 推出的笛卡儿方块复合. 所得复合笛卡儿方块表明了 $\varphi$ 与所考虑投射极限的相容性.

注意有典范同构
\begin{equation*}
\Psi(X')=\varphi(j(X')){\quoi},\tag{5.11.1.2}
\end{equation*}\etiquette[5.11.1.2]{IV.eq5.11.1.2}
就能从 $\varphi$ 在同构意义下重构对 $(\Psi,u)$; 这个同构通过在 \eqhyperref[IV.eq5.11.1.1]{5.11.1.1} 中用 $j(X')=(X'\times X,\pr_2)$ 代替 $(X',p)$, 并注意 $\Psi(X'\times X)\simeq\Psi(X')\times\Psi(X)$ 得出. 上述同构显然关于 $X'$ 是函子性的, 即给出同构
$$
\Psi\simeq \varphi\circ j{\quoi},
$$
同样可验证, 在此同构下, $u:e'\rightarrow\Psi(X)$ 等同于 $\varphi(\delta_X)$. 至此得到了以下结果的主要部分.

\begin{enonce*}{引理 5.11.2}\etiquette[5.11.2]{IV.lem5.11.2}
设 $E$ 与 $E'$ 为有限投射极限可表的两个范畴, $X$ 为 $E$ 的对象, $\Ssex(E,E')$ 为从 $E$ 到 $E'$ 的左正合函子 $\Psi$ 的范畴, $\Ssex(E_{/X},E')$ 为类似的从 $E_{/X}$ 到 $E'$ 的左正合函子 $\varphi$ 的范畴. 则范畴 $\Ssex(E_{/X},E')$ 与由对 $(\Psi,u)$ 组成的范畴 $\Ssex(E,E')_{/\Gamma}$ 之间有等价, 其中 $\Psi\in\ob\Ssex(E,E')$, $u:e'\rightarrow\Psi(X)$ ($e'$ 为 $E'$ 的终对象). 此等价及一个拟逆函子在 \hyperref[IV.subsubsec5.11.1]{5.11.1} 中已有明确描述. 若 $\varphi$ 与 $(\Psi,u)$ 相对应, 则 $\varphi$ 与某种给定投射极限交换\pageoriginale 当且仅当 $\Psi$ 也如此. 若 $E$ 与 $E'$ 中的基变换函子与某种给定归纳极限交换, 则 $\varphi$ 与该种归纳极限交换当且仅当 $\Psi$ 也如此.
\end{enonce*}

使 $\Ssex(E,E')_{/\Gamma}$ 对应于对 $(\Psi,u)$ 的范畴结构的明确说明, 留给读者. 若用 $\Gamma(\Psi)=\Hom(e',\Psi(X))$ 定义函子 $\Gamma:\Ssex(E,E')\rightarrow(\Ens)$, 则对 $(\Psi,u)$ 中的 $u$ 可解释为 $\Gamma(\Psi)$ 的元素, 即 $\Ssex(E,E')^{\circ}$ 上的预层范畴中的同态 $\Psi\rightarrow\Gamma$. 这说明了 \hyperref[IV.subsubsec5.11.1]{5.11.1} 陈述中所用记号 ${}_{/\Gamma}$ 的根据. 为完成 \hyperref[IV.lem5.11.2]{5.11.2} 的证明, 还须在 \hyperref[IV.subsubsec5.11.1]{5.11.1} 中明确所考虑的构造在 $\varphi$, 分别在 $(\Psi,u)$, 变化时的函子性. 这基本上是显然的, 留给读者. 还须验证关于与归纳极限交换的断言; 这也直接由联系这些函子的显式公式 \eqhyperref[IV.eq5.11.1.1]{5.11.1.1} 与 \eqhyperref[IV.eq5.11.1.2]{5.11.1.2} 得出.

特别地, 当 $E$ 与 $E'$ 为 $𝒰$-拓扑斯, 并用相应逆像函子解释拓扑斯态射时, 得到:

\begin{enonce*}{命题 5.12}\etiquette[5.12]{IV.prop5.12}
设 $E$ 与 $E'$ 为两个 $𝒰$-拓扑斯, $X$ 为 $E$ 的对象, 并在 $\Hhomtop(E',E)$ 上考虑反变函子
$$
\gamma:f\longmapsto
\Gamma(E',f^{\ast}(X))=\Hom(e',f^{\ast}(X)):\Hhomtop(E',E)^{\circ}\longrightarrow
(\Ens){\quoi}.
$$
则有范畴等价
\begin{equation*}
\Hhomtop(E',E_{/X})\xrightarrow{\approx}\Hhomtop(E',E)_{/\gamma}{\quoi},\tag{5.12.1}
\end{equation*}\etiquette[5.12.1]{IV.eq5.12.1}
其中\pageoriginale 右边为 $\Hhomtop(E',E)^{\widehat{\phantom{,}}}_{/\gamma}$ 中由满足 $f\in\Homtop(E',E)$ 及 $u\in\gamma(f)$, 即 $u\in\Gamma(E',f^{\ast}(X))$ 的对 $(f,u)$ 所成的满子范畴. 此函子如下得到: 将每个拓扑斯态射 $h:E'\rightarrow E_{/X}$ 配以由复合
$$
f=j_X\circ h:E'\rightarrow E_{/X}\rightarrow E{\quoi},
$$
以及态射
$$
u:e'=h^{\ast}(X,\id_X)\longrightarrow
f^{\ast}(X)=h^{\ast}(j^{\ast}_X(X))=h^{\ast}(X\times X)
$$
所成的对 $(f,u)$; 后一个态射由对角态射 $\delta_X:X\rightarrow X\times X$ 应用 $h^{\ast}$ 得出.
\end{enonce*}

\setcounter{subsection}{12}
\subsubsection{}\etiquette[5.12.2]{IV.subsubsec5.12.2}
用 \hyperref[IV.prop5.12]{5.12}, \hyperref[IV.prop5.11]{5.11} 的证明就变得几乎显然, 实质上归结为以下内容 (它代替一份并不会更有启发性的形式证明): 给定拓扑斯态射 $g:F\rightarrow E'_{/X'}$, 与给定拓扑斯态射 $g_2:F\rightarrow E'$ 及 $g^{\ast}_2(X')$ 的截面 $u$, 「是一回事」. 但 $g^{\ast}_2(X')=g^{\ast}_2(f^{\ast}(X))=(fg_2)^{\ast}(X)$, 所以给定 $u$ 又等价于给定拓扑斯态射 $fg_2:F\rightarrow E$ 到拓扑斯态射 $g_1:F\rightarrow E_{/X}$ 的提升. 这恰好表明给定拓扑斯态射 $g:F\rightarrow E'_{/X'}$, 实质上等价于给定 \hyperref[IV.prop5.11]{5.11} 中那样的三元组 $(g_1,g_2,\alpha)$.

\begin{enonce*}[remark]{注 5.13}\etiquette[5.13]{IV.rem5.13}
我们将看到 (\S 15\nde{这一节并不存在. 不过, 关于拓扑斯纤维积存在性的一般陈述, 参见 [16], 命题 4.47 (注意 Grothendieck 拓扑斯之间的一切几何态射都是有界的).}), 对每个拓扑斯图表
\[
\xymatrix@C=2cm@R=1.25cm{
E_1 \ar[d]&\\
E   & E_2 \ar[l] }
\]
在\pageoriginale 属于 $𝒱$ 的 $𝒰$-拓扑斯的 $2$-范畴 ($𝒱$-$𝒰$-$\Top$) 的意义下存在 $2$-纤维积 (在等价意义下, 与满足 $𝒰\in𝒱$ 的宇宙 $𝒱$ 的选取无关). 除 \hyperref[IV.prop5.11]{5.11} 之外的「拓扑斯纤维积」的自然例子, 见 \textsc{Giraud} 的著作 [3] 最后一章.
\end{enonce*}

\begin{enonce*}[remark]{练习 5.14}\etiquette[5.14]{IV.exer5.14}
\begin{enumerate}
\renewcommand{\theenumi}{\alph{enumi}}
\renewcommand{\labelenumi}{\theenumi)}
\item 设 $F$ 与 $G$ 为拓扑斯 $E$ 上的两个拓扑斯. 定义由对 $(f,\alpha)$ 组成的范畴 $\Hhomtop_E(F,G)$, 其中 $f:F\rightarrow G$ 为拓扑斯态射, $\alpha:p\rightarrow q\circ f$ 为从 $F$ 到 $E$ 的拓扑斯态射的同构 (其中 $p:F\rightarrow E$ 与 $q:G\rightarrow E$ 为结构拓扑斯态射). 定义配对函子 $\Hhomtop_E(F,G)\times\Hhomtop_E(G,H)\rightarrow\Hhomtop_E(F,H)$ 及结合同构.

\item 设 $X$ 与 $Y$ 为拓扑斯 $E$ 的两个对象. 定义一个从集合 $\Hom(X,Y)$ 定义的离散范畴到范畴 $\Hhomtop_E(E_{/X},E_{/Y})$ 的自然函子. 建立它与 $E$ 中态射复合以及 a) 中所考虑配对的相容性.

\item 证明 b) 中所考虑的函子为范畴等价. 特别地, 当「在 $E$-等价的意义下」知道了作为 \emph{$E$ 上拓扑斯}的诱导拓扑斯 $E_{/X}$ 时, 就可在唯一同构意义下重构拓扑斯 $E$ 的对象 $X$.~\footnote[1]{此结果归于 P. \textsc{Deligne}. $E$ 为点拓扑斯的情形先前已由 \textsc{Hakim} 女士处理.}
\end{enumerate}
\end{enonce*}

% END INCLUDED SGA4_English_current_buildable_source_20260906/source/04/drafts/SGA4_IV_LINES_2988_3412_English.texfrag


% BEGIN INCLUDED SGA4_English_current_buildable_source_20260906/source/04/drafts/SGA4_IV_LINES_3413_3631_English.texfrag
\section{拓扑斯的点与纤维函子}\etiquette[6]{IV.sec6}

\setcounter{subsection}{-1}
\subsection{}\etiquette[6.0]{IV.subsec6.0}
设 $P$ 为标准 $𝒰$-\emph{点拓扑斯} (\hyperref[IV.subsec2.2]{2.2})
$$
P=(𝒰\text{-}\Ens){\quoi}.
$$
符号\pageoriginale $P$ 代表具有点的性质的几何对象, 而 $(𝒰\text{-}\Ens)$ 则代表一个应解释为「$P$ 上的层范畴」的「大范畴」; 二者所联系的直观颇不相同, 因而应完全根据语境使用其中一个符号, 尽管从严格逻辑观点看, 它们表示同一个对象.

\begin{enonce*}{定义 6.1}\etiquette[6.1]{IV.def6.1}
设 $E$ 为 $𝒰$-拓扑斯. $E$ 的一个点是从点拓扑斯 $P$ (\hyperref[IV.subsec6.0]{6.0}) 到 $E$ 的拓扑斯态射 $P\rightarrow E$. 范畴 $\Hhomtop(P,E)$ (\hyperref[IV.subsec3.2]{3.2}) 称为 $E$ 的点范畴, 记为 $\Ppoints(E)$ 或 $\sheaf{\Pt}(E)$. 若 $p:P\rightarrow E$ 为 $E$ 的点, 相应逆像函子为 $p^{\ast}:E\rightarrow(𝒰\text{-}\Ens)$, 且 $F$ 为 $E$ 的对象, 则集合 $p^{\ast}(F)$ 称为 $F$ 在 $p$ 处的纤维, 记为 $F_p$.
\end{enonce*}

\setcounter{subsection}{1}
\subsubsection{}\etiquette[6.1.1]{IV.subsubsec6.1.1}
当然, 当 $F$ 为 $E$ 的群对象 (分别为环对象, 分别为\ldots) 时, 它在 $p$ 处的纤维 $F_p$ 为群 (分别为环, 分别为\ldots) (\hyperref[IV.subsubsec3.1.2]{3.1.2}).

\setcounter{subsubsection}{1}

\subsubsection{}\etiquette[6.1.2]{IV.subsubsec6.1.2}
由\hyperlink{sga4.exp.IV.subsubsection.3.2.1}{3.2.1}, $E$ 的点范畴通过函子 $p\mapsto p^{\ast}$ 等价于与 $𝒰$-归纳极限交换且左正合的函子
$$
\varphi:E\longrightarrow (𝒰\text{-}\Ens)
$$
的范畴的\emph{反范畴}. 因而给定 $E$ 的一个点, 实质上与给定这样的函子 $\varphi$ 是一回事.

\begin{enonce*}{定义 6.2}\etiquette[6.2]{IV.def6.2}
$𝒰$-拓扑斯 $E$ 的纤维函子是与 $𝒰$-归纳极限交换且左正合的函子 $\varphi:E\rightarrow(𝒰\text{-}\Ens)$. $E$ 上的纤维函子范畴记为 $\Ffib(E)$, 是\pageoriginale $\Hhom(E,(𝒰\text{-}\Ens))$ 中由 $E$ 上的纤维函子组成的满子范畴.
\end{enonce*}

\setcounter{subsection}{2}
\subsubsection{}\etiquette[6.2.1]{IV.subsubsec6.2.1}
由 \hyperref[IV.subsubsec6.1.2]{6.1.2}, $E$ 上的纤维函子范畴 $\Ffib(E)$ 因而通过从后者到前者的函子 $p\mapsto p^{\ast}$, 等价于 $E$ 的点范畴 $\Ppoint(E)$ 的\emph{反范畴}:
\begin{equation*}
\Ppoint(E)\xrightarrow{\approx}\Ffib(E)^{\circ}{\quoi}.
\tag{6.2.1.1}
\end{equation*}\etiquette[6.2.1.1]{IV.eq6.2.1.1}

函子
$$
\varphi:E\longrightarrow (𝒰\text{-}\Ens)
$$
为纤维函子, 当且仅当它是对应于 $E$ 的某个点 $p$ 的纤维函子 $F\mapsto F_p$ (\hyperref[IV.def6.1]{6.1}). 换言之, 函子 \eqhyperref[IV.eq6.2.1.1]{6.2.1.1} 为满射. 还注意到, 函子 $\varphi$ 为纤维函子当且仅当它左正合且把覆盖族映为满射族 (\hyperref[IV.cor1.7]{1.7}). 最后这个条件也可表述为 $\varphi$ 与 $𝒰$-和交换, 且把满态射映为满态射.

\subsection{}\etiquette[6.3]{IV.subsec6.3}
设 $C\in 𝒰$ 为景. \emph{景 $C$ 的点}是拓扑斯 $C^{\sim}$ 的点, 而\emph{景 $C$ 的点范畴}为范畴
$$
\Ppoint(C)=\Ppoint(C^{\sim}){\quoi}.
$$
另一方面, 考虑函子
\begin{equation*}
\varphi\longrightarrow \varphi|
C=\varphi\circ\epsilon_C:\Ffib(C^{\sim})\longrightarrow
\Hhom(C,(𝒰\text{-}\Ens)){\quoi},\tag{6.3.1}
\end{equation*}\etiquette[6.3.1]{IV.eq6.3.1}
它全忠实, 其本质像为满子范畴 $\Mmorsite((𝒰\text{-}\Ens),C)^{\circ}$ (\hyperref[IV.prop4.9.4]{4.9.4}), 也将记为 $\Ffib(C)$. 函子
$$
\Psi:C\longrightarrow (𝒰\text{-}\Ens)
$$
称\pageoriginale 为\emph{景 $C$ 上的纤维函子}, 如果它属于刚才描述的 \eqhyperref[IV.eq6.3.1]{6.3.1} 的本质像. 这样的函子因而\emph{连续} (从而更有 (III \hyperref[III.prop1.6]{1.6}), 它把覆盖族映为覆盖族, 即映为满射族), 并且左正合. 当有限投射极限在 $C$ 中可表时 (实践中遇到的情形几乎都满足此条件), 上述性质\emph{刻画}了 $C$ 上的纤维函子: 它们此时恰是左正合且把覆盖族映为满射族的函子 $\Psi:C\rightarrow(𝒰\text{-}\Ens)$ (\hyperref[IV.prop4.9.4]{4.9.4}).

应记住函子 $\varphi\rightarrow\varphi|C$ 为范畴等价
$$
\Ffib(C^{\sim})\xrightarrow{\approx}\Ffib(C){\quoi},
$$
因而\emph{给定拓扑斯 $C^{\sim}$ 上的纤维函子} (或等价地, 该拓扑斯的点), \emph{与给定景 $C$ 上的纤维函子, 实质上是一回事}. 给定 $C$ 上的纤维函子 $\Psi$, 它在同构意义下来自 $C^{\sim}$ 上的纤维函子 $\varphi$, 后者可由 $\Psi$ 用公式
\begin{equation*}
\varphi(F)\simeq \varinjlim_{C_{/F}}\Psi (X){\quoi},\tag{6.3.2}
\end{equation*}\etiquette[6.3.2]{IV.eq6.3.2}
在典范同构意义下重构, 其中 $C_{/F}$ 为 $C$ 的对象 $X$ 配有 (在 $\widehat{C}$ 中的) 态射 $X\rightarrow F$, 即 $F(X)$ 的元素, 所成的范畴. 公式 \eqhyperref[IV.eq6.3.2]{6.3.2} 直接来自 $\varphi$ 与归纳极限交换这一事实, 以及关于 $F$ 函子性的典范同构 (\hyperrefext[II][II.cor4.1.1]{4.1.1}):
$$
F\simeq \varinjlim_{C_{/F}}\epsilon_C(X){\quoi}.
$$
更\pageoriginale 一般地, 当 $G$ 为 $C$ 上的预层时, 有关于 $G$ 函子性的典范同构:
\begin{equation*}
\varphi(\underline{a}G)\simeq
\varinjlim_{C_{/G}}\Psi(X){\quoi},\tag{6.3.3}
\end{equation*}\etiquette[6.3.3]{IV.eq6.3.3}
其中 $\underline{a}G$ 为与 $G$ 相伴的层 (\hyperrefext[II][II.cor4.1.1]{4.1.1}). 事实上, 有公式
$$
\underline{a}G\simeq \varinjlim_{C_{/G}}\epsilon_C(X){\quoi}.
$$

\setcounter{subsection}{4} \setcounter{subsubsection}{-1}
\subsubsection{}\etiquette[6.4.0]{IV.subsubsec6.4.0}
关于\emph{保守函子族}
$$
\varphi_i:E\longrightarrow E_i\qquad i\in I{\quoi}.
$$
的概念, 参见 \hyperrefext[I][I.def6.1]{6.1}. 注意当 $E$ 与各 $E_i$ 为 $𝒰$-拓扑斯, 各函子 $\varphi_i$ 为拓扑斯态射
$$
f_i:E_i\longrightarrow E{\quoi},
$$
所对应的逆像函子 $f_i^{\ast}$ 时, \hyperrefext[I][I.prop6.2]{6.2}, \hyperrefext[I][I.cor6.3]{6.3} 及 \hyperrefext[I][I.prop6.4]{6.4} 中所假设的一切正合性质均满足, 只须在 \hyperrefext[I][I.prop6.2]{6.2} (v) 中假定 $D$ 有限. 此时以下条件等价: $(f_i^{\ast})_{i\in I}$ 保守; 它对\emph{单态射}保守 (\hyperrefext[I][I.def6.1]{6.1}); 它对\emph{满态射}保守; 它\emph{忠实}.

当\pageoriginale 函子族 $f_i^{\ast}$ 保守时, 有时也说\emph{拓扑斯态射族 $(f_i)_{i\in I}$ 保守}. 特别地:

\begin{enonce*}{定义 6.4.1}\etiquette[6.4.1]{IV.def6.4.1}
设 $E$ 为 $𝒰$-拓扑斯, $(p_i:P\rightarrow E)_{i\in I}$ 为 $E$ 的点族. 若相应的从 $E$ 到 $(𝒰\text{-}\Ens)$ 的纤维函子族 $F\mapsto F_{p_i}$ 保守 (\hyperref[IV.subsubsec6.4.0]{6.4.0}), 则称此点族保守. 若拓扑斯 $E$ 容许保守点族 (即该拓扑斯的一切点所成的族保守), 则称 $E$ 有足够多的点.
\end{enonce*}

\setcounter{subsection}{4} \setcounter{subsubsection}{1}
\subsubsection{}\etiquette[6.4.2]{IV.subsubsec6.4.2}
注意至此所用的一切 $𝒰$-拓扑斯都有足够多的点. 不过可以「特意」构造没有足够多的点的拓扑斯 (\hyperref[IV.exer7.2.6]{7.2.6} e) 及 \hyperref[IV.subsec7.4]{7.4}); 这样的拓扑斯必定不在 \hyperref[IV.subsec2.2]{2.2} 的几何意义下为「空的」. 有足够多的点的拓扑斯, 容许以某个 $I\in𝒰$ 为指标的保守点族 (\hyperref[IV.prop6.5]{6.5}). 不过注意, 若 $E$ 为 $𝒰$-拓扑斯, 则 $E$ 的点的同构类集不一定 $𝒰$-小 (\hyperref[IV.subsec7.3]{7.3}), 尽管实践中遇到的许多情形确实如此. 最后, 关于一个有意义的、提供足够多的点且涵盖代数几何中一切情形的存在定理 (归于 P. Deligne), \sisi{.}{参见 \hyperrefext[VI][VI.sec9]{9}.}

\subsubsection{}\etiquette[6.4.3]{IV.subsubsec6.4.3}
当 $(p_i)_{i\in I}$ 为拓扑斯 $E$ 的保守点族时, \hyperrefext[I][I.prop6.2]{6.2} 的讨论适用. 特别地, 这些讨论蕴涵: $E$ 中的两个箭头 $u,v:F\rightrightarrows G$ 相等当且仅当对每个 $i\in I$, 它们在纤维上诱导的箭头 $u_{p_i},v_{p_i}:F_{p_i}\rightrightarrows G_{p_i}$ 相\pageoriginale 等; $E$ 的箭头 $u:F\rightarrow G$ 为单态射 (分别为满态射) 当且仅当对每个 $i\in I$, 映射 $u_{p_i}:F_{p_i}\rightarrow G_{p_i}$ 如此; $E$ 的对象 $F$ 为始对象 (分别为终对象) 当且仅当对每个 $i\in I$, $F_{p_i}$ 为空 (分别只有一点); 对象 $F$ 为终对象 $e_E$ 的子对象当且仅当对每个 $i\in I$, $F_{p_i}$ 至多有一点.


% END INCLUDED SGA4_English_current_buildable_source_20260906/source/04/drafts/SGA4_IV_LINES_3413_3631_English.texfrag


% BEGIN INCLUDED SGA4_English_current_buildable_source_20260906/source/04/drafts/SGA4_IV_LINES_3632_3863_English.texfrag
\begin{enonce*}{命题 6.5}\etiquette[6.5]{IV.prop6.5}
\begin{enumerate}
\renewcommand{\theenumi}{\alph{enumi}}
\renewcommand{\labelenumi}{\theenumi)}
\item 设 $C$ 为 $𝒰$-景, $(\varphi_i)_{i\in I}$ 为 $C$ 上的纤维函子族 (\hyperref[IV.subsec6.3]{6.3}). 拓扑斯 $E=C^{\sim}$ 的相应点族 $(p_i)_{i\in I}$ 保守当且仅当: 对 $C$ 中的每个族 $(X_j\rightarrow X)_{j\in J}$, 如果对每个 $i\in I$, 相应的族 $(\varphi_i(X_j)\rightarrow\varphi_i(X))_{j\in J}$ 为满射族, 则该族为覆盖族.

\item 设 $E$ 为 $𝒰$-拓扑斯. 若 $E$ 有足够多的点 (\hyperref[IV.def6.4.1]{6.4.1}), 则 $E$ 容许一个 $𝒰$-小的保守点族.
\end{enumerate}
\end{enonce*}



断言 b) 为 \hyperrefext[I][I.prop7.7]{7.7} 的特例. 为证明 a), 对由 $\epsilon(X)$, $X\in\ob C$ 组成的 $E$ 的生成族应用 \hyperrefext[I][I.prop7.7]{7.7}. 回想 (\hyperrefext[II][II.theo4.4]{4.4}), 族 $X_j\rightarrow X$ 为覆盖族当且仅当族 $\epsilon(X_j)\rightarrow\epsilon(X)$ 为满态射族. 若点族 $(p_i)_{i\in I}$ 保守, 这就相当于 (\hyperref[IV.subsubsec6.4.3]{6.4.3}) 对每个 $i\in I$, 族 $(X_j)_{p_i}\rightarrow(X)_{p_i}$ 为满射族, 即族 $\varphi_i(X_j)\rightarrow\varphi_i(X)$ 为满射族. 这证明了 a) 的必要性. 对充分性, 应用判别准则 \hyperrefext[I][I.prop7.7]{7.7}, 将问题化为验证: 每个使 $F_{p_i}\rightarrow\epsilon(X)_{p_i}$ 对每个 $i\in I$ 均为同构的单态射 $F\rightarrow\epsilon(X)$\pageoriginale 都是同构, 或等价地, 都为满态射. 现在可以找到态射的覆盖族 $\epsilon(X_j)\rightarrow F$, 再进一步加细, 可假设态射 $\epsilon(X_j)\rightarrow\epsilon(X)$ 由态射 $X_j\rightarrow X$ 诱导. 按假设, 对每个 $i\in I$, 复合态射 $(X_j)_{p_i}\rightarrow F_{p_i}\rightarrow(X)_{p_i}$ 构成满射族 (因为是两个满射族的复合), 即族 $\varphi_i(X_j)\rightarrow\varphi_i(X)$ 为满射族. 由假设可知族 $X_j\to X$ 为覆盖族, 从而族 $\epsilon(X_j)\rightarrow\epsilon(X)$ 为满态射族, 故更有 $F\rightarrow\epsilon(X)$ 为满态射.

\begin{enonce*}{推论 6.5.1}\etiquette[6.5.1]{IV.cor6.5.1}
设 $C$ 为范畴. 使 $C$ 成为有足够多纤维函子的 $𝒰$-景 (即相应拓扑斯有足够多的点) 的 $C$ 上的拓扑, 完全由 $\Hhom(C,\Ens)$ 中由 $C$ 上的纤维函子组成的满子范畴 $\Phi$ 确定.
\end{enonce*}

事实上, 由 \hyperref[IV.prop6.5]{6.5} a), 此时可用纤维函子族描述覆盖族.

事实上, 将在下面 \eqhyperref[IV.subsubsec6.8.5]{6.8.5} 看到, 子范畴 $\Phi$ 包含在 $C$ 上的 pro-可表函子范畴中, 从而 $\Phi^{\circ}\simeq\Ppoint(C^{\sim})$ 在等价意义下等同于 $\Pro-C$ 的满子范畴.

将 \hyperref[IV.cor6.5.1]{6.5.1} 与 \textsc{Giraud} 定理 \hyperrefext[II][II.theo5.5]{5.5} 比较, 后者在 $C$ 上的拓扑集与 $\Hhom(C^{\circ},\Ens)$ 的某个满子范畴集之间建立了一一对应.

\begin{enonce*}[remark]{练习 6.5.2}\etiquette[6.5.2]{IV.exer6.5.2}
设\pageoriginale $C$ 为等价于 $𝒰$ 中某范畴的范畴, $P$ 为 $\Pro(C)$ 的子范畴. 对每个 $p\in P$, 用 $X\mapsto X_p$ 表示由 $p$ pro-表示的函子. $C$ 中的族 $(X_i\rightarrow X)_{i\in I}$ 称为 $P$-\emph{覆盖族}, 如果对每个 $p\in P$, 族 $(X_{ip}\rightarrow X_p)_{i\in I}$ 为满射族. 证明 $C$ 上有拓扑 $T_P$, 其覆盖族恰为 $P$-覆盖族, 且 $T_P$ 是使函子 $X\mapsto X_p$, $p\in P$, 成为纤维函子的 $C$ 上最细拓扑. 证明拓扑 $T_P$ 使 $C$ 成为有\emph{足够多}纤维函子的景. 证明对 $C$ 上的每个拓扑 $T$, 在比 $T$ 细且有足够多纤维函子的拓扑 $T'$ 中有一个最粗者: 它为拓扑 $T_P$, 其中 $P$ 为下面 \eqhyperref[IV.subsubsec6.8.5]{6.8.5} 中描述的, 等价于 $\Ppoint(C)$ 的 $\Pro(C)$ 的严格满子范畴.
\end{enonce*}

\begin{enonce*}[remark]{问题 6.5.3}\etiquette[6.5.3]{IV.prop6.5.3}
刻画 $\Pro(C)$ 的那些可由 $C$ 上某个拓扑作为 $\Ppoint(C^{\sim})$ 的本质像出现的子范畴 (严格满、对 $𝒰$-滤过投射极限稳定, \ldots) $P$ (\hyperref[IV.subsubsec6.8.5]{6.8.5}).\footnote[1]{必要条件参见 \hyperref[IV.exer9.1.8]{9.1.8} a) 及 e); 对 $\Top(X)$ 且 $X$ 为素的局部 Noether 空间的情形, 充分必要条件参见 \hyperref[IV.exer9.1.12]{9.1.12} b).}
\end{enonce*}

\setcounter{subsection}{5}
\subsection{}\etiquette[6.6]{IV.subsec6.6}
设
$$
f:E\longrightarrow E'
$$
为 $𝒰$-拓扑斯态射. 它诱导典范函子
$$
\Ppoint(f)=(p\longmapsto f\circ p):\Ppoint(E)\longrightarrow
\Point(E'){\quoi}.
$$
给定两个可复合拓扑斯态射
$$
E\xrightarrow{f} E'\xrightarrow{g} E''{\quoi},
$$
立即\pageoriginale 可验证
$$
\Ppoint(gf)=\Ppoint(g)\circ\Ppoint(f):\Ppoint(E)\longrightarrow
\Ppoint(E')\longrightarrow \Ppoint(E''){\quoi}.
$$
这明确了 $\Ppoint(E)$ 对拓扑斯 $E$ 的函子性依赖 (这次终于完全没有语言滥用): 若 $𝒱$ 为满足 $𝒰\in𝒱$ 的宇宙, 则得到一个从作为 $𝒱$ 的元素的 $𝒰$-拓扑斯范畴到作为 $𝒱$ 的元素的范畴的范畴的 (真正的!) 函子
$$
\Ppoint:(𝒱\text{-}𝒰\text{-}\text{top})\longrightarrow
(𝒱\text{-}\cat)
$$
.

\subsection{}\etiquette[6.7]{IV.subsec6.7}
设 $E$ 为 $𝒰$-拓扑斯. 则 $E$ 的每个对象 $F$ 定义反变函子
$$
p\longmapsto F_p=p^{*}(F):\Ppoint(E)^{\circ}\longrightarrow
(𝒰\text{-}\Ens){\quoi},
$$
因而当 $F$ 变化时给出典范函子
\begin{equation*}
E\longrightarrow
\Ppoint(E)^{\widehat{\phantom{,}}}=\Hhom(\Ppoint(E)^{\circ},
(𝒰\text{-}\Ens)){\quoi},\tag{6.7.1}
\end{equation*}\etiquette[6.7.1]{IV.eq6.7.1}
按定义 (\hyperref[IV.prop6.5]{6.5}), 此函子忠实当且仅当 $E$ 有足够多的点. (函子 \eqhyperref[IV.eq6.7.1]{6.7.1} 与 Fourier 变换具有某种形式上的类比\ldots)

设 $X$ 为 $E$ 的对象, 因而定义了
$$
C=\Ppoint(E){\quoi}.
$$
上的预层 $\widehat{X}$. 因而可定义 $C$ 上的预层范畴 $\widehat{C}$ 中的态射 $p\rightarrow\widehat{X}$ 所成的范畴
$$
C_{/\widehat{X}}=\Ppoint(E)_{/\widehat{X}}
$$
即由对 $(p,\xi)$ 所成的范畴, 其中 $p$ 为 $E$ 的点, $\xi$ 为 $X$ 在 $p$ 处的纤维 $\widehat{X}(p)=X_p$ 的元素. 据此, 我断言\emph{存在典范范畴等价}
\begin{equation*}
C'=\Ppoint(E_{/X})\xrightarrow{\approx}C_{/\widehat{X}}\quav \mbox{
其中 }\quad C=\Ppoint(E){\quoi},\tag{6.7.2}
\end{equation*}\etiquette[6.7.2]{IV.eq6.7.2}
它与包含函子 $C_{/\widehat{X}}\rightarrow C$ 的\pageoriginale 复合, 正是由局部化态射 \eqhyperref[IV.eq5.2.1]{5.2.1}
$$
j_X:E_{/X}\longrightarrow E{\quoi}.
$$
函子性地 (\hyperref[IV.subsec6.6]{6.6}) 诱导的函子
$$
\Ppoint(j_X):C'=\Ppoint(E_{/X})\longrightarrow C=\Ppoint(E)
$$
这就是 \hyperref[IV.prop5.12]{5.12} 中取 $E'=P=(\Ens)$ 所得的特例.

\begin{enonce*}{推论 6.7.3}\etiquette[6.7.3]{IV.cor6.7.3}
设 $E$ 为拓扑斯. 若 $E$ 有足够多的点, 则每个诱导拓扑斯 $E_{/X}$, $X\in\ob E$, 也有足够多的点. 反之, 若 $(X_i)_{i\in I}$ 为 $E$ 中覆盖终对象的对象族, 且对每个 $i\in I$, $E_{/X_i}$ 有足够多的点, 则 $E$ 有足够多的点.
\end{enonce*}

对第一个断言, 设 $f:X'\rightarrow X''$ 为 $E_{/X}$ 中的非同构态射. 我们证明存在 $E_{/X}$ 的点 $q$, 使得 $f_q$ 不是同构. 因为 $E$ 有足够多的点, 存在 $E$ 的点 $p$ 使得 $f_p:X'_p\rightarrow X''_p$ 不是同构. 这蕴涵存在 $q\in X_p$, 使得 $f_p$ 在 $X'_p$ 与 $X''_p$ 于 $X_p$ 上的 $q$ 处的纤维 $X'_q$ 与 $X''_q$ 上诱导的映射不是同构. 但通过等价 \eqhyperref[IV.eq6.7.2]{6.7.2}, 可将 $q$ 等同于 $E_{/X}$ 的点, 而刚考虑的映射 $X'_q\rightarrow X''_q$ 正是 $f:X'\rightarrow X''$ 在 $q$ 处的纤维上诱导的映射, 这就证明了断言.

反之, 假设各 $X_i$ 覆盖 $E$ 的终对象, 且各 $E_{/X_i}$ 有足够多的点. 我们证明 $E$ 也如此. 设\pageoriginale $f:X'\rightarrow X''$ 为 $E$ 中的非同构态射. 则存在某个 $i\in I$ 使得 $f_i:X'_i\rightarrow X''_i$ 不是同构, 其中下标 $i$ 表示限制到 $X_i$. 因而存在 $E_{/X_i}$ 的点 $p_i$, 使得 $f_{ip_i}:X'_{ip_i}\rightarrow X''_{ip_i}$ 不是同构. 若 $p$ 表示 $p_i$ 通过局部化态射 $E_{/X_i}\rightarrow E$ 在 $E$ 中的像, 则这就是说 $X'_p\rightarrow X''_p$ 不是同构, 从而证明 $E$ 有足够多的点.

\begin{enonce*}[remark]{注 6.7.4}\etiquette[6.7.4]{IV.rem6.7.4}
上述论证表明, 若 $(p_{\alpha})$ 为 $E$ 的保守点族, 则位于某个 $p_{\alpha}$ 上方的 $E_{/X}$ 的点所成的族 (以各 $X_{p_{\alpha}}$ 的不交并为指标的族) 保守. 同样, 若各 $X_i$ 覆盖 $E$ 的终对象, 且对每个 $i\in I$, $P_i$ 为 $E_{/X_i}$ 的保守点集, 则 $E$ 中由 $p\in P_i$ 通过局部化态射 $E_{/X_i}\rightarrow E$ 所得的\emph{像}组成的点集 (其中 $i\in I$ 且 $p$ 变化) 保守.
\end{enonce*}

\setcounter{subsection}{7}

% END INCLUDED SGA4_English_current_buildable_source_20260906/source/04/drafts/SGA4_IV_LINES_3632_3863_English.texfrag


% BEGIN INCLUDED SGA4_English_current_buildable_source_20260906/source/04/drafts/SGA4_IV_LINES_3864_4120_English.texfrag
\subsection{}\etiquette[6.8]{IV.subsec6.8}
设 $E$ 为拓扑斯, $p$ 为 $E$ 的点. \emph{拓扑斯 $E$ 的点 $p$ 的邻域}是一对 $(X,u)$, 其中 $X\in \ob E$ 且 $u\in X_p$. 由 \eqhyperref[IV.eq6.7.2]{6.7.2}, 也可将 $u$ 解释为 $p$ 到诱导拓扑斯 $E_{/X}$ 的某个点的提升. 用 $\varphi_p$ 表示与点 $p$ 相联系的纤维函子, 也可将点 $p$ 的邻域解释为范畴 $E_{/\varphi_p}$ 的对象 (这是 $(E^\circ)^{\widehat{\phantom{,}}}_{/\varphi_p}$ 中由源属于 $E\subset \widehat{E}$ 的对象组成的满子范畴). 自然称后一个范畴, 即由上述对 $(X,u)$ 所成的范畴, 为\emph{拓扑斯 $E$ 的点 $p$ 的邻域范畴}; 还记为 $\Vvois(p)$. 因为 $\varphi_p$\pageoriginale 左正合, 且有限投射极限在 $E$ 中可表, 所以有限投射极限在范畴 $\Vvois(p)$ 中可表, 且典范函子 $\Vvois(p)\rightarrow E$ 与它们交换 (\hyperrefext[I][I.sec8.10.10]{8.10.10}). 因而更有反范畴 $\Vvois(p)^{\circ}$ 为滤过的. 此外, 显然 (\hyperrefext[I][I.prop3.4]{3.4}) 有关于 $F\in \ob E$ 函子性的典范同构
\begin{equation*}
\varphi_p(F)=F_p\simeq
\varinjlim_{\Vvois(p)^{\circ}}F(X){\quoi}.\tag{6.8.1}
\end{equation*}\etiquette[6.8.1]{IV.eq6.8.1}
换言之, \emph{函子 $\varphi_p$, 即「在点 $p$ 处的纤维」, 同构于拓扑斯 $E$ 中由 $E$ 的点 $p$ 的各邻域所表的函子的滤过归纳极限}. 事实上, \emph{纤维函子 pro-可表} (\hyperrefext[I][I.sec8]{8}), 这也表示它由 $E$ 的 pro-对象 $(X_i)_{i\in I}$ 所表, 其中 $I$ 为满足 $I\in𝒰$ 的滤过有序集. 事实上, 由该处, 这等价于说范畴 $\Vvois(p)$ 容许小满共尾子范畴. 现在设 $C$ 为 $E$ 的小满生成子范畴; 我断言, $\Vvois(p)=E_{/\varphi_p}$ 中由满足 $X\in \ob C$ 的 $p$ 的邻域 $(X,u)$ 组成的满子范畴 $C_{/\varphi_{p}}$ (此范畴显然为小的), 在 $\Vvois(p)$ 中共尾. 事实上, 设 $(X,u)$ 为 $p$ 的邻域; 须找到位于 $(X,u)$ 上方的 $p$ 的邻域 $(X',u')$, 使得 $X'\in \ob C$. 存在覆盖族 $X'_i\rightarrow X$, 其中各 $X'_i$ 属于 $C$, 从而各 $X'_{ip}$ 覆盖 $X_p$; 故存在 $X'=X'_{i_{\circ}}$ 及 $u'\in X'_p$, 使得 $(X',u')$ 位于 $(X,u)$ 上方, 如所断言的那样.

\setcounter{subsubsection}{1}
\subsubsection{}\etiquette[6.8.2]{IV.subsubsec6.8.2}
设\pageoriginale $C$ 为 $𝒰$-景, $p$ 为 $C$ 的点, 即拓扑斯 $C^{\sim}$ 的点. 滥用记号, 仍用 $\varphi_p$ 表示纤维函子 $\varphi_p$ 到 $C$ 的「限制」, 即复合 $\varphi_p\circ \epsilon$; 对 $C$ 的对象 $X$, 还常写
$$
X_p=\varphi_p(X)=\epsilon(X)_p{\quoi}.
$$
\emph{景 $C$ 中点 $p$ 的邻域}是一对 $(X,u)$, 其中 $X\in \ob C$ 且 $u\in X_p=\varphi_p(X)$. 这些邻域仍构成范畴 $C_{/\varphi_p}$, 可记为 $\Vvois_C(p)$. 来证明\emph{范畴 $\Vvois_C(p)^{\circ}$ 仍滤过, 容许一个 $𝒰$-小的满共尾子范畴, 且有关于 $C$ 上的层 $F$ 函子性的同构:}
\begin{equation*}
F_p\simeq \varinjlim_{\Vvois_C(p)^{\circ}} F(X){\quoi}.\tag{6.8.3}
\end{equation*}\etiquette[6.8.3]{IV.eq6.8.3}
先注意, 重复 \eqhyperref[IV.eq6.8.1]{6.8.1} 的论证, 若 $C'$ 为 $C$ 的拓扑生成满子范畴 (II \hyperref[II.def3.0.1]{3.0.1}), 则 $\Vvois_{C'}(p)^{\circ}$ 在 $\Vvois_C(p)^{\circ}$ 中共尾. 取 $C'\in𝒰$ 为小的, 就把断言的证明化归到 $C\in𝒰$ 的情形. 此时 $\widehat{C}$ 为 $𝒰$-拓扑斯, 且有相伴层函子 $\underline{a}:\widehat{C}\rightarrow C^{\sim}$, 从而可在 $\widehat{C}$ 上构造复合纤维函子 $\varphi_p\underline{a}$. 对拓扑斯 $\widehat{C}$ 及其满生成子范畴 $C$ 应用 \eqhyperref[IV.eq6.8.1]{6.8.1}, 可知范畴 $C_{/\varphi_p\underline{a}}^{\circ},$ 显然同构于 $\Vvois_C(p)^{\circ}$, 是滤过的, 并且有关于 $C$ 上的预层 $P$ 函子性的同构:
$$
\varphi_p\underline{a}(P)\simeq \varinjlim_X P(X){\quoi}.
$$
若 $F$ 为 $C$ 上的层, 则 $F=\underline{a}F$, 对看作预层的 $F$ 应用上述同构,\pageoriginale 即得 \eqhyperref[IV.eq6.8.3]{6.8.3}. 同时, 这个论证建立了更一般的公式
\begin{equation*}
(\underline{a}P)_p\simeq \varinjlim_{\Vvois_C(p)^{\circ}}
P(X){\quoi},\tag{6.8.4}
\end{equation*}\etiquette[6.8.4]{IV.eq6.8.4}
此同构关于 $C$ 上的任意预层 $P$ 为函子性的, 至少当 $C$ 为小时如此. 对一般情形, 引入满足 $C\in𝒱$ 的宇宙 $𝒱$, 把 $C$ 上的 $\varphi_p$ 看作取值于 $(𝒱\text{-}\Ens)$ 的纤维函子, 并用相伴层的构造与扩大宇宙的相容性 (\hyperrefext[II][II.prop3.6]{3.6}), 即得结论.

\setcounter{subsubsection}{4}
\subsubsection{}\etiquette[6.8.5]{IV.subsubsec6.8.5}
我们将拓扑斯 $C^{\sim}$ 的每个点 $p$ 配以 $𝒰$-景 $C$ 的一个 pro-对象, 由典范函子 $\Vvois_C(p)\rightarrow C$ 定义, 这里考虑到了 $p$ 在 $C$ 中的邻域范畴为余滤过的, 且容许小满共尾子范畴这一事实. 立即可验证, 当 $p$ 变化时, 这给出函子
\begin{equation*}
\Ppoint(C^{\sim})\longrightarrow \Pro(C){\quoi}.\tag{6.8.5.1}
\end{equation*}\etiquette[6.8.5.1]{IV.eq6.8.5.1}
此函子\emph{全忠实}. 通过相应的纤维函子, 可将它解释为函子 $\Fib(C)\rightarrow \Fib(C^{\sim})\rightarrow \Pro(C)^{\circ}$. 由 \hyperref[IV.prop4.9.4]{4.9.4} (在 \hyperref[IV.subsec6.3]{6.3} 中对纤维函子作过回顾), 第一个函子为等价, 第二个由 pro-对象的一般形式理论 (\hyperrefext[I][I.sec8]{8}) 为全忠实函子.

\subsubsection{}\etiquette[6.8.6]{IV.subsubsec6.8.6}
考虑 $C$ 为 $𝒰$-小且赋有混沌拓扑, 因而 $C^{\sim}=\widehat{C}$ 的特殊情形. 我断言此时上述函子甚至为范畴等价
\begin{equation*}
\Ppoint(\widehat{C})\xrightarrow{\approx}\Pro(C){\quoi}.\tag{6.8.6.1}
\end{equation*}\etiquette[6.8.6.1]{IV.eq6.8.6.1}
事实上,\pageoriginale 只须证明此函子本质满射, 这由形如 $P\mapsto P(X)$ 的 $\widehat{C}$ 上的函子为纤维函子这一事实, 以及拓扑斯上的纤维函子的每个滤过归纳极限仍是纤维函子这一显然事实得出.

将 $C$ 等同于 $\Pro(C)$ 的满子范畴 (\hyperrefext[I][I.sec8]{8}), 由 \sisi{(6.8.5.2)}{\eqhyperref[IV.eq6.8.6.1]{6.8.6.1}} 的一个拟逆诱导的函子
$$
C\longrightarrow \Ppoint(\widehat{C})
$$
显然同构于 \eqhyperref[IV.eq4.6.2.2]{4.6.2.2} 中已经考虑的典范函子; 考虑到 $\Ppoint(\widehat{C})$ 对小投射极限稳定, 所说拟逆可以看作它的典范扩充 (\hyperrefext[I][I.sec8]{8}).

\subsubsection{}\etiquette[6.8.7]{IV.subsubsec6.8.7}
更一般地, 设 $(X_i)_{i\in I}$ 为 $𝒰$-景 $C$ 的 pro-对象, 从而给出 $C^{\sim}$ 上的函子
\begin{equation*}
\varphi(F)=\varinjlim_I F(X_i){\quoi},\tag{6.8.7.1}
\end{equation*}\etiquette[6.8.7.1]{IV.eq6.8.7.1}
由 (\hyperref[IV.subsubsec6.8.5]{6.8.5}), 自然要问此函子何时为纤维函子. 以下条件等价:
\begin{enumerate}
\renewcommand{\theenumi}{\roman{enumi}}
\renewcommand{\labelenumi}{(\theenumi)}
\item 上述函子 $\varphi$ 为 $C^{\sim}$ 上的纤维函子.

\item $\varphi$ 到 $C$ 的「限制」为景 $C$ 上的纤维函子 (\hyperref[IV.subsec6.3]{6.3}).

\item 对 $C$ 的对象 $Y$ 的每个覆盖族 $Y_{\alpha}\rightarrow Y$, 每个 $i_{\circ}\in I$ 以及每个态射 $X_{i_{\circ}}\rightarrow Y$, 存在 $i\geq i_{\circ}$, $\alpha$ 及态射 $X_i\rightarrow Y_{\alpha}$, 使图表
\[
\xymatrix@C=2cm@R=1.25cm{
X_i \ar[d] \ar[r]& Y_{\alpha} \ar[d]\\
X_{{i}_{o}} \ar[r]     &  Y \quad . }
\]
交换.

事实上,\pageoriginale 条件 (iii) 仅表示 $\varphi| C$ 将覆盖族映为覆盖族. 蕴涵 (i) $\Rightarrow$ (ii) $\Rightarrow$ (iii) 显然, 只须证明 (iii) $\Rightarrow$ (i). 因为函子 $\varphi$ 显然左正合, 只须证明 (\hyperref[IV.prop4.9.4]{4.9.4}) 它把 $C^{\sim}$ 中的覆盖族 $F_{\alpha}\rightarrow F$ 映为覆盖族, 即对每个 $i_{\circ}$ 和每个 $x_{\circ}\in F(X_{i_{\circ}})$, 存在 $i\geq i_{\circ}$, $\alpha$ 及 $y\in F_{\alpha}(X_i)$, 使得 $x_{\circ}$ 与 $y$ 在 $F(X_i)$ 中的像相同. 因为族 $F_{\alpha}\rightarrow F$ 为覆盖族, 存在 $X_{i_{\circ}}$ 的覆盖族 $Z_{\gamma}\rightarrow X_{i_{\circ}}$, 使对每个 $\gamma$, $x_{\circ}$ 在 $F(Z_{\gamma})$ 中的逆像提升为某个 $F_{\alpha}(Z_{\gamma})$ 的元素. 由假设 (iii), 存在 $i\geq i_{\circ}$ 及 $X_{i_{\circ}}$-态射 $X_i\rightarrow Z_{\gamma}$. 因而 $x_{\circ}$ 在 $F(X_i)$ 中的像提升为某个 $F_{\alpha}(X_i)$ 的元素 $y$, \cqfd
\end{enumerate}

\begin{enonce*}[remark]{练习 6.9}\etiquette[6.9]{IV.exer6.9}
设 $C\in𝒰$ 为景, $\varphi$ 为 $C$ 上的纤维函子, 由 $\Pro(C)$ 的对象 $(X_i)_{i\in I}$ pro-表示. 对每个指标 $i\in I$, $C$ 的每个对象 $Y$, 每个态射 $u:X_i\rightarrow Y$ 以及 $Y$ 上的每个覆盖筛 $R$, 选取指标 $j\geq i$ 使得复合 $X_j\rightarrow X_i\xrightarrow{u}Y$ 经过 $R$ 分解. 对固定的 $i$, 设 $J(i)$ 为由 $i$ 以及 $Y$, $u$, $R$ 变化时所得到的各 $j$ 组成的集合; 对 $I$ 的子集 $I'$, 类似地设 $J(I')$ 为 $i\in I'$ 的各 $J(i)$ 的并. 另一方面, 若 $I'$ 为 $I$ 的有限子集, 设 $m(I')\in I$ 为 $I'$ 的一个上界; 对 $I$ 的任意子集 $I'$, 设 $M(I')$ 为一切 $m(I'')$ 的集合, 其中 $I''$ 遍历\pageoriginale $I'$ 的有限子集集. 假定函数 $I'\mapsto m(I')$ 选为当 $I'$ 只有一个元素 $i$ 时, 有 $m(I')=i$; 这蕴涵对 $I$ 的每个子集 $I'$, 有 $I'\subset M(I')$. 考虑从 $I$ 的子集集到自身的映射 $MJ$ 的迭代 $(MJ)^n$, 并对 $I$ 的每个子集 $I'$, 设 $P(I')$ 为各 $(MJ)^n(I')$ 的并.
\begin{enumerate}
\renewcommand{\theenumi}{\alph{enumi}}
\renewcommand{\labelenumi}{\theenumi)}
\item 证明对 $I$ 的每个子集 $I'$, $P(I')$ 按诱导序为 $I$ 的滤过子集; 且当 $I'$ 遍历 $I$ 的\emph{有限}子集集时, $I$ 是各 $P(I')$ 的递增滤过并.

\item 证明存在基数 $c\in 𝒰$, 只依赖于基数 $\card \text{fl}(C)=a$, 使得对 $I$ 的每个有限子集 $I'$, 有 $\card(P(I'))\leq c$. (若 $a$ 无限, 可取 $c=2^{a}$.)

\item 用 \hyperref[IV.subsubsec6.8.7]{6.8.7} 证明, 对 $I$ 的每个子集 $I'$, 由 pro-对象 $(X_i)_{i\in P(I')}$ pro-表示的 $C$ 上的函子 $\varphi_{I'}$ 为 $C$ 上的纤维函子, 且纤维函子 $\varphi$ 是当 $I'$ 遍历 $I$ 的有限子集集时各函子 $\varphi_{I'}$ 的滤过归纳极限.

\item 推出存在 $\Ffib(C)\approxeq \Ffib(C^{\sim})$ 的小满子范畴 $\Phi$, 使 $\Ffib(C)$ 的每个对象都是 $\Phi$ 的对象的滤过归纳极限. (若 $a$ 无限, 可取 $\Phi$ 的基数为 $\leq 2^{2^{a}}$; 若 $a$ 有限, 当然可取 $\Phi=C$.)
\end{enumerate}
\end{enonce*}

\begin{enonce*}[remark]{练习 6.10}\etiquette[6.10]{IV.exer6.10}
设 $C$ 为 $\mathcal{U}$-景, $D$ 为小滤过 $\varinjlim$ 可表的 $\sheaf{\mathcal{U}}$-范畴; 设 $\Fais(C,D)$ 为 $C$ 上取值于 $D$ 的层范畴 (\hyperrefext[II][II.def6.1]{6.1}). 扩充 \eqhyperref[IV.eq6.8.3]{6.8.3}, 定义「纤维函子」
\begin{equation*}
F\longrightarrow F_p:\Fais(C,D)\longrightarrow D{\quoi}.\tag{6.10.1}
\end{equation*}\etiquette[6.10.1]{IV.eq6.10.1}
若有限 $\varprojlim$ 在 $D$ 中可表, 且与滤过 $\varinjlim$ 交换, 则上述函子左正合.
\end{enonce*}

% END INCLUDED SGA4_English_current_buildable_source_20260906/source/04/drafts/SGA4_IV_LINES_3864_4120_English.texfrag


% BEGIN INCLUDED SGA4_English_current_buildable_source_20260906/source/04/drafts/SGA4_IV_LINES_4121_4475_English.texfrag
\section{纤维函子与拓扑斯的点的例子}\etiquette[7]{IV.sec7}\pageoriginale

\subsection{拓扑空间 $X$ 的 $\Top(X)$ 的情形}\etiquette[7.1]{IV.subsec7.1}

设 $x$ 为 $X$ 的点. 将单元素集 $\{x\}$ 看作拓扑空间, 并考虑包含
$$
i_x:\{x\}\longrightarrow X{\quoi}.
$$
由 \hyperref[IV.subsubsec4.1.1]{4.1.1}, 它定义拓扑斯态射 (在唯一同构意义下定义)
\begin{equation*}
\Top(i_x):\Top(\{x\})\longrightarrow \Top(X){\quoi}.\tag{7.1.1}
\end{equation*}\etiquette[7.1.1]{IV.eq7.1.1}
另一方面, 可将 $P$ 等同于 $\Top(\{x\})$ (\hyperref[IV.subsec2.2]{2.2}), 从而得到 $\Top(X)$ 的点 $p_x$:
\begin{equation*}
p_x:P\longrightarrow \Top(X){\quoi},\tag{7.1.2}
\end{equation*}\etiquette[7.1.2]{IV.eq7.1.2}
由 $x$ 在唯一同构意义下定义. 相应纤维函子由熟悉的公式 $[TF]$ 给出, 而该公式还可直接由 \hyperrefext[I][I.prop5.1]{5.1} 得出:
\begin{equation*}
p^{\ast}_x(F)=F_x=F_{p_x}=\varinjlim_{U\ni x}F(U){\quoi},\tag{7.1.3}
\end{equation*}\etiquette[7.1.3]{IV.eq7.1.3}
其中极限遍历 $x$ 在 $X$ 中的开邻域 $U$.

因为已知 $\Top(X)$ 同构于 $\Top(X_{\sob})$ (\hyperref[IV.subsubsec4.2.1]{4.2.1}), 更一般地可知 $X_{\sob}$ 的每个点, 即 $X$ 的每个不可约闭子集 $Z$, 都定义 $\Top(X)$ 的点, 或等价地, 定义 $\Top(X)$ 上的纤维函子, 仍记为 $F\mapsto F_Z$. 回到它在 $X_{\sob}$ 上通过公式 \eqhyperref[IV.eq7.1.3]{7.1.3} 的定义, 得到
\begin{equation*}
F_Z=\varinjlim_{U\cap Z\neq \phi} F(U){\quoi},\tag{7.1.4}
\end{equation*}\etiquette[7.1.4]{IV.eq7.1.4}
其中\pageoriginale 归纳极限遍历 $X$ 的与 $Z$ 相交的开子集 $U$.

\setcounter{subsubsection}{4}
\subsubsection{}\etiquette[7.1.5]{IV.subsubsec7.1.5}
众所周知, 纤维函子 $F\mapsto F_x$ $(x\in X)$ 已经构成保守族. 将 $X$ 上的层解释为平展空间时, 这一点特别清楚, 因为此时 $F$ 在 $x$ 处的纤维就是作为 $X$ 上的纤维化空间的意义下的纤维. 注意所考虑的保守纤维函子族的指标集为 $X$, 因而为 $𝒰$-小的.

\subsubsection{}\etiquette[7.1.6]{IV.subsubsec7.1.6}
由 \hyperref[IV.subsubsec4.2.3]{4.2.3}, 范畴 $\Ppoint(\Top(X))$ 等价于按特殊化关系排序的集合 $X_{\sob}$ 所对应的范畴. 换言之, $\Top(X)$ 上的每个纤维函子都同构于一个纤维函子 $F\mapsto F_Z$, 其中 $Z$ 为唯一确定的 $X_{\sob}$ 的元素, 即唯一确定的 $X$ 的不可约闭子集; 此外, 若 $Z$, $Z'\in X_{\sob}$, 则集合 $\Hom(F_Z,F_{Z'})$ 或者为空, 或者只有一点, 后一情形当且仅当 $Z$ 在 $X_{\sob}$ 中为 $Z'$ 的特殊化, 即当且仅当作为 $X$ 的子集有 $Z\subset Z'$. 当 $X$ 本身为素空间时, 可在这些陈述中以 $X$ 自身代替 $X_{\sob}$.

注意, $\Top(X)$ 的纤维函子的自同构群 (与 $\Top(X)$ 的相应点的自同构群互为反群) 总是平凡群 (更一般地, \hyperref[IV.subsubsec4.2.3]{4.2.3} 对与拓扑空间相联系的每个拓扑斯态射的自同构群给出同样结论). 这是所考虑情形中特别特殊的现象: 见 \hyperref[IV.subsec7.2]{7.2} 的例子, 以及\pageoriginale VIII 7. 对与「模问题」相联系的拓扑斯 (例如参见 [10]), 纤维函子的自同构群还有一个显著解释, 即所要分类的 (代数闭域上的) 代数结构的自同构群.

\subsubsection{}\etiquette[7.1.7]{IV.subsubsec7.1.7}
我们刚看到如何从拓扑斯 $\Top(X)$ (甚至只须在等价意义下知道它), 将素空间 $X$ (或与任意拓扑空间相联系的素空间 $X_{\sob}$) 重构为 $\Top(X)$ 的「点」的同构类集, 至少作为一个按特殊化关系排序的集合. 由 \hyperref[IV.subsec2.1]{2.1} 所述, 还能如下重构 $X$ 的拓扑, 即它的开子集族: 对 $E$ 的终对象 $e_E$ 的每个子对象 $U$, 设 $\Pt(U)$ 为满足 $U_x\neq\phi$ 的 $x\in X$ 所成的集合 (由 \eqhyperref[IV.eq6.7.2]{6.7.2}, 它还等同于诱导拓扑斯 $E_{/U}$ 的点的同构类集). 则 $U\mapsto \Pt(U)$ 是从 $e_E$ 的子对象集到 $X$ 的开子集集的有序集同构.

\subsubsection{}\etiquette[7.1.8]{IV.subsubsec7.1.8}
\hyperref[IV.subsubsec7.1.6]{7.1.6} 对拓扑空间 $X$ 所定义的拓扑斯 $\Top(X)$ 的点范畴的确定, 导致关于任意拓扑斯 $E$ 的点 $p$, $p'$ 的以下术语: 当存在从 $p'$ 到 $p$ 的态射 (在 \hyperref[IV.def6.1]{6.1} 的范畴 $\Ppoint(E)$ 的意义下) 时, 说 $p$ 为 $p'$ 的\emph{特殊化}, 或 $p'$ 为 $p$ 的\emph{一般化}. 这些关系仍是传递的, 但用 \hyperref[IV.subsubsec6.8.6]{6.8.6}, 很容易找到 $p$ 与 $p'$ 互为对方的特殊化\pageoriginale, 而 $p$ 与 $p'$ 不同构 (更不相等) 的例子. 拓扑斯 $E$ 的点 $p$ 的一般化范畴, 即我们所指的范畴 $\Ppoint(E)_{/p}$, 在某些方面起着类似于对概形在一点处作局部化的作用. 这个作用将在 \sisi{第 VI 章}{\hyperrefext[VI][VI.subsec8.4]{8.4}} 中, 通过构造 $E$ 在点 $p$ 处局部化所得到的拓扑斯而明确, 它的点范畴典范等价于 $p$ 的一般化范畴.

\begin{enonce*}[remark]{练习 7.1.9}\etiquette[7.1.9]{IV.exer7.1.9}
设 $E$ 为 $𝒰$-拓扑斯. 证明 $E$ 等价于形如 $\Top(X)$ 的拓扑斯, 其中 $X$ 为 $\in𝒰$ 的拓扑空间, 当且仅当它满足以下两个条件:
\begin{enumerate}
\renewcommand{\theenumi}{\alph{enumi}}
\renewcommand{\labelenumi}{\theenumi)}
\item 终对象 $e_E$ 的子对象族为拓扑斯 $E$ 的生成族.

\item $E$ 有足够多的点 (\hyperref[IV.prop6.5]{6.5}). (此条件并非多余: 参见 \hyperref[IV.subsec7.4]{7.4}.)
\end{enumerate}

与 \hyperref[IV.exer7.8]{7.8} c) 比较.
\end{enonce*}

\begin{enonce*}[remark]{练习 7.1.10}\etiquette[7.1.10]{IV.exer7.1.10}
设 $X$ 为赋有算子群 $G$ 的拓扑空间 $(X,G\in𝒰)$, 并设 $E=\Top(X,G)$ (\hyperref[IV.subsec2.3]{2.3}).
\begin{enumerate}
\renewcommand{\theenumi}{\alph{enumi}}
\renewcommand{\labelenumi}{\theenumi)}
\item 证明对 $E$ 的每个对象 $(X',G)$, 诱导拓扑斯 $E_{/(X',G)}$ 典范等价于拓扑斯 $\Top(X',G)$ (比较 \hyperref[IV.subsec5.7]{5.7}).

\item 当 $G$ 在 $X'$ 上的作用为适当且自由的时 (Bourbaki, Gen. Top. 第 \hyperrefext[III][III.sec4]{4} 章), 证明带算子空间的态射 $(X',G)\rightarrow (X'/G,e)$ 诱导 (\hyperref[IV.subsubsec4.1.2]{4.1.2}) 拓扑斯等价
$$
\Top(X',G)\longrightarrow \Top(X'/G){\quoi}.
$$

\item 由 b) 推出\pageoriginale 存在 $E$ 的对象 $Z$, 使诱导拓扑斯 $E_{/Z}$ 等价于 $\Top(X)$, 而局部化函子 $E\rightarrow E_{/Z}$ 同构于「遗忘 $G$ 作用」的函子. (仿照 \hyperref[IV.subsubsec5.8.3]{5.8.3} 的论证.)

\item 由 c) 推出, $E$ 上的每个纤维函子都通过「遗忘 $G$ 作用」的函子, 由 $\Top(X)$ 上的某个纤维函子诱导, 从而可由 $X_{\sob}$ 的某个点定义.

\item 根据带算子空间 $(X_{\sob},G)$ 确定范畴 $\Ppoint(E)$ 的结构 (在等价意义下). 特别地推出, $X_{\sob}$ 的两点定义 $E$ 上同构的纤维函子当且仅当它们在 $G$ 的作用下共轭.
\end{enumerate}
\end{enonce*}

\begin{enonce*}[remark]{练习 7.1.11}\etiquette[7.1.11]{IV.exer7.1.11}
设 $X\in 𝒰$ 为拓扑空间. 证明 $𝒱$-拓扑斯 $\TOP(X)$ (\hyperref[IV.subsec2.5]{2.5}) 有足够多的点. 更确切地, 对 $\TOP(X)$ 的每个对象 $X'$ 及每个 $x'\in X'$, 在 $\TOP(X)$ 上定义一个只依赖于空间 $(X',x')$ 在 $X$ 上的「芽」的纤维函子 $F\mapsto F_{X',x'}$, 并证明此纤维函子族保守 (且以 $𝒱$-小的指标集为指标). 用一个适当的滤过投射系 $(X'_i,x'_i)_{i\in I}$, 其中 $I$ 不是 $𝒰$-小的, 证明这种方式并不能得到 $𝒱$-拓扑斯 $\TOP(X)$ 的全部纤维函子. 证明此类纤维函子的同构类集的势不 $\in 𝒱$.
\end{enonce*}

\subsection{分类拓扑斯 $B_G$ 的点}\etiquette[7.2]{IV.subsec7.2}\pageoriginale

设 $E$ 为 $𝒰$-拓扑斯, $G$ 为 $E$ 的群, 从而有分类拓扑斯 $B_G$ (\hyperref[IV.subsec2.4]{2.4}). 若 $e$ 为单位群, 则态射 $e\rightarrow G\rightarrow e$ 定义 (\hyperref[IV.subsec4.5]{4.5}) 拓扑斯态射 (考虑到 $B_e\approx E$)
\begin{equation*}
E\longrightarrow B_G \longrightarrow E{\quoi},\tag{7.2.1}
\end{equation*}\etiquette[7.2.1]{IV.eq7.2.1}
其复合同构于 $\id_E$; 转到点范畴就得到函子
\begin{equation*}
\Ppoint(E)\longrightarrow \Ppoint(B_G)\longrightarrow
\Ppoint(E){\quoi},\tag{7.2.2}
\end{equation*}\etiquette[7.2.2]{IV.eq7.2.2}
其复合同构于恒等函子. 用 \eqhyperref[IV.eq7.2.1]{7.2.1} 的第一个拓扑斯态射等同于相对于 $B_G$ 的对象 $E_G=G$ 的局部化态射这一事实 (\hyperref[IV.subsec5.8]{5.8}); 由 \eqhyperref[IV.eq6.7.2]{6.7.2} 可知 \eqhyperref[IV.eq7.2.2]{7.2.2} 的第一个函子等同于自然「包含」函子
\begin{equation*}
\Ppoint(B_G)_{/\widehat{E}_G}\longrightarrow
\Ppoint(B_G){\quoi},\tag{7.2.3}
\end{equation*}\etiquette[7.2.3]{IV.eq7.2.3}
其中 $\widehat{E}_G$ 表示 $\Ppoint(B_G)$ 上的预层 $p'\mapsto (E_G)_{p'}$. 因为 $E_G$ 显然覆盖 $B_G$ 的终对象, 它的纤维 $(E_G)_p$ 非空, 从而 \eqhyperref[IV.eq7.2.3]{7.2.3} 本质满射; 这表示 \emph{$B_G$ 上的每个纤维函子都 (在同构意义下) 经由「遗忘 $G$ 作用」的函子 $B_G\rightarrow E$, 由 $E$ 上的纤维函子诱导}. 特别地得到, \emph{若 $E$ 有足够多的点 (分别有小保守点族)}, 则 $B_G$ 也如此.

\begin{enonce*}[remark]{注 7.2.4}\etiquette[7.2.4]{IV.rem7.2.4}
可以进一步根据范畴 $C=\Ppoint(E)$ 及其上的群预层
$$
\widehat{G}:p\longmapsto G_p: C^{\circ}\longrightarrow
(\text{群})
$$
确定范畴 $C'=\Ppoint(B_G)$ 的结构 (在等价意义下).\pageoriginale 事实上, 定义一个新范畴 $C_1$, 具有与 $C$ 相同的对象, 且对 $C$ 的两个对象 $p$, $q$, 有
$$
\Hom_{C_1}(p,q)=\Hom_C(p,q)\times \widehat{G}(p){\quoi},
$$
$C_1$ 中箭头的复合由 $C$ 中的复合以及预层 $\widehat{G}$ 的群律诱导. 此时给定到任意范畴 $D$ 的函子 $\varphi_1:C_1\rightarrow D$, 等价于给定一个配有预层 $\widehat{G}$ 作用的函子 $\varphi:C\rightarrow D$ (即对每个 $p\in \ob C$, 给定 $\widehat{G}(p)$ 在 $\varphi(p)$ 上的作用, 并在 $p$ 变化时满足一个显然的函子性条件). 利用这个观察, 定义典范函子
\begin{equation*}
\varphi_1:C_1\longrightarrow C'{\quoi},\tag{7.2.4.1}
\end{equation*}\etiquette[7.2.4.1]{IV.eq7.2.4.1}
它对应于 \eqhyperref[IV.eq7.2.2]{7.2.2} 的函子 $\varphi:C\rightarrow C'$, 即将 $E$ 的每个点 $p$ 配以 $B_G$ 上诱导的点 $\varphi(p)$, 以及后者的自然 $\widehat{G}(p)=G_p$ 作用; 此作用来自 $X$ 遍历 $B_G$ 时, $G_p$ 在各 $X_p$ 上的作用. 留给读者用 \eqhyperref[IV.eq7.2.2]{7.2.2} 的第一个函子的结构 \eqhyperref[IV.eq7.2.3]{7.2.3}, 并注意自然包含函子 $C\rightarrow C_1$ 相对于 $C'$ 上的预层 $P'=\widehat{E}_G$ 的逆像 $P_1$ 具有类似结构, 来证明 \eqhyperref[IV.eq7.2.4.1]{7.2.4.1} 为范畴等价.
\end{enonce*}

\setcounter{subsubsection}{4}
\subsubsection{}\etiquette[7.2.5]{IV.subsubsec7.2.5}
特别地\pageoriginale 得到, \emph{$B_G$ 的点的同构类恰与 $E$ 的点的同构类对应: 分类拓扑斯 $B_G$ 的每个点都由 $E$ 的点在一个}不唯一的\emph{同构意义下诱导}. 特别地, 当 $E$ 为点拓扑斯 $P$, 因而 $G$ 为通常的群时, 范畴 $\Ppoint(B_G)$ 是基本群为 $G$ 的连通群胚: $B_G$ 上的每个纤维函子都 (非典范地) 同构于「遗忘 $G$ 作用」的函子 $\omega$, 而后一个函子的自同态幺半群为群 $G$ (从而 $\omega$ 的每个自同态都是自同构).

\begin{enonce*}[remark]{练习 7.2.6}\etiquette[7.2.6]{IV.exer7.2.6}
\begin{enumerate}
\renewcommand{\theenumi}{\alph{enumi}}
\renewcommand{\labelenumi}{\theenumi)}
\item 设 (Tors) 为由对 $(\Gamma,P)\in 𝒰$ 组成的范畴, 其中 $\Gamma$ 为群, $P$ 为 $\Gamma$ 下的右挠子. 函子 $(\Gamma,P)\mapsto \Gamma$
\begin{equation*}
(\Tors)\longrightarrow (\Groupes)\tag{7.2.6.1}
\end{equation*}\etiquette[7.2.6.1]{IV.eq7.2.6.1}
为余纤维化函子. 用 \hyperref[IV.rem7.2.4]{7.2.4} 的记号, 考虑函子
$$
\widehat{G}:C^{\circ}\longrightarrow (\Groupes){\quoi},
$$
并设 $D'$ 为从余纤维化范畴 \eqhyperref[IV.eq7.2.6.1]{7.2.6.1} 作逆像得到的 $C^{\circ}$ 上的余纤维化范畴, $D$ 为相应的 $C$ 上的纤维化范畴, 与 $D'$ 有相同的纤维范畴, 因而对 $p\in \ob C$ 有
$$
D_p=\mbox{ $G_p$ 下的右挠子范畴}{\quoi}.
$$
构造典范函子
\begin{equation*}
D\longrightarrow C'=\Ppoint(B_G)\quav C'\longrightarrow
D{\quoi},\tag{7.2.6.2}
\end{equation*}\etiquette[7.2.6.2]{IV.eq7.2.6.2}
使它们\emph{互为拟逆}. 特别地推出, $B_G$ 在 $E$ 的一个给定点 $p$ 上方的点的范畴, 典范等价于群 $G_p$ 下的右挠子范畴.

\item 设\pageoriginale $𝒢=(G_i)_{i\in I}$ 为 $𝒰$-拓扑斯 $E$ 中群的严格投射系 ($I\in𝒰$ 为滤过预序集), 从而仿照 \hyperref[IV.subsubsec2.7.1]{2.7.1} 的定义得到分类拓扑斯 $B_{𝒢}$. 仿照 a) 确定 $B_{𝒢}$ 的点范畴. (注意. 将范畴 $(\Tors)$ 与 $(\Groupes)$ 换成以 $I$ 为指标的这些范畴中的投射系的范畴.) 推出若 $I$ 容许可数共尾子集, 则 $B_{𝒢}$ 上的每个纤维函子都同构于由 $E$ 上的纤维函子 (通过「遗忘 $G$ 作用」的函子) 所诱导的函子, 后者在不唯一的同构意义下确定.

\item 取 $E$ 为点拓扑斯, 从而 $𝒢$ 为通常的 pro-群, 证明即使 $I$ 任意, 只要 $𝒢$ 为 pro-有限的 (即各 $G_i$ 有限), b) 的结论仍成立: 每个纤维函子都同构于「遗忘 $G$ 作用」的函子.

\item 仍取 $E$ 为点拓扑斯的情形, 对一个适当的投射系 $𝒢=(G_i)_{i\in I}$, 给出一个 $B_{𝒢}$ 上的纤维函子的例子, 它不同构于「遗忘 $G$ 作用」的函子.

\item 将 $𝒢=(G_i)$ 看作拓扑斯 $\widehat{I}$ 中的群, 并设 $𝒯$ 为 $I$ 上\emph{带}为 $G$ 的\emph{团} (gerbe) [\hyperref[IV.sec3]{3}]. 说明如何借助团 $T$「扭曲」分类拓扑斯 $B_G$, 得到拓扑斯 $B^T_{𝒢}$, 它是非典范等价于分类拓扑斯 $B_{G_i}$ 的各子范畴 $B^T(i)$ 的归纳极限. 通过在 $B^T_{𝒢}$ 上的纤维函子范畴与 $T$ 在 $I$ 上的截面范畴之间建立\pageoriginale 范畴等价, 证明拓扑斯 $B^T_{𝒢}$ 容许纤维函子当且仅当团 $T$ 为「中性的」. 推出, 若各 $G_i$ 交换, 因而 $I$ 上带为 $𝒢$ 的团由 $\displaystyle{\varprojlim_I}^{(2)}G_i=H^{2}(I,𝒢)$ 分类 (见前引文献), 则存在一个形如 $B^T_{𝒢}$ 的 (非「空」(\hyperref[IV.subsec2.2]{2.2})) 拓扑斯的例子, 它没有点. (取一个 $\displaystyle{\varprojlim_I}^{(2)}G_i\neq 0$ 的例子.)
\end{enumerate}
\end{enonce*}

% END INCLUDED SGA4_English_current_buildable_source_20260906/source/04/drafts/SGA4_IV_LINES_4121_4475_English.texfrag


% BEGIN INCLUDED SGA4_English_current_buildable_source_20260906/source/04/drafts/SGA4_IV_LINES_4476_4830_English.texfrag
\subsection{预层拓扑斯 $\widehat{C}$ 的点; 点范畴不等价于小范畴的 $𝒰$-预层拓扑斯 $\widehat{C}$ 的例子}\etiquette[7.3]{IV.subsec7.3}

设 $E$ 为 $𝒰$-拓扑斯, $(\varphi_i)_{i\in I}$ 为 $E$ 上的滤过纤维函子族. 由函子 $\varinjlim$ 的正合性质 (\sisi{I\quad 及}{\hyperrefext[I][I.prop2.8]{2.8}}), 函子 $\varphi=\varinjlim\varphi_i$ 也为纤维函子: \emph{纤维函子的每个滤过归纳极限都是纤维函子}. 因而范畴 $\Ffib(E)$ 容许 $𝒰$-滤过归纳极限 (且包含函子 $\Ffib(E)\rightarrow \Hhom(E,(𝒰\text{-}\Ens))$ 与它们交换); 换言之, 范畴 $\Ppoint(E)$ 容许 $𝒰$-滤过投射极限. 练习 \sisi{\hyperref[IV.exer7.1.10]{7.1.10}}{\hyperref[IV.exer7.1.11]{7.1.11}} 已用这个事实给出「大」纤维函子范畴的例子.

\hypertarget{sga4.exp.IV.subsubsection.7.3.3}{}用 \eqhyperref[IV.eq6.8.6.1]{6.8.6.1}, 可以通过形如 $E=\widehat{C}$, $C\in𝒰$ 的预层拓扑斯构造简单得多的例子. 这立即给出 $\Ffib(\widehat{C})$ 不等价于 $\in𝒰$ 的范畴的例子, 即这个 $𝒰$-范畴的对象的同构类集的基数不 $\in𝒰$ 的例子. 这表示 pro-对象范畴 $\Pro-(C)$ 不等价于 $\in 𝒰$ 的范畴. 例如, 只须取 $D=C^{\circ}$ 为形如 $[0,n]$ 的有限集的范畴, 其中 $n\geq 0$ 为\pageoriginale 整数; 此时 $\Ind(D)\simeq \Pro(C)^{\circ}$ 等价于非空集范畴. 拓扑斯 $E$ 此时为熟悉的\emph{余单纯集}范畴. 也可以取 $C$ 为非空有限集范畴; 此时发现单纯集的预层拓扑斯 $\widehat{C}$ 的点范畴等价于 pro-有限集范畴, 或等价地, 等价于全不连通紧空间范畴.

\subsection{没有点的非空拓扑斯}\etiquette[7.4]{IV.subsec7.4}

以下例子归于 $P$. Deligne. (另一个例子参见 \hyperref[IV.exer7.2.6]{7.2.6} e).) 设 $K$ 为配有测度 $\mu$ 的紧空间, $U$ 为 $K$ 的可测子集模去零测集所成的有序集. 使 $U$ 成为范畴 $\underline{U}$, 其中 $\ob \underline{U}=U$, $\underline{U}$ 的态射为 $U$ 的元素之间的「包含态射」. 取如下预拓扑使 $\underline{U}$ 成为景: 对 $E\in U$, $\Cov(E)$ 由以 $E$ 为上界、且各 $E_i$ 的并模去某个零测集为 $E$ 的 $U$ 的元素的\emph{可数}族组成. 这给出拓扑斯 $\Top(\mu)=\underline{U}^{\sim}$, 其终对象的子对象集为生成族 (这个拓扑斯看来未曾受到概率学家的注意). 此拓扑斯为「空拓扑斯」(\hyperref[IV.subsec2.2]{2.2}) 当且仅当 $\mu=0$. 另一方面, 此拓扑斯的点范畴等价于满足 $\mu(\{x\})\neq 0$ 的点 $x\in K$ 所成的集定义的离散范畴 (证明留给读者). 因而当 $K$ 没有这样的点时, 它为空, 例如 $K$ 为实直线中的单位区间, 赋予 Lebesgue 测度诱导的测度.

\begin{enonce*}[remark]{练习 7.5}\etiquette[7.5]{IV.exer7.5}
(\emph{\sisi{Karoubi}{Karoubi} 范畴\pageoriginale 与拓扑斯态射})\footnote[1]{比较 \hyperrefext[I][I.exer8.7.8]{8.7.8}.}
\begin{enumerate}
\renewcommand{\theenumi}{\alph{enumi}}
\renewcommand{\labelenumi}{\theenumi)}
\item 设 $C$ 为范畴, $X$ 为 $C$ 的对象, $p$ 为 $X$ 的自同态. 若 $p^2=p$, 则称 $p$ 为\emph{投影子}. 证明若 $p$ 为投影子, 则 $\Ker(\id_X,p)$ 可表当且仅当 $\Coker(\id_X,p)$ 可表, 且如此得到的 $C$ 的两个对象典范同构. 此时说投影子 $p$ \emph{有像}, 并称 $\Ker(\id_X,p)\simeq \Coker(\id_X,p)$ 为\emph{投影子 $p$ 的像}; 依语境而定, 将它等同于 $C$ 的子对象或商对象. 与 $X$ 中某个适当投影子 $p$ 的 $\text{Im\,} p$ 同构的对象, 称为\emph{$C$ 的对象 $X$ 的直因子}. 如果 $C$ 的对象中的每个投影子都有像, 就称 $C$ 为\emph{有直因子的范畴}, 或 \emph{\sisi{Karoubi}{Karoubi} 范畴}. 若 $F:C\rightarrow C'$ 为与核或余核交换的函子, 则 $F$ 把有像的投影子映为有像的投影子\nde{这里对 $F$ 不需加条件: 每个函子都保持有像的投影子.}.

\item 证明对每个范畴 $C$, 都能找到从 $C$ 到一个 Karoubi 范畴 (在等价意义下确定) 的函子 $\varphi:C\rightarrow \kar (C)$, 使对每个 Karoubi 范畴 $C'$, 函子
$$
f\mapsto f\circ \varphi:\Hhom(\kar(C),C')\longrightarrow \Hhom(C,C')
$$
为范畴等价; $\kar(C)$ 称为范畴 $C$ 的 \emph{Karoubi 包络}. (提示: 取 $\ob \kar (C)$ 为对 $(X,p)$ 的集合, 其中 $X\in \ob C$ 且 $p$ 为 $X$ 中的投影子; 取 $\Hom((X,p),(Y,q))$ 为 $\Hom(X,Y)$ 中由满足 $f=qfp.$ 的那些 $f:X\rightarrow Y$ 所成的子集.) 证明 $\varphi$ 全忠实.

\item 设 $F:E\rightarrow E'$ 为范畴之间的函子. 证明若 $F$ 与某种给定归纳或投射极限交换,\pageoriginale 则 $F$ 的每个直因子也如此. 特别地, 若 $E$ 与 $E'$ 为 $𝒰$-拓扑斯, 且 $F$ 为拓扑斯态射 $E'\rightarrow E$ 的逆像函子, 则 $F$ 的每个直因子也如此; 因而范畴 $\Hhomtop(E',E)$ 为 Karoubi 范畴, 特别地范畴 $\Ppoint(E)$ 为 Karoubi 范畴.

\item 证明对每个 $𝒰$-范畴 $C$, 范畴 $\Ind(C)$ 为 Karoubi 范畴 (其中 $\Ind(C)$ 使用以滤过预序集 $I\in 𝒰$ 为指标的 $C$ 中的归纳系构成). (若 $C\in 𝒰$, 例如可用 $\Ind(C)$ 等价于范畴 $\Ffib((C^{\circ})^{\widehat{\phantom{,}}})=\Ppoint((C^{\circ})^{\widehat{\phantom{,}}})^{\circ}.$ 这一事实 (\hyperref[IV.subsec7.3]{7.3}).) 推出 $\kar(C)$ 的另一种构造: 它为 $\Ind(C)$ 中由 $C$ 的对象 $X$ 的投影子在 $\Ind(C)$ 中的像组成的满子范畴.
\end{enumerate}
\end{enonce*}

\begin{enonce*}[remark]{练习 7.6}\etiquette[7.6]{IV.exer7.6}
(\emph{拓扑斯的本质态射, 本质点}).
\begin{enumerate}
\renewcommand{\theenumi}{\alph{enumi}}
\renewcommand{\labelenumi}{\theenumi)}
\item 设 $f:E\rightarrow F$ 为拓扑斯态射. 证明以下条件等价:
\begin{itemize}
\item[i)] $f_{!}$ 存在, 即 $f^{\ast}$ 有左伴随.

\item[ii)] $f^{\ast}$ 与 $𝒰$-投射极限交换.

\item[${\rm ii')}$] $f^{\ast}$ 与 $𝒰$-积交换.
\end{itemize}
此时称 $f$ 为从拓扑斯 $E$ 到拓扑斯 $F$ 的\emph{本质态射}.

证明若 $f$ 满足这些条件, 则 $f$ 的每个直因子也如此. (用 \hyperref[IV.rem1.8]{1.8} 及 \hyperref[IV.exer7.5]{7.5} c).)

\item 设 $E$ 为拓扑斯. \emph{拓扑斯 $E$ 的本质点}是 $E$ 中使 $p_{!}$ 存在的点 $p:P\rightarrow E$. 证明若 $E=\Top(X)$, 其中 $X$ 为拓扑空间, 且 $p$ 为 $E$ 中由 $x\in X$ 定义的点, 则 $p$ 为本质点\pageoriginale 当且仅当 $x$ 有一个最小开邻域 $U$ (当 $X$ 的点都闭时, 这表示 $x$ 为 $X$ 的\emph{孤立点}). 拓扑斯态射 $f:E\rightarrow F$ 为本质态射 (a)) 的必要条件是 $f$ 将本质点映为本质点; 当 $E$ 有足够多的本质点 (即这些点所成的族保守) 时, 此条件也充分 (参见 d)).

\item 拓扑斯 $E$ 的点 $p$ 为本质点, 当且仅当相应的纤维函子由 $E$ 的对象 $X$ 所表. 由 $E$ 的给定对象 $X$ 所表的协变函子 $\varphi:E\rightarrow (\Ens)$ 为纤维函子 (即定义 $E$ 的点, 且必为本质点) 的充分必要条件是 $X$ 非空、连通 (\hyperref[IV.subsubsec4.3.5]{4.3.5}) 且投射 (即函子 $\varphi$ 把满态射映为满态射). (提示: 先证明 $\varphi$ 与余积交换当且仅当 $X$ 非空且连通, 再用判别准则 \hyperref[IV.prop4.9.4]{4.9.4}.) 推出范畴 $\Ppointess(E)$ 与 $E$ 中由非空连通投射对象组成的满子范畴 $P$ 之间的等价.

\item 证明 $E$ 的拓扑在 $P$ 上诱导的拓扑为混沌拓扑. 推出 $E$ 的以下条件等价 (归于 J. E. ROOS [12 c), 命题 1]): (i) $E$ 的本质点族保守; (ii) $E$ 中由非空连通投射对象组成的满子范畴 $P$ 为生成子范畴; (iii)~$E$ 等价于形如 $\widehat{C}$ 的预层拓扑斯, 其中 $C$ 为等价于 $\in 𝒰$ 的范畴的范畴 (或者若愿意, $C\in 𝒰$). (用 $c$) 以及以下事实: 若 $P$ 为生成子范畴, 则自然函子 $E\rightarrow P^{\sim}$ 为范畴\pageoriginale 等价.)

\item 拓扑斯 $E$ 的本质点范畴为 Karoubi 范畴 (参见 \hyperref[IV.exer7.5]{7.5} c)). 设 $C$ 为等价于 $\in 𝒰$ 的范畴的范畴. 证明典范全忠实函子 (\hyperref[IV.subsubsec4.6.2]{4.6.2})
$$
C\longrightarrow \Ppoint(\widehat{C})
$$
经过函子
$$
C\longrightarrow \Ppointess(\widehat{C}){\quoi},
$$
分解, 后者使 $\Ppointess(\widehat{C})$ 成为 $C$ 的 \emph{Karoubi 包络} (\hyperref[IV.exer7.5]{7.5} b)). 特别地, 此函子为范畴等价当且仅当范畴 $C$ 为 Karoubi 范畴. (提示: 用 $c$), 证明 $\widehat{C}$ 的每个本质点都同构于某个 $X\in \ob C$ 的直因子.)

\item 设 $C$ 与 $C'$ 为两个等价于 $\in 𝒰$ 的范畴的范畴. 证明典范全忠实函子 (\hyperref[IV.subsubsec4.6.2]{4.6.2})
$$
\Hhom(C,C')\longrightarrow \Hhomtop(\widehat{C},\widehat{C}')
$$
取值于本质\emph{拓扑斯}态射的满子范畴 $\Hhomtopess(\widehat{C},\widehat{C}')$ (a)), 且诱导函子
$$
\Hhom(C,C')\longrightarrow \Hhomtopess(\widehat{C},\widehat{C}')
$$
为范畴等价当且仅当 $C$ 为空或 $C'$ 为 Karoubi 范畴. (用下面的 $g$).)

\item 设 $C$ 为等价于 $\in 𝒰$ 的范畴的范畴, $E$ 为拓扑斯. 在本质拓扑斯态射 $\widehat{C}\rightarrow E$ 的范畴 $\Hhomtopess(\widehat{C},E)$ 与范畴 $\Hhom(C, \Ppointess(E))$ 之间定义等价. (用 $b$)\pageoriginale 以及 $\widehat{C}$ 有足够多的本质点这一事实.)

\item 设 $C$ 为有限范畴. 证明预层拓扑斯 $\widehat{C}$ 的每个点都为本质点, 且 $\Ppoint(\widehat{C})$ 等价于有限范畴. (注意 $C$ 的 Karoubi 包络有限, 并用 \eqhyperref[IV.eq6.8.6.1]{6.8.6.1}, 将问题化为证明: 若 $C$ 为有限 Karoubi 范畴, 则典范函子 $C\rightarrow \Pro(C)$ 为范畴等价.) 用 b) 推出, \emph{每个}拓扑斯态射 $\widehat{C}'\rightarrow \widehat{C}$ 都是本质态射, 且 $\Hhom(C,C')\rightarrow \Hhomtop(\widehat{C},\widehat{C}')$ 为等价当且仅当 $C$ 为空或 $C'$ 为 Karoubi 范畴.

\item 设 $X$ 为拓扑空间, $f:\Top(X)\rightarrow P$ 为从 $X$ 到单点空间的连续映射所诱导的拓扑斯态射. 证明 $f$ 为本质态射当且仅当空间 $X$ \emph{局部连通}, 即满足以下条件: 对每个 $x\in X$ 以及 $x$ 的每个开邻域 $U$, $x$ 在 $U$ 中的连通分支是 $x$ 的邻域; 或等价地, 对 $X$ 的每个开子集 $U$, $U$ 的连通分支都是开的. (推广见 \hyperref[IV.exer8.7]{8.7} l).)
\end{enumerate}
\end{enonce*}

\begin{enonce*}[remark]{练习 7.7}\etiquette[7.7]{IV.exer7.7}
(\emph{分类拓扑斯的不寻常的点})\footnote[1]{也参见 \hyperref[IV.exer7.2.6]{7.2.6} d).}.
\begin{enumerate}
\renewcommand{\theenumi}{\alph{enumi}}
\renewcommand{\labelenumi}{\theenumi)}
\item 设 $G$ 为幺半群. 幺半群 $G$ 含有满足 $g\neq e$ 及 $g^2=g$ 的元素 $g$, 当且仅当分类拓扑斯 $B_G$ 容许一个不同构于平凡点 (对应于纤维函子「遗忘 $G$ 作用」) 的本质点. (用 \hyperref[IV.exer7.6]{7.6} e).)

\item 设 $G$ 为整数 $\geq 0$ 的加法幺半群. 证明分类拓扑斯 $B_G$ 的每个本质点都同构于平凡点. 构造 $B_G$ 的一个不同构于平凡点的点. 证明范畴 $\Ppoint(B_G)$ 不等价于 $\in 𝒰$ 的范畴.
\end{enumerate}
\end{enonce*}

\begin{enonce*}[remark]{练习 7.8}\etiquette[7.8]{IV.exer7.8}
(\emph{$\Point(E)$ 上的拓扑, 以及与有序集相联系的拓扑斯}).
\begin{enumerate}
\renewcommand{\theenumi}{\alph{enumi}}
\renewcommand{\labelenumi}{\theenumi)}
\item 设 $E$ 为拓扑斯. 用 $\Point(E)$ 表示 $E$ 的点的同构类集. 对 $E$ 的每个开对象 $U$, 设 $U'$ 为\pageoriginale $E$ 中满足 $U_p\neq \phi$ 的点 $p$ 的类所成的集合. 证明 $U\mapsto U'$ 为从 $E$ 的开对象有序集 $\Ouv(E)$ 到 $X=\Point(E)$ 的子集集的保序映射, 且此映射与有限下确界及任意 $\Sup$ 交换. 推出 $X$ 的形如 $U'$ 的子集定义 $X$ 上的拓扑 (称为\emph{拓扑斯 $E$ 的点的类集上的典范拓扑}). 证明若 $E$ 有足够多的点, 则映射 $U\mapsto U'$ 是从有序集 $\Ouv(E)$ 到有序集 $\Ouv(X)$ 的同构.

\item 设 $\Oo\in 𝒰$ 为容许有限下确界和任意 $\Sup$ 的有序集, 从而定义一个具有有限积及任意余积的范畴 $\cat(\Oo)$. 说有共同靶 $U$ 的态射族 $U_i\rightarrow U$ 为覆盖族, 如果 $U$ 为各 $U_i$ 的 $\Sup$. 假设 $\cat(\Oo)$ 中的余积是泛的, 即 $\Oo$ 中的任意 $\Sup$ 与下确界交换. 这定义了 $\cat(\Oo)$ 上的拓扑, 使它成为景, 因而给出拓扑斯 $\cat(\Oo)^{\sim}$, 也简记为 $\Oo^{\sim}$. 证明 $\Oo$ 典范同构于 $\Ouv(\Oo^{\sim})$.

证明范畴 $\Hhomtop(E,\Oo^{\sim})$ 等价于与有限下确界及任意 $\Sup$ 交换的有序集态射 $\Oo\rightarrow \Ouv(E)$ 所成的有序集对应的范畴. (用判别准则 \hyperref[IV.prop4.9.4]{4.9.4}.) 若 $\Oo'$ 为满足与 $\Oo$ 同样条件的有序集, 推出 $\Hhomtop(\Oo^{\sim}, {\Oo'}^{\sim})$ 等价于与有限下确界及任意 $\Sup$ 交换的一切有序集态射 ${\Oo}'\rightarrow\Oo$ 所成的有序集对应的范畴.

\item 证明\pageoriginale 拓扑斯 $E$ 等价于形如 $\Oo^{\sim}$ 的拓扑斯, 其中 $\Oo$ 如 b) 中所述, 当且仅当 $E$ 的开对象族在 $E$ 中为生成族. 它等价于某个 $\Top(X)$, 其中 $X\in 𝒰$, 当且仅当它满足前述条件且有足够多的点. 假设第一个条件成立, 用 b) 的最后一个断言, 根据有序集 $\Oo=\Ouv(E)$ 的结构表述第二个条件.

\item 定义一个典范拓扑斯态射
$$
E\longrightarrow \Ouv(E)^{\sim}{\quoi},
$$
使它对从 $E$ 到由其开对象生成的拓扑斯的态射 (参见 c)) 为泛的 (在等价意义下).

\item 设 $X'$ 为 $X=\Point(E)$ 的\emph{小}子集, 赋予 $X$ 的拓扑诱导的拓扑. 定义典范拓扑斯态射
$$
\Ouv(E)^{\sim}\longrightarrow \Top(X'){\quoi},
$$
使之在 $E$ 的点族 $X'$ 保守时为等价. 再推出典范拓扑斯态射
$$
E\longrightarrow \Top(X'){\quoi}.
$$
就从拓扑斯 $E$ 到形如 $\Top(Y)$ 的拓扑斯的态射而言, 给出此拓扑斯态射的泛刻画 (在等价意义下). (与此相关, 注意在等价意义下, $\Top(X')$ 不依赖于选定的小保守纤维函子族, 或等价地 (\hyperref[IV.subsec4.2]{4.2}), $X'_{\sob}$ 在同胚意义下与这种族的选取无关.)

\item 假设\pageoriginale 范畴 $\Ppoint(E)$ 不等价于小范畴 (参见 \hyperref[IV.subsec7.3]{7.3}), 且 $E$ 容许小保守纤维函子族. 证明若将这种族取得充分大, 则它定义的 $\Ppoint(E)$ 的子集 $X'$ 在诱导拓扑下不是素拓扑空间.

\item 与练习 \hyperref[IV.exer7.6]{7.6} 的关系参见 \hyperref[IV.sec8]{8}.
\end{enumerate}
\end{enonce*}

% END INCLUDED SGA4_English_current_buildable_source_20260906/source/04/drafts/SGA4_IV_LINES_4476_4830_English.texfrag


% BEGIN INCLUDED SGA4_English_current_buildable_source_20260906/source/04/drafts/SGA4_IV_LINES_4831_5055_English.texfrag
\section{局部化. 拓扑斯的开对象}\etiquette[8]{IV.sec8}

\subsection{}\etiquette[8.1]{IV.subsec8.1}
设 $C$ 为范畴, $T(X)$ 为涉及 $C$ 的一个对象 $X$ 的关系. 如果对 $C$ 的每个态射 $X'\rightarrow X$, 关系 $T(X)$ 蕴涵 $T(X')$, 就说关系 $T(X)$ (关于 $X$) \emph{对基变换稳定}. 这相当于说, $C$ 中由使 $T(X)$ 成立的对象 $X$ 组成的满子范畴 $R$ 为\emph{筛} (\hyperrefext[I][I.def4.1]{4.1}).

\begin{enonce*}[remark]{注 8.1.1}\etiquette[8.1.1]{IV.rem8.1.1}
假设 $C$ 为 $𝒰$-范畴 (\hyperrefext[I][I.def1.1]{1.1}), 从而可将 $C$ 等同于 $\widehat{C}$ 的满子范畴 (\hyperrefext[I][I.prop1.4]{1.4}). 有时将 $\ob C$ 上关系 $T$ 的定义扩充到 $\widehat{C}$ 的对象 $F$ 上的关系 $\widehat{T}$, 会较方便: 定义 $\widehat{T}(F)$ 表示对 $\widehat{C}$ 中每个满足 $X\in\ob C$ 的箭头 $X\rightarrow F$, 都有 $T(X)$ 成立. 关系 $\widehat{T}(F)$ 显然仍对 $F$ 的基变换稳定. 仍用 $R$ 表示由 $C$ 的筛 $R$ 定义的 $\widehat{C}$ 的终对象 $e$ 的子对象 (\hyperrefext[I][I.rem4.2.1]{4.2.1}), 则关系 $\widehat{T}(F)$ 也可表述为 $\Hom_C^{\widehat{\phantom{,}}}(F,R)\neq∅$, 即它表示唯一的态射 $F\rightarrow e$ 经过 $e$ 的子对象 $R$ 分解.
\end{enonce*}

\subsection{}\etiquette[8.2]{IV.subsec8.2}
现在假设 $C$ 为\emph{景}. 以 $X\in\ob C$ 为自变量的关系 $T(X)$ 称为\emph{对下降稳定}, 如果对 $C$ 的对象 $X$ 的每个覆盖\pageoriginale 族 $(X_i:X_i\rightarrow X)_{i\in I}$, 关系「对每个 $i\in I$ 有 $T(X_i)$」蕴涵 $T(X)$. 若 $T$ 对基变换稳定 (\hyperref[IV.subsec8.1]{8.1}) 且对下降稳定, 则称它为\emph{局部性质的关系}. 当 $C$ 为 $𝒰$-范畴, 并已假定 $T$ 对基变换稳定, 即由 $C$ 的筛 $R$ 定义, 且将该筛等同于 $C$ 上终预层的子预层时, 立即可见 $T$ 具有局部性质当且仅当 $R$ 为\emph{层}. 因而\emph{以 $X\in\ob C$ 为自变量的局部性质的关系, 恰好对应于 $C$ 的终层的子层, 即拓扑斯 $C^{\sim}$ 的终对象的子对象}. 所以它们实质上只依赖于由景 $C$ 定义的拓扑斯 $C^{\sim}$ (与景相联系的一切重要概念都是如此!). 用终层的子层 $R$, 关系 $T(X)$ 表为 $\Hom(\epsilon(X),R)\neq∅$. 它典范地扩充为拓扑斯 $C^{\sim}$ 上的关系 $T^{\sim}(F):\Hom(F,R)\neq∅$.

\begin{enonce*}[remark]{注 8.2.1}\etiquette[8.2.1]{IV.rem8.2.1}
用 \hyperref[IV.rem8.1.1]{8.1.1} 中引入的记号, 对 $C$ 上的预层 $F$, 因为 $\Hom(F,R)\simeq\Hom(\underline{a}(F),R)$, 关系 $\widehat{T}(F)$ 等价于关系 $\widehat{T}(\underline{a}(F))$.
\end{enonce*}

\begin{enonce*}{定义 8.3}\etiquette[8.3]{IV.def8.3}
$𝒰$-拓扑斯 $E$ 的一个开对象, 是 $E$ 的终对象 $e_E$ 的子对象; 若 $X$ 为 $E$ 的对象, 有时把诱导拓扑斯 $E_{/X}$ 的开对象, 即 $X$ 的子对象, 称为 $X$ 的开对象.

$𝒰$-景 $C$ 的开对象, 是相应 $𝒰$-拓扑斯 $C^{\sim}$ 的开对象.
\end{enonce*}

因为 $E$ 的终对象典范同构, \hyperref[IV.def8.3]{8.3} 中选取 $e_E$ 所引入的歧义是无害的. 还\pageoriginale 可以用更内在但实践中较不方便的方式, 把 $E$ 的开对象定义为 $E$ 的这样的满子范畴 $R$: 对 $E$ 的对象 $F$, 关系 $F\in\ob R$ 具有局部性质. 同样, $𝒰$-景 $C$ 的开对象与 $C$ 的那些对 $C$ 的对象使关系 $F\in\ob R$ 具有局部性质的满子范畴一一对应; 因而此概念实质上与所选宇宙 $𝒰$ 无关. 最后, 由 \hyperref[IV.subsec8.2]{8.2} 所述, 拓扑斯 (或景) 上的局部性质的关系, 实质上与该拓扑斯 (或景) 的开对象是一回事.

\setcounter{subsection}{3}
\subsection{开对象的例子}\etiquette[8.4]{IV.subsec8.4}

\subsubsection{}\etiquette[8.4.1]{IV.subsubsec8.4.1}
设 $X$ 为拓扑空间, 并将 $\Top(X)$ (\hyperref[IV.subsec2.1]{2.1}) 解释为 $X$ 上的平展空间范畴. 因为 $\Top(X)$ 中的单态射显然就是单射, 所以拓扑斯 $\Top(X)$ 的开对象等同于空间 $X$ 的开子集. 更一般地, 拓扑斯 $\Top(X,G)$ 的开对象 (\hyperref[IV.subsec2.4]{2.4}) 等同于 $X$ 中在 $G$ 作用下\emph{不变}的开子集.

\subsubsection{}\etiquette[8.4.2]{IV.subsubsec8.4.2}
$𝒰$-拓扑斯 $E$ 的开对象构成 $𝒰$-小集 (\hyperrefext[I][I.prop7.4]{7.4}), 按包含关系排序, 容许任意上确界和有限下确界. 它的最小元为始对象 (或「空对象」) $\phi_E$, 最大元为 $E$ 的终对象 $e_E$. 这两个元素只有在 $E$ 为「空拓扑斯」(\hyperref[IV.subsec2.2]{2.2}) 时才相同.

\subsubsection{}\etiquette[8.4.3]{IV.subsubsec8.4.3}
设 $G$ 为 $E$ 中的幺半群对象, 从而给出拓扑斯 $B_G$ (\sisi{\hyperref[IV.subsec2.4]{2.4}}{\hyperref[IV.subsec2.6]{2.6}}), 即 $E$ 中具有 $G$ 左作用的对象范畴. $B_G$ 的终对象 $e_G$ 为\pageoriginale $E$ 的终对象 $e_E$ 配上 $G$ 的平凡作用. 因而 $B_G$ 的开对象有序集典范同构于 $E$ 的开对象有序集. 特别地, 若 $E$ 为点拓扑斯, 从而 $G$ 为通常的幺半群, 则 $B_G$ 的开对象集恰有两个元素, 即 $\phi_{B_G}$ 与 $e_G$.

\subsubsection{}\etiquette[8.4.4]{IV.subsubsec8.4.4}
若 $E$ 为形如 $\widehat{C}$ 的拓扑斯, 其中 $C$ 为 $𝒰$-范畴, 则 (\hyperref[IV.subsec8.1]{8.1}) $E$ 的开对象有序集同构于 $C$ 的筛的有序集.

\subsubsection{}\etiquette[8.4.5]{IV.subsubsec8.4.5}
设 $X$ 为非离散的素拓扑空间. 证明 $\TOP(X)$ (\hyperref[IV.subsec2.5]{2.5}) 有一些开对象不是通过典范态射 \eqhyperref[IV.eq4.10.1]{4.10.1} $\TOP(X)\rightarrow\Top(X)$ 从 $X$ 的开子集, 即从 $\Top(X)$ 的开对象, 作逆像得到的.

\subsection{局部性质的关系的例子}\etiquette[8.5]{IV.subsec8.5}

对拓扑斯 $E$ 的每个对象 $X$, 用 $F\rightarrow F_X$ 表示\emph{局部化函子} $j_X^{\ast}:E\rightarrow E_{/X}$, $Y\mapsto Y\times X$ (\hyperref[IV.subsec5.2]{5.2}).

\subsubsection{}\etiquette[8.5.1]{IV.subsubsec8.5.1}
设 $u:F\rightarrow G$ 为 $E$ 的态射. 以下关于自变量 $X\in\ob E$ 的关系:

$u_X:F_X\rightarrow G_X$ 为单态射 (分别为满态射, 分别为同构)

\noindent 为局部性质的关系. 特别地, 若 $G$ 为 $E$ 的群对象, 则关系「$G_X$ 为 $X$ 上的平凡群」具有局部性质; 相应的 $E$ 的开对象称为\emph{群 $G$ 的余支集}.

\subsubsection{}\etiquette[8.5.2]{IV.subsubsec8.5.2}
设\pageoriginale $u,v:F\rightrightarrows G$ 为 $E$ 的两个态射. 关于自变量 $X\in\ob E$ 的关系
$$
u_X=v_X:F_X\longrightarrow G_X
$$
为局部性质的关系. 特别地, 若 $G$ 为 $E$ 的群对象, $u$ 为 $G$ 的一个\emph{截面} (\hyperref[IV.subsubsec4.3.6]{4.3.6}), 则关于 $X$ 的关系
$$
\mbox{$E_{/X}$ 的群 $G_X$ 的截面 $u_X$ 为零}
$$
具有局部性质; 相应的 $E$ 的开对象称为\emph{群 $G$ 的截面 $u$ 的余支集}.

\begin{enonce*}[remark]{注 8.5.3}\etiquette[8.5.3]{IV.rem8.5.3}
当 $E=\Top(X)$ 时, 群层 $G$, 分别地这样的层 $G$ 的一个截面的余支集, 按经典术语 [TF], 就是 $X$ 中分别与 \emph{$G$ 的支集}、\emph{$u$ 的支集}互补的开子集.
\end{enonce*}

\setcounter{subsubsection}{3}
\subsubsection{}\etiquette[8.5.4]{IV.subsubsec8.5.4}
设 $P(f)$ 为以 $f\in{\rm f\ell\,}C$ 为自变量的关系, 其中 $C$ 为景. 当纤维积在 $C$ 中可表时, 关系 $P(f)$ 称为\emph{在靶上} (或\emph{在基上}, 或「\emph{在下方}」) \emph{对基变换稳定} (分别为\emph{对下降稳定}、\emph{具有局部性质}), 如果对 $C$ 的每个态射 $f:X\rightarrow Y$, 以下关于自变量 $Y'\in\ob C_{/Y}$ 的关系:

「从 $f$ 沿结构态射 $Y'\rightarrow Y$ 作基变换得到的态射 $f':X'=\fprod{X}{Y'}{Y}\rightarrow Y'$ 满足 $P(f')$」

对基变换稳定, 分别对下降稳定、具有局部性质. 当 $C$ 的拓扑通过预拓扑定义时, 如果对 $C$ 的每个态射 $f:X\rightarrow Y$ 及 $\Cov(X)$ 中的每个族 $(g_i:X_i\rightarrow X)_{i\in I}$,\pageoriginale 关系 $P(f)$ 等价于关系「对每个 $i\in I$ 有 $P(fg_i)$」, 就说关系 $P(f)$ \emph{在源上} (或\emph{「在上方」}) \emph{具有局部性质}. 当 $\Cov(X)$ 由 $C$ 的全部覆盖族组成时, 此条件也表示诱导景 $C_{/X}$ 的对象 $X'$ 上的下述关系
\begin{center}
$P(fg)$, 其中 $g:X'\rightarrow X$ 为结构态射
\end{center}
具有局部性质. 注意若关系 $P(f)$ 在上方具有局部性质, 则它更在下方具有局部性质.

\subsubsection{}\etiquette[8.5.5]{IV.subsubsec8.5.5}
例如, 取 $C=(\Sch)$ 为 $𝒰$ 中的概形范畴, 赋予忠实平坦拟紧拓扑 (\SGA 3 IV 6.3) 或更粗的拓扑. 此时下列每种态射性质都在下方具有局部性质:

满射、根性、泛开、泛闭、固有、拟紧、拟紧且稠密、泛同胚、分离、拟分离、局部有限型、局部有限表示、有限型、有限表示、同构、单态射、开浸入、闭浸入、拟紧浸入、仿射、拟仿射、有限、拟有限、整、平坦、忠实平坦、非分歧、光滑、平展

\noindent (\EGA IV 2.6.4, 2.7.1 及 \EGA IV 17.7.1). 另一方面, 给 $(\Sch)$ 赋予如下预拓扑: 对每个概形 $X\in𝒰$, $\Cov(X)$ 由满射的态射族 $(f_i:X_i\rightarrow X)_{i\in I}$ 组成, 其中各 $f_i$ 平坦且局部有限表示. 此时下列每种性质都\emph{在上方具有局部性质}:

局部有限型、局部有限表示、有限型、平坦、非分歧、光滑、平展.

\noindent (\EGA IV 17.7.5, 17.7.7).

% END INCLUDED SGA4_English_current_buildable_source_20260906/source/04/drafts/SGA4_IV_LINES_4831_5055_English.texfrag


% BEGIN INCLUDED SGA4_English_current_buildable_source_20260906/source/04/drafts/SGA4_IV_LINES_5056_5253_English.texfrag
\begin{enonce*}[remark]{练习 8.6}\etiquette[8.6]{IV.exer8.6}
设\pageoriginale $f:E\rightarrow F$ 为 $𝒰$-拓扑斯态射. 考虑关于自变量 $X\in\ob F$ 的下列关系: 诱导态射 $E_{/f^{\ast}(X)}\rightarrow F_{/X}$ 为拓扑斯等价. 证明此关系具有局部性质.
\end{enonce*}

\begin{enonce*}[remark]{练习 8.7}\etiquette[8.7]{IV.exer8.7}
(\emph{拓扑斯的分拆, 拓扑斯的和}). 设 $E$ 为 $𝒰$-拓扑斯. $E$ 的开对象族 $(e_i)_{i\in I}$, 即 $E$ 的终对象 $e$ 的子对象族, 称为 \emph{$e$ (或拓扑斯 $E$) 的分拆}, 如果典范态射 $\amalg_{i\in I}e_i\rightarrow e$ 为同构, 即 (\hyperrefext[II][II.cor4.6.2]{4.6.2}) $e$ 为各 $e_i$ 的上确界, 且 $i\neq j$ 蕴涵 $e_i\cap e_j=∅_E$.
\begin{enumerate}
\renewcommand{\theenumi}{\alph{enumi}}
\renewcommand{\labelenumi}{\theenumi)}
\item 证明 $E$ 的对象族 $(e_i)_{i\in I}$ 为 $E$ 的分拆, 当且仅当由局部化函子 $E\rightarrow E_{/e_i}$ 定义的函子
\begin{equation*}
E\longrightarrow \prod_i E_{/e_i}\tag{8.7.1}
\end{equation*}\etiquette[8.7.1]{IV.eq8.7.1}
为范畴等价.

\item 反之, 设 $(E_i)_{i\in I}$ 为 $𝒰$-拓扑斯族, 其中 $I\in𝒰$. 证明 $E=\Pi_{i\in I}E_i$ 为 $𝒰$-拓扑斯 (称为族 $(E_i)_{i\in I}$ 的\emph{和拓扑斯}). 定义 $E$ 的分拆 $(e_i)_{i\in I}$ 及范畴等价 $E_{/e_i}\approx E_i$, 使投影函子 $E\rightarrow E_i$ 等同于局部化函子 $E\rightarrow E_{/e_i}$.

\item 设 $E$ 为族 $(E_i)_{i\in I}$ 的和拓扑斯. 证明对每个 $i\in I$, 投影函子 $E\rightarrow E_i$ 形如 $u_i^{\ast}$, 其中
\begin{equation*}
u_i:E_i\longrightarrow E\tag{8.7.2}
\end{equation*}\etiquette[8.7.2]{IV.eq8.7.2}
为\pageoriginale 拓扑斯态射. 证明对每个 $𝒰$-拓扑斯 $F$, 函子 $v\mapsto(v\circ u_i)_{i\in I}$
\begin{equation*}
\Hhomtop(E,F)\longrightarrow \prod_i\Hhomtop(E_i,F)\tag{8.7.3}
\end{equation*}\etiquette[8.7.3]{IV.eq8.7.3}
为范畴等价.

\item 设 $(e_i)_{i\in I}$ 为拓扑斯 $E$ 的分拆. 对每个拓扑斯态射 $g:F\rightarrow E$, 其中 $f_i=g^{\ast}(e_i)$ 的族 $(f_i)_{i\in I}$ 为 $F$ 的分拆; 诱导态射 $g_i:F_{/f_i}\rightarrow E_{/e_i}$ 使得可以 (在唯一同构意义下) 重构 $g$, 函子 $g^{\ast}$ 等同于各函子 $g_i^{\ast}$ 的笛卡儿积 (这里考虑到 \hyperref[IV.exer8.7]{8.7} 1) 类型的等价). 推出可用形如 $\Hhomtop(F',E_i)$ 的范畴完整描述范畴 $\Hhomtop(F,E)$, 其中 $F'$ 为形如 $F_{/f}$ 的拓扑斯 ($f$ 为 $F$ 的终对象的直和项), $E_i=E_{/e_i}$.

\item 特别地, 若 $F$ 连通且非空 (参见 \hyperref[IV.subsubsec4.3.5]{4.3.5}), 即对 $F$ 的每个分拆 $(f_i)_{i\in I}$, 存在唯一的 $i\in I$ 使得 $f_i\neq\phi_F$, 证明由拓扑斯态射 \eqhyperref[IV.eq8.7.2]{8.7.2} 诱导的典范函子
\begin{equation*}
\coprod_{i\in I}\Hhomtop(F,E_i)\longrightarrow
\Hhomtop(F,E)\tag{8.7.4}
\end{equation*}\etiquette[8.7.4]{IV.eq8.7.4}
为范畴等价. 更特别地, 拓扑斯态射族 \eqhyperref[IV.eq8.7.2]{8.7.2} 诱导范畴等价
$$
\coprod_{i\in I}\Ppoint(E_i)\xrightarrow{\approx}\Ppoint(E){\quoi}.
$$

\item $E$ 的开对象集的子集\pageoriginale $I$ 称为 $E$ 的\emph{约化分拆}, 如果以 $I$ 为指标的恒等族 $(e_i)_{i\in I}$ 为分拆, 且各 $e_i$ 均异于 $\phi_E$. 证明约化分拆之间的加细关系 (\hyperrefext[I][I.sec4.3.2]{4.3.2}), 使约化分拆集 $P$ 成为 $𝒰$-小的有序集, 其中任意两个元素的上确界存在. 这给出严格投射系 $(I_p)_{p\in P}$, 看作 pro-集, 记为 $\pi_{\circ}(E)$. 称为\emph{拓扑斯 $E$ 的 pro-$\pi_{\circ}$}.

\item 证明 $\pi_{\circ}(E)$ pro-表示与从 $E$ 到点拓扑斯 $P$ 的典范态射 $f:E\rightarrow P$ (\hyperref[IV.subsec4.3]{4.3}) 相联系的函子
$$
T\longmapsto f_{\ast}f^{\ast}(T)=\Gamma(E,
T_E):(𝒰\text{-}\Ens)\longrightarrow (𝒰\text{-}\Ens)
$$
特别地, 此函子可表当且仅当 $\pi_{\circ}(E)$ 本质常值, 即 (作为 pro-集) 同构于一个仍记为 $\pi_{\circ}(E)$ 的「通常」集合. pro-集 $\pi_{\circ}(E)$ 只有一点当且仅当 $E$「连通且非空」. 它为空当且仅当 $E$ 为「空拓扑斯」(\hyperref[IV.subsec2.2]{2.2}).

\item 假设 $E$ 拟紧, 即其终对象 $e$ 的每个覆盖都容许有限子覆盖. 证明此时 $\pi_{\circ}(E)$ 为 pro-有限集, 因而可将它 (利用有限集范畴的 pro-对象与紧全不连通空间之间的熟悉等价) 等同于紧全不连通空间, 仍记为 $\pi_{\circ}(E)$.

\item 设 $X$ 为拓扑空间, $\pi_{\circ}(X)$ 为 $X$ 的连通分支集, 看作常值 pro-集. 定义典范\pageoriginale pro-集态射
\begin{equation*}
\pi_{\circ}(X)\longrightarrow
\pi_{\circ}(\Top(X)){\quoi},\tag{8.7.5}
\end{equation*}\etiquette[8.7.5]{IV.eq8.7.5}
或者等价地, 定义典范映射
\begin{equation*}
\pi_{\circ}(X)\longrightarrow \varprojlim
\pi_{\circ}(\Top(X)){\quoi}.\tag{8.7.6}
\end{equation*}\etiquette[8.7.6]{IV.eq8.7.6}
证明第一个态射为同构当且仅当 $X$ 同胚于连通空间的和空间 (特别地, 当 $X$ 局部连通时如此).

\item 假设 $X$ 为拟紧概形. 证明 \eqhyperref[IV.eq8.7.6]{8.7.6} 为双射.

\item 明确「函子性」的意义后, 建立态射 \eqhyperref[IV.eq8.7.5]{8.7.5} 与 \eqhyperref[IV.eq8.7.6]{8.7.6} 关于 $X$ 的「函子性」.


\item (比较 \hyperref[IV.exer7.6]{7.6} i).) 证明关于 $𝒰$-拓扑斯 $E$ 的以下两个条件等价: (i) 从 $E$ 到点拓扑斯的典范态射为「本质的」(练习 \hyperref[IV.exer7.6]{7.6} a)), 即从 $(𝒰\text{-}\Ens)$ 到 $E$ 的函子 $I\mapsto I_E$ 与以 $𝒰$ 中的集合为指标的积交换. (ii) 对 $E$ 的每个对象 $X$, 诱导拓扑斯 $E_{/X}$ 的 $\pi_{\circ}(E_{/X})$ 是通常的集合, 即 $X$ 同构于 $E$ 的一个连通对象族的和. 在这种情形, 称 $E$ 为\emph{局部连通拓扑斯} (比较 \sisi{IV~}{}\hyperref[IV.exer2.7.5]{2.7.5}).
\end{enumerate}
\end{enonce*}

\begin{enonce*}[remark]{练习 8.8}\etiquette[8.8]{IV.exer8.8}
(\emph{拓扑斯的稠密态射}).
\begin{enumerate}
\renewcommand{\theenumi}{\alph{enumi}}
\renewcommand{\labelenumi}{\theenumi)}
\item 设 $f:E'\rightarrow E$ 为拓扑斯态射. 证明以下条件等价:

\begin{itemize}
\item[(i)] $f_{\ast}(\phi_{E'})=\phi_E$ (其中 $\phi_E$ 与 $\phi_{E'}$ 分别表示 $E$ 与 $E'$ 的始对象).

\item[(ii)] 对\pageoriginale $E$ 的每个「非空」对象 $X$, 即不同构于 $\phi_E$ 的对象, $f^{\ast}(X)$ 为 $E'$ 中的非空对象.

\item[$\rm(ii')$] 同 (ii), 其中 $X$ 为 $E$ 的终对象 $e_E$ 的子对象, 即 $X$ 为拓扑斯 $E$ 的开对象.
\end{itemize}

此时称拓扑斯态射 $f$ 为\emph{稠密的}.

\item 设 $f:X'\rightarrow X$ 为拓扑空间的连续映射, 给出拓扑斯态射 $\Top(f):\Top(X')\rightarrow\Top(X)$ (\hyperref[IV.subsec4.1]{4.1}). 证明 $\Top(f)$ 稠密当且仅当 $f$ 稠密, 即 $f(X')$ 在 $X$ 中稠密.

\item 证明若 $f:E'\rightarrow E$ 为「保守」的拓扑斯态射 (\hyperref[IV.subsubsec6.4.0]{6.4.0}), 即 $f^{\ast}$ 忠实, 则 $f$ 稠密. 证明逆命题不一定成立, 即使对局部化态射 $E_{/U}\rightarrow E$ 也如此, 其中 $U$ 为拓扑斯 $E$ 的开对象. (证明在此情形, $f$ 只有在 $U$ 为 $E$ 的终对象时才保守 (从而只有在 $f$ 为拓扑斯等价时才保守), 并用 b) 中的例子.)

\item 设 $f:X'\rightarrow X$ 为拓扑斯 $E$ 中的箭头. 若相应的诱导拓扑斯态射 $E_{/X'}\rightarrow E_{/X}$ \eqhyperref[IV.eq5.5.2]{5.5.2} 稠密, 则称 $f$ 为拓扑斯 $E$ 中的\emph{稠密态射}. 证明此性质只依赖于 $X'$ 在 $X$ 中的像 $f(X')$ 作为 $X$ 的子对象有序集 $\ouv(X)$ 的元素. 证明态射 $E_{/X'}\rightarrow E_{/X}$ 保守当且仅当 $f$ 为覆盖.
\end{enumerate}
\end{enonce*}

% END INCLUDED SGA4_English_current_buildable_source_20260906/source/04/drafts/SGA4_IV_LINES_5056_5253_English.texfrag


% BEGIN INCLUDED SGA4_English_current_buildable_source_20260906/source/04/drafts/SGA4_IV_LINES_5254_5490_English.texfrag
